%% Ngân hàng đề Toán 12 (KNTT) -- Chương I, Bài 5, Dạng 01
%% Dạng: Ứng dụng đạo hàm trong vật lí chuyển động (vận tốc, gia tốc, độ cao, quãng đường)
%% Phương án ĐẦU TIÊN trong \cauTN luôn là đáp án đúng (gói trondethi tự đảo).
%% Phần II: \y{nội dung}{d|s}. Phần III: đáp án đặt ở đối số cuối.

\cauTN[id=C1-B5-D01-NB-001,muc=NB,loigiai={Vận tốc $v(t)=s'(t)=-2t^{2}+36t+5$, gia tốc $a(t)=v'(t)=-4t+36$. $v'(t)=0\Leftrightarrow t=9$; $v$ đồng biến trên $(0;9)$, nghịch biến trên $(9;14)$ nên $v_{\max}=v(9)=167$.}]{Một vật chuyển động theo quy luật $s(t)=-\dfrac{2}{3}t^3+18t^2+5t$, trong đó $t$ (giây) là khoảng thời gian tính từ lúc vật bắt đầu chuyển động và $s$ (mét) là quãng đường vật đi được trong khoảng thời gian đó, $0\le t\le 14$. Vận tốc lớn nhất của vật trong khoảng thời gian trên bằng}
  {$167\,\mathrm{m/s}$.}
  {$9\,\mathrm{m/s}$.}{$117\,\mathrm{m/s}$.}{$147\,\mathrm{m/s}$.}

\cauTN[id=C1-B5-D01-NB-002,muc=NB,loigiai={Vận tốc $v(t)=s'(t)=24-5t$. Ô tô dừng khi $v(t)=0\Leftrightarrow t=\dfrac{24}{5}$.}]{Một ô tô đang chạy với vận tốc $24$ m/s thì người lái đạp phanh. Từ thời điểm đó, ô tô chuyển động chậm dần đều và quãng đường đi được sau $t$ giây là $s(t)=24t-\dfrac{5}{2}t^2$ (mét). Thời gian từ lúc đạp phanh đến khi ô tô dừng hẳn bằng}
  {$\dfrac{24}{5}\,\mathrm{s}$.}
  {$\dfrac{5}{24}\,\mathrm{s}$.}{$\dfrac{29}{5}\,\mathrm{s}$.}{$\dfrac{19}{5}\,\mathrm{s}$.}

\cauTN[id=C1-B5-D01-NB-003,muc=NB,loigiai={Ta có $h(\dfrac{1}{2})=\dfrac{99}{4}$.}]{Một quả bóng được ném thẳng đứng lên trên từ độ cao $20$ m với vận tốc ban đầu $12$ m/s. Độ cao của quả bóng (tính bằng mét) sau $t$ giây là $h(t)=20+12t-5t^2$ (bỏ qua sức cản không khí, lấy $g=10$ m/s$^2$). Độ cao của quả bóng tại thời điểm $t=\dfrac{1}{2}$ giây bằng}
  {$\dfrac{99}{4}\,\mathrm{m}$.}
  {$99\,\mathrm{m}$.}{$\dfrac{99}{2}\,\mathrm{m}$.}{$\dfrac{4}{99}\,\mathrm{m}$.}

\cauTN[id=C1-B5-D01-NB-004,muc=NB,loigiai={Ta có $h(\dfrac{1}{2})=\dfrac{65}{4}$.}]{Một quả bóng được ném thẳng đứng lên trên từ độ cao $5$ m với vận tốc ban đầu $25$ m/s. Độ cao của quả bóng (tính bằng mét) sau $t$ giây là $h(t)=5+25t-5t^2$ (bỏ qua sức cản không khí, lấy $g=10$ m/s$^2$). Độ cao của quả bóng tại thời điểm $t=\dfrac{1}{2}$ giây bằng}
  {$\dfrac{65}{4}\,\mathrm{m}$.}
  {$\dfrac{4}{65}\,\mathrm{m}$.}{$65\,\mathrm{m}$.}{$\dfrac{61}{4}\,\mathrm{m}$.}

\cauTN[id=C1-B5-D01-NB-005,muc=NB,loigiai={Vận tốc $v(t)=s'(t)=30-5t$. Ô tô dừng khi $v(t)=0\Leftrightarrow t=6$.}]{Một ô tô đang chạy với vận tốc $30$ m/s thì người lái đạp phanh. Từ thời điểm đó, ô tô chuyển động chậm dần đều và quãng đường đi được sau $t$ giây là $s(t)=30t-\dfrac{5}{2}t^2$ (mét). Thời gian từ lúc đạp phanh đến khi ô tô dừng hẳn bằng}
  {$6\,\mathrm{s}$.}
  {$12\,\mathrm{s}$.}{$10\,\mathrm{s}$.}{$2\,\mathrm{s}$.}

\cauTN[id=C1-B5-D01-NB-006,muc=NB,loigiai={Vận tốc $v(t)=s'(t)=-t^{2}+16t$, gia tốc $a(t)=v'(t)=-2t+16$. Do đó $v(2)=28$.}]{Một vật chuyển động theo quy luật $s(t)=-\dfrac{1}{3}t^3+8t^2$, trong đó $t$ (giây) là khoảng thời gian tính từ lúc vật bắt đầu chuyển động và $s$ (mét) là quãng đường vật đi được trong khoảng thời gian đó, $0\le t\le 11$. Vận tốc của vật tại thời điểm $t=2$ giây bằng}
  {$28\,\mathrm{m/s}$.}
  {$29\,\mathrm{m/s}$.}{$14\,\mathrm{m/s}$.}{$30\,\mathrm{m/s}$.}

\cauTN[id=C1-B5-D01-NB-007,muc=NB,loigiai={Ta có $v(t)=h'(t)=20-10t$. $h'(t)=0\Leftrightarrow t=2$.}]{Một quả bóng được ném thẳng đứng lên trên từ độ cao $2$ m với vận tốc ban đầu $20$ m/s. Độ cao của quả bóng (tính bằng mét) sau $t$ giây là $h(t)=2+20t-5t^2$ (bỏ qua sức cản không khí, lấy $g=10$ m/s$^2$). Thời điểm quả bóng đạt độ cao lớn nhất là}
  {$2\,\mathrm{s}$.}
  {$1\,\mathrm{s}$.}{$3\,\mathrm{s}$.}{$4\,\mathrm{s}$.}

\cauTN[id=C1-B5-D01-NB-008,muc=NB,loigiai={Ta có $v(t)=h'(t)=30-10t$. Do đó $v(2)=10$.}]{Một quả bóng được ném thẳng đứng lên trên từ độ cao $0$ m với vận tốc ban đầu $30$ m/s. Độ cao của quả bóng (tính bằng mét) sau $t$ giây là $h(t)=30t-5t^2$ (bỏ qua sức cản không khí, lấy $g=10$ m/s$^2$). Vận tốc của quả bóng tại thời điểm $t=2$ giây bằng}
  {$10\,\mathrm{m/s}$.}
  {$14\,\mathrm{m/s}$.}{$5\,\mathrm{m/s}$.}{$8\,\mathrm{m/s}$.}

\cauTN[id=C1-B5-D01-NB-009,muc=NB,loigiai={Ta có $h(\dfrac{1}{2})=\dfrac{55}{4}$.}]{Một quả bóng được ném thẳng đứng lên trên từ độ cao $0$ m với vận tốc ban đầu $30$ m/s. Độ cao của quả bóng (tính bằng mét) sau $t$ giây là $h(t)=30t-5t^2$ (bỏ qua sức cản không khí, lấy $g=10$ m/s$^2$). Độ cao của quả bóng tại thời điểm $t=\dfrac{1}{2}$ giây bằng}
  {$\dfrac{55}{4}\,\mathrm{m}$.}
  {$\dfrac{59}{4}\,\mathrm{m}$.}{$\dfrac{51}{4}\,\mathrm{m}$.}{$\dfrac{55}{2}\,\mathrm{m}$.}

\cauTN[id=C1-B5-D01-NB-010,muc=NB,loigiai={Vận tốc $v(t)=s'(t)=-2t^{2}+8t$, gia tốc $a(t)=v'(t)=-4t+8$. $v'(t)=0\Leftrightarrow t=2$ và $v'$ đổi dấu từ dương sang âm nên vận tốc lớn nhất tại $t=2$.}]{Một vật chuyển động theo quy luật $s(t)=-\dfrac{2}{3}t^3+4t^2$, trong đó $t$ (giây) là khoảng thời gian tính từ lúc vật bắt đầu chuyển động và $s$ (mét) là quãng đường vật đi được trong khoảng thời gian đó, $0\le t\le 5$. Thời điểm vật đạt vận tốc lớn nhất là}
  {$2\,\mathrm{s}$.}
  {$8\,\mathrm{s}$.}{$4\,\mathrm{s}$.}{$1\,\mathrm{s}$.}

\cauTN[id=C1-B5-D01-NB-011,muc=NB,loigiai={Ta có $v(t)=h'(t)=15-10t$. $h'(t)=0\Leftrightarrow t=\dfrac{3}{2}$.}]{Một quả bóng được ném thẳng đứng lên trên từ độ cao $1$ m với vận tốc ban đầu $15$ m/s. Độ cao của quả bóng (tính bằng mét) sau $t$ giây là $h(t)=1+15t-5t^2$ (bỏ qua sức cản không khí, lấy $g=10$ m/s$^2$). Thời điểm quả bóng đạt độ cao lớn nhất là}
  {$\dfrac{3}{2}\,\mathrm{s}$.}
  {$\dfrac{3}{4}\,\mathrm{s}$.}{$3\,\mathrm{s}$.}{$\dfrac{2}{3}\,\mathrm{s}$.}

\cauTN[id=C1-B5-D01-NB-012,muc=NB,loigiai={Vận tốc $v(t)=s'(t)=25-4t$. Gia tốc $a(t)=v'(t)=-4$.}]{Một ô tô đang chạy với vận tốc $25$ m/s thì người lái đạp phanh. Từ thời điểm đó, ô tô chuyển động chậm dần đều và quãng đường đi được sau $t$ giây là $s(t)=25t-2t^2$ (mét). Gia tốc của ô tô bằng}
  {$-4\,\mathrm{m/s^2}$.}
  {$-2\,\mathrm{m/s^2}$.}{$0\,\mathrm{m/s^2}$.}{$4\,\mathrm{m/s^2}$.}

\cauTN[id=C1-B5-D01-NB-013,muc=NB,loigiai={Vận tốc $v(t)=s'(t)=24-5t$. Gia tốc $a(t)=v'(t)=-5$.}]{Một ô tô đang chạy với vận tốc $24$ m/s thì người lái đạp phanh. Từ thời điểm đó, ô tô chuyển động chậm dần đều và quãng đường đi được sau $t$ giây là $s(t)=24t-\dfrac{5}{2}t^2$ (mét). Gia tốc của ô tô bằng}
  {$-5\,\mathrm{m/s^2}$.}
  {$-3\,\mathrm{m/s^2}$.}{$5\,\mathrm{m/s^2}$.}{$-10\,\mathrm{m/s^2}$.}

\cauTN[id=C1-B5-D01-NB-014,muc=NB,loigiai={Vận tốc $v(t)=s'(t)=24-6t$. Ô tô dừng khi $v(t)=0\Leftrightarrow t=4$.}]{Một ô tô đang chạy với vận tốc $24$ m/s thì người lái đạp phanh. Từ thời điểm đó, ô tô chuyển động chậm dần đều và quãng đường đi được sau $t$ giây là $s(t)=24t-3t^2$ (mét). Thời gian từ lúc đạp phanh đến khi ô tô dừng hẳn bằng}
  {$4\,\mathrm{s}$.}
  {$5\,\mathrm{s}$.}{$6\,\mathrm{s}$.}{$3\,\mathrm{s}$.}

\cauTN[id=C1-B5-D01-NB-015,muc=NB,loigiai={Vận tốc $v(t)=s'(t)=30-3t$. $v(2)=24$.}]{Một ô tô đang chạy với vận tốc $30$ m/s thì người lái đạp phanh. Từ thời điểm đó, ô tô chuyển động chậm dần đều và quãng đường đi được sau $t$ giây là $s(t)=30t-\dfrac{3}{2}t^2$ (mét). Vận tốc của ô tô sau $2$ giây kể từ lúc đạp phanh bằng}
  {$24\,\mathrm{m/s}$.}
  {$19\,\mathrm{m/s}$.}{$48\,\mathrm{m/s}$.}{$12\,\mathrm{m/s}$.}

\cauTN[id=C1-B5-D01-NB-016,muc=NB,loigiai={Vận tốc $v(t)=s'(t)=-3t^{2}+60t+2$, gia tốc $a(t)=v'(t)=-6t+60$. $v'(t)=0\Leftrightarrow t=10$; $v$ đồng biến trên $(0;10)$, nghịch biến trên $(10;20)$ nên $v_{\max}=v(10)=302$.}]{Một vật chuyển động theo quy luật $s(t)=-\dfrac{3}{3}t^3+30t^2+2t$, trong đó $t$ (giây) là khoảng thời gian tính từ lúc vật bắt đầu chuyển động và $s$ (mét) là quãng đường vật đi được trong khoảng thời gian đó, $0\le t\le 20$. Vận tốc lớn nhất của vật trong khoảng thời gian trên bằng}
  {$302\,\mathrm{m/s}$.}
  {$10\,\mathrm{m/s}$.}{$2\,\mathrm{m/s}$.}{$252\,\mathrm{m/s}$.}

\cauTN[id=C1-B5-D01-NB-017,muc=NB,loigiai={Vận tốc $v(t)=s'(t)=-t^{2}+16t+2$, gia tốc $a(t)=v'(t)=-2t+16$. Do đó $a(7)=2$.}]{Một vật chuyển động theo quy luật $s(t)=-\dfrac{1}{3}t^3+8t^2+2t$, trong đó $t$ (giây) là khoảng thời gian tính từ lúc vật bắt đầu chuyển động và $s$ (mét) là quãng đường vật đi được trong khoảng thời gian đó, $0\le t\le 16$. Gia tốc của vật tại thời điểm $t=7$ giây bằng}
  {$2\,\mathrm{m/s^2}$.}
  {$3\,\mathrm{m/s^2}$.}{$4\,\mathrm{m/s^2}$.}{$1\,\mathrm{m/s^2}$.}

\cauTN[id=C1-B5-D01-NB-018,muc=NB,loigiai={Vận tốc $v(t)=s'(t)=-t^{2}+18t+3$, gia tốc $a(t)=v'(t)=-2t+18$. Do đó $a(17)=-16$.}]{Một vật chuyển động theo quy luật $s(t)=-\dfrac{1}{3}t^3+9t^2+3t$, trong đó $t$ (giây) là khoảng thời gian tính từ lúc vật bắt đầu chuyển động và $s$ (mét) là quãng đường vật đi được trong khoảng thời gian đó, $0\le t\le 18$. Gia tốc của vật tại thời điểm $t=17$ giây bằng}
  {$-16\,\mathrm{m/s^2}$.}
  {$-11\,\mathrm{m/s^2}$.}{$-21\,\mathrm{m/s^2}$.}{$-32\,\mathrm{m/s^2}$.}

\cauTN[id=C1-B5-D01-NB-019,muc=NB,loigiai={Vận tốc $v(t)=s'(t)=-3t^{2}+60t+3$, gia tốc $a(t)=v'(t)=-6t+60$. Do đó $v(13)=276$.}]{Một vật chuyển động theo quy luật $s(t)=-\dfrac{3}{3}t^3+30t^2+3t$, trong đó $t$ (giây) là khoảng thời gian tính từ lúc vật bắt đầu chuyển động và $s$ (mét) là quãng đường vật đi được trong khoảng thời gian đó, $0\le t\le 20$. Vận tốc của vật tại thời điểm $t=13$ giây bằng}
  {$276\,\mathrm{m/s}$.}
  {$326\,\mathrm{m/s}$.}{$376\,\mathrm{m/s}$.}{$176\,\mathrm{m/s}$.}

\cauTN[id=C1-B5-D01-NB-020,muc=NB,loigiai={Ta có $h(\dfrac{1}{2})=\dfrac{75}{4}$.}]{Một quả bóng được ném thẳng đứng lên trên từ độ cao $0$ m với vận tốc ban đầu $40$ m/s. Độ cao của quả bóng (tính bằng mét) sau $t$ giây là $h(t)=40t-5t^2$ (bỏ qua sức cản không khí, lấy $g=10$ m/s$^2$). Độ cao của quả bóng tại thời điểm $t=\dfrac{1}{2}$ giây bằng}
  {$\dfrac{75}{4}\,\mathrm{m}$.}
  {$75\,\mathrm{m}$.}{$\dfrac{71}{4}\,\mathrm{m}$.}{$\dfrac{79}{4}\,\mathrm{m}$.}

\cauTN[id=C1-B5-D01-NB-021,muc=NB,loigiai={Ta có $v(t)=h'(t)=18-10t$. $h'(t)=0\Leftrightarrow t=\dfrac{9}{5}$, $h$ đạt giá trị lớn nhất bằng $h(\dfrac{9}{5})=\dfrac{86}{5}$.}]{Một quả bóng được ném thẳng đứng lên trên từ độ cao $1$ m với vận tốc ban đầu $18$ m/s. Độ cao của quả bóng (tính bằng mét) sau $t$ giây là $h(t)=1+18t-5t^2$ (bỏ qua sức cản không khí, lấy $g=10$ m/s$^2$). Độ cao lớn nhất mà quả bóng đạt được bằng}
  {$\dfrac{86}{5}\,\mathrm{m}$.}
  {$\dfrac{91}{5}\,\mathrm{m}$.}{$\dfrac{172}{5}\,\mathrm{m}$.}{$86\,\mathrm{m}$.}

\cauTN[id=C1-B5-D01-NB-022,muc=NB,loigiai={Quãng đường $s(10)=1420$.}]{Một vật chuyển động theo quy luật $s(t)=-\dfrac{3}{3}t^3+24t^2+2t$, trong đó $t$ (giây) là khoảng thời gian tính từ lúc vật bắt đầu chuyển động và $s$ (mét) là quãng đường vật đi được trong khoảng thời gian đó, $0\le t\le 12$. Quãng đường vật đi được trong $10$ giây đầu tiên bằng}
  {$1420\,\mathrm{m}$.}
  {$710\,\mathrm{m}$.}{$2840\,\mathrm{m}$.}{$1920\,\mathrm{m}$.}

\cauTN[id=C1-B5-D01-NB-023,muc=NB,loigiai={Vận tốc $v(t)=s'(t)=16-2t$. Ô tô dừng khi $v(t)=0\Leftrightarrow t=8$.}]{Một ô tô đang chạy với vận tốc $16$ m/s thì người lái đạp phanh. Từ thời điểm đó, ô tô chuyển động chậm dần đều và quãng đường đi được sau $t$ giây là $s(t)=16t-t^2$ (mét). Thời gian từ lúc đạp phanh đến khi ô tô dừng hẳn bằng}
  {$8\,\mathrm{s}$.}
  {$12\,\mathrm{s}$.}{$16\,\mathrm{s}$.}{$6\,\mathrm{s}$.}

\cauTN[id=C1-B5-D01-NB-024,muc=NB,loigiai={Vận tốc $v(t)=s'(t)=-2t^{2}+16t+5$, gia tốc $a(t)=v'(t)=-4t+16$. $v'(t)=0\Leftrightarrow t=4$; $v$ đồng biến trên $(0;4)$, nghịch biến trên $(4;8)$ nên $v_{\max}=v(4)=37$.}]{Một vật chuyển động theo quy luật $s(t)=-\dfrac{2}{3}t^3+8t^2+5t$, trong đó $t$ (giây) là khoảng thời gian tính từ lúc vật bắt đầu chuyển động và $s$ (mét) là quãng đường vật đi được trong khoảng thời gian đó, $0\le t\le 8$. Vận tốc lớn nhất của vật trong khoảng thời gian trên bằng}
  {$37\,\mathrm{m/s}$.}
  {$4\,\mathrm{m/s}$.}{$5\,\mathrm{m/s}$.}{$42\,\mathrm{m/s}$.}

\cauTN[id=C1-B5-D01-NB-025,muc=NB,loigiai={Quãng đường $s(2)=\dfrac{16}{3}$.}]{Một vật chuyển động theo quy luật $s(t)=-\dfrac{1}{3}t^3+2t^2$, trong đó $t$ (giây) là khoảng thời gian tính từ lúc vật bắt đầu chuyển động và $s$ (mét) là quãng đường vật đi được trong khoảng thời gian đó, $0\le t\le 4$. Quãng đường vật đi được trong $2$ giây đầu tiên bằng}
  {$\dfrac{16}{3}\,\mathrm{m}$.}
  {$\dfrac{19}{3}\,\mathrm{m}$.}{$\dfrac{32}{3}\,\mathrm{m}$.}{$\dfrac{13}{3}\,\mathrm{m}$.}

\cauTN[id=C1-B5-D01-NB-026,muc=NB,loigiai={Ta có $v(t)=h'(t)=30-10t$. Do đó $v(\dfrac{1}{2})=25$.}]{Một quả bóng được ném thẳng đứng lên trên từ độ cao $15$ m với vận tốc ban đầu $30$ m/s. Độ cao của quả bóng (tính bằng mét) sau $t$ giây là $h(t)=15+30t-5t^2$ (bỏ qua sức cản không khí, lấy $g=10$ m/s$^2$). Vận tốc của quả bóng tại thời điểm $t=\dfrac{1}{2}$ giây bằng}
  {$25\,\mathrm{m/s}$.}
  {$26\,\mathrm{m/s}$.}{$23\,\mathrm{m/s}$.}{$50\,\mathrm{m/s}$.}

\cauTN[id=C1-B5-D01-NB-027,muc=NB,loigiai={Vận tốc $v(t)=s'(t)=25-5t$. Gia tốc $a(t)=v'(t)=-5$.}]{Một ô tô đang chạy với vận tốc $25$ m/s thì người lái đạp phanh. Từ thời điểm đó, ô tô chuyển động chậm dần đều và quãng đường đi được sau $t$ giây là $s(t)=25t-\dfrac{5}{2}t^2$ (mét). Gia tốc của ô tô bằng}
  {$-5\,\mathrm{m/s^2}$.}
  {$-3\,\mathrm{m/s^2}$.}{$-7\,\mathrm{m/s^2}$.}{$5\,\mathrm{m/s^2}$.}

\cauTN[id=C1-B5-D01-NB-028,muc=NB,loigiai={Vận tốc $v(t)=s'(t)=24-2t$. $v(1)=22$.}]{Một ô tô đang chạy với vận tốc $24$ m/s thì người lái đạp phanh. Từ thời điểm đó, ô tô chuyển động chậm dần đều và quãng đường đi được sau $t$ giây là $s(t)=24t-t^2$ (mét). Vận tốc của ô tô sau $1$ giây kể từ lúc đạp phanh bằng}
  {$22\,\mathrm{m/s}$.}
  {$20\,\mathrm{m/s}$.}{$11\,\mathrm{m/s}$.}{$26\,\mathrm{m/s}$.}

\cauTN[id=C1-B5-D01-NB-029,muc=NB,loigiai={Vận tốc $v(t)=s'(t)=-t^{2}+16t$, gia tốc $a(t)=v'(t)=-2t+16$. $v'(t)=0\Leftrightarrow t=8$; $v$ đồng biến trên $(0;8)$, nghịch biến trên $(8;13)$ nên $v_{\max}=v(8)=64$.}]{Một vật chuyển động theo quy luật $s(t)=-\dfrac{1}{3}t^3+8t^2$, trong đó $t$ (giây) là khoảng thời gian tính từ lúc vật bắt đầu chuyển động và $s$ (mét) là quãng đường vật đi được trong khoảng thời gian đó, $0\le t\le 13$. Vận tốc lớn nhất của vật trong khoảng thời gian trên bằng}
  {$64\,\mathrm{m/s}$.}
  {$8\,\mathrm{m/s}$.}{$39\,\mathrm{m/s}$.}{$62\,\mathrm{m/s}$.}

\cauTN[id=C1-B5-D01-NB-030,muc=NB,loigiai={Ta có $h(1)=10$.}]{Một quả bóng được ném thẳng đứng lên trên từ độ cao $0$ m với vận tốc ban đầu $15$ m/s. Độ cao của quả bóng (tính bằng mét) sau $t$ giây là $h(t)=15t-5t^2$ (bỏ qua sức cản không khí, lấy $g=10$ m/s$^2$). Độ cao của quả bóng tại thời điểm $t=1$ giây bằng}
  {$10\,\mathrm{m}$.}
  {$8\,\mathrm{m}$.}{$5\,\mathrm{m}$.}{$20\,\mathrm{m}$.}

\cauTN[id=C1-B5-D01-NB-031,muc=NB,loigiai={Ta có $v(t)=h'(t)=40-10t$. $h'(t)=0\Leftrightarrow t=4$.}]{Một quả bóng được ném thẳng đứng lên trên từ độ cao $10$ m với vận tốc ban đầu $40$ m/s. Độ cao của quả bóng (tính bằng mét) sau $t$ giây là $h(t)=10+40t-5t^2$ (bỏ qua sức cản không khí, lấy $g=10$ m/s$^2$). Thời điểm quả bóng đạt độ cao lớn nhất là}
  {$4\,\mathrm{s}$.}
  {$8\,\mathrm{s}$.}{$6\,\mathrm{s}$.}{$2\,\mathrm{s}$.}

\cauTN[id=C1-B5-D01-NB-032,muc=NB,loigiai={Ta có $v(t)=h'(t)=30-10t$. Do đó $v(1)=20$.}]{Một quả bóng được ném thẳng đứng lên trên từ độ cao $15$ m với vận tốc ban đầu $30$ m/s. Độ cao của quả bóng (tính bằng mét) sau $t$ giây là $h(t)=15+30t-5t^2$ (bỏ qua sức cản không khí, lấy $g=10$ m/s$^2$). Vận tốc của quả bóng tại thời điểm $t=1$ giây bằng}
  {$20\,\mathrm{m/s}$.}
  {$24\,\mathrm{m/s}$.}{$22\,\mathrm{m/s}$.}{$16\,\mathrm{m/s}$.}

\cauTN[id=C1-B5-D01-NB-033,muc=NB,loigiai={Ta có $v(t)=h'(t)=10-10t$. $h'(t)=0\Leftrightarrow t=1$, $h$ đạt giá trị lớn nhất bằng $h(1)=25$.}]{Một quả bóng được ném thẳng đứng lên trên từ độ cao $20$ m với vận tốc ban đầu $10$ m/s. Độ cao của quả bóng (tính bằng mét) sau $t$ giây là $h(t)=20+10t-5t^2$ (bỏ qua sức cản không khí, lấy $g=10$ m/s$^2$). Độ cao lớn nhất mà quả bóng đạt được bằng}
  {$25\,\mathrm{m}$.}
  {$26\,\mathrm{m}$.}{$27\,\mathrm{m}$.}{$50\,\mathrm{m}$.}

\cauTN[id=C1-B5-D01-NB-034,muc=NB,loigiai={Vận tốc $v(t)=s'(t)=12-6t$. Ô tô dừng tại $t=2$, quãng đường $s(2)=12$.}]{Một ô tô đang chạy với vận tốc $12$ m/s thì người lái đạp phanh. Từ thời điểm đó, ô tô chuyển động chậm dần đều và quãng đường đi được sau $t$ giây là $s(t)=12t-3t^2$ (mét). Quãng đường ô tô đi được từ lúc đạp phanh đến khi dừng hẳn bằng}
  {$12\,\mathrm{m}$.}
  {$24\,\mathrm{m}$.}{$6\,\mathrm{m}$.}{$14\,\mathrm{m}$.}

\cauTN[id=C1-B5-D01-NB-035,muc=NB,loigiai={Vận tốc $v(t)=s'(t)=20-6t$. $v(1)=14$.}]{Một ô tô đang chạy với vận tốc $20$ m/s thì người lái đạp phanh. Từ thời điểm đó, ô tô chuyển động chậm dần đều và quãng đường đi được sau $t$ giây là $s(t)=20t-3t^2$ (mét). Vận tốc của ô tô sau $1$ giây kể từ lúc đạp phanh bằng}
  {$14\,\mathrm{m/s}$.}
  {$7\,\mathrm{m/s}$.}{$16\,\mathrm{m/s}$.}{$10\,\mathrm{m/s}$.}

\cauTN[id=C1-B5-D01-NB-036,muc=NB,loigiai={Ta có $v(t)=h'(t)=18-10t$. $h'(t)=0\Leftrightarrow t=\dfrac{9}{5}$.}]{Một quả bóng được ném thẳng đứng lên trên từ độ cao $15$ m với vận tốc ban đầu $18$ m/s. Độ cao của quả bóng (tính bằng mét) sau $t$ giây là $h(t)=15+18t-5t^2$ (bỏ qua sức cản không khí, lấy $g=10$ m/s$^2$). Thời điểm quả bóng đạt độ cao lớn nhất là}
  {$\dfrac{9}{5}\,\mathrm{s}$.}
  {$\dfrac{9}{10}\,\mathrm{s}$.}{$\dfrac{5}{9}\,\mathrm{s}$.}{$\dfrac{14}{5}\,\mathrm{s}$.}

\cauDS[id=C1-B5-D01-NB-037,muc=NB,loigiai={Vận tốc $v(t)=s'(t)=-3t^{2}+30t+3$, gia tốc $a(t)=v'(t)=-6t+30$. $v'(t)=0\Leftrightarrow t=5$; $v$ đồng biến trên $(0;5)$, nghịch biến trên $(5;10)$ nên $v_{\max}=v(5)=78$. Quãng đường $s(3)=117$. Vận tốc $v(t)=s'(t)=-3t^{2}+30t+3$, gia tốc $a(t)=v'(t)=-6t+30$. Do đó $v(3)=66$. Vận tốc $v(t)=s'(t)=-3t^{2}+30t+3$, gia tốc $a(t)=v'(t)=-6t+30$. $v'(t)=0\Leftrightarrow t=5$ và $v'$ đổi dấu từ dương sang âm nên vận tốc lớn nhất tại $t=5$. a) Đúng. b) Đúng. c) Sai: giá trị đúng là $117\,\mathrm{m}$. d) Sai: giá trị đúng là $78\,\mathrm{m/s}$.}]{Một vật chuyển động theo quy luật $s(t)=-\dfrac{3}{3}t^3+15t^2+3t$, trong đó $t$ (giây) là khoảng thời gian tính từ lúc vật bắt đầu chuyển động và $s$ (mét) là quãng đường vật đi được trong khoảng thời gian đó, $0\le t\le 10$. Xét tính đúng sai của các mệnh đề sau:}{%
  \y{Vận tốc của vật tại thời điểm $t=3$ giây là $66\,\mathrm{m/s}$.}{d}
  \y{Vật đạt vận tốc lớn nhất tại thời điểm $t=5\,\mathrm{s}$.}{d}
  \y{Quãng đường vật đi được trong $3$ giây đầu là $67\,\mathrm{m}$.}{s}
  \y{Vận tốc lớn nhất của vật là $39\,\mathrm{m/s}$.}{s}}

\cauDS[id=C1-B5-D01-NB-038,muc=NB,loigiai={Ta có $h(\dfrac{1}{2})=\dfrac{31}{4}$. Ta có $v(t)=h'(t)=18-10t$. $h'(t)=0\Leftrightarrow t=\dfrac{9}{5}$. Ta có $v(t)=h'(t)=18-10t$. $h'(t)=0\Leftrightarrow t=\dfrac{9}{5}$, $h$ đạt giá trị lớn nhất bằng $h(\dfrac{9}{5})=\dfrac{81}{5}$. Ta có $v(t)=h'(t)=18-10t$. Do đó $v(\dfrac{1}{2})=13$. a) Sai: giá trị đúng là $\dfrac{81}{5}\,\mathrm{m}$. b) Đúng. c) Đúng. d) Đúng.}]{Một quả bóng được ném thẳng đứng lên trên từ độ cao $0$ m với vận tốc ban đầu $18$ m/s. Độ cao của quả bóng (tính bằng mét) sau $t$ giây là $h(t)=18t-5t^2$ (bỏ qua sức cản không khí, lấy $g=10$ m/s$^2$). Xét tính đúng sai của các mệnh đề sau:}{%
  \y{Độ cao lớn nhất của quả bóng là $\dfrac{81}{10}\,\mathrm{m}$.}{s}
  \y{Quả bóng đạt độ cao lớn nhất tại thời điểm $t=\dfrac{9}{5}\,\mathrm{s}$.}{d}
  \y{Sau $\dfrac{1}{2}$ giây quả bóng ở độ cao $\dfrac{31}{4}\,\mathrm{m}$.}{d}
  \y{Vận tốc của quả bóng tại thời điểm $t=\dfrac{1}{2}$ giây là $13\,\mathrm{m/s}$.}{d}}

\cauDS[id=C1-B5-D01-NB-039,muc=NB,loigiai={Vận tốc $v(t)=s'(t)=-t^{2}+8t+3$, gia tốc $a(t)=v'(t)=-2t+8$. Do đó $a(2)=4$. Vận tốc $v(t)=s'(t)=-t^{2}+8t+3$, gia tốc $a(t)=v'(t)=-2t+8$. Do đó $v(2)=15$. Quãng đường $s(2)=\dfrac{58}{3}$. Vận tốc $v(t)=s'(t)=-t^{2}+8t+3$, gia tốc $a(t)=v'(t)=-2t+8$. $v'(t)=0\Leftrightarrow t=4$ và $v'$ đổi dấu từ dương sang âm nên vận tốc lớn nhất tại $t=4$. a) Đúng. b) Đúng. c) Đúng. d) Sai: giá trị đúng là $4\,\mathrm{m/s^2}$.}]{Một vật chuyển động theo quy luật $s(t)=-\dfrac{1}{3}t^3+4t^2+3t$, trong đó $t$ (giây) là khoảng thời gian tính từ lúc vật bắt đầu chuyển động và $s$ (mét) là quãng đường vật đi được trong khoảng thời gian đó, $0\le t\le 7$. Xét tính đúng sai của các mệnh đề sau:}{%
  \y{Vận tốc của vật tại thời điểm $t=2$ giây là $15\,\mathrm{m/s}$.}{d}
  \y{Vật đạt vận tốc lớn nhất tại thời điểm $t=4\,\mathrm{s}$.}{d}
  \y{Quãng đường vật đi được trong $2$ giây đầu là $\dfrac{58}{3}\,\mathrm{m}$.}{d}
  \y{Gia tốc của vật tại thời điểm $t=2$ giây là $5\,\mathrm{m/s^2}$.}{s}}

\cauDS[id=C1-B5-D01-NB-040,muc=NB,loigiai={Vận tốc $v(t)=s'(t)=-t^{2}+12t+3$, gia tốc $a(t)=v'(t)=-2t+12$. $v'(t)=0\Leftrightarrow t=6$; $v$ đồng biến trên $(0;6)$, nghịch biến trên $(6;10)$ nên $v_{\max}=v(6)=39$. Vận tốc $v(t)=s'(t)=-t^{2}+12t+3$, gia tốc $a(t)=v'(t)=-2t+12$. $v'(t)=0\Leftrightarrow t=6$ và $v'$ đổi dấu từ dương sang âm nên vận tốc lớn nhất tại $t=6$. Quãng đường $s(8)=\dfrac{712}{3}$. Vận tốc $v(t)=s'(t)=-t^{2}+12t+3$, gia tốc $a(t)=v'(t)=-2t+12$. Do đó $v(8)=35$. a) Sai: giá trị đúng là $35\,\mathrm{m/s}$. b) Đúng. c) Sai: giá trị đúng là $\dfrac{712}{3}\,\mathrm{m}$. d) Đúng.}]{Một vật chuyển động theo quy luật $s(t)=-\dfrac{1}{3}t^3+6t^2+3t$, trong đó $t$ (giây) là khoảng thời gian tính từ lúc vật bắt đầu chuyển động và $s$ (mét) là quãng đường vật đi được trong khoảng thời gian đó, $0\le t\le 10$. Xét tính đúng sai của các mệnh đề sau:}{%
  \y{Vận tốc của vật tại thời điểm $t=8$ giây là $70\,\mathrm{m/s}$.}{s}
  \y{Vận tốc lớn nhất của vật là $39\,\mathrm{m/s}$.}{d}
  \y{Quãng đường vật đi được trong $8$ giây đầu là $712\,\mathrm{m}$.}{s}
  \y{Vật đạt vận tốc lớn nhất tại thời điểm $t=6\,\mathrm{s}$.}{d}}

\cauDS[id=C1-B5-D01-NB-041,muc=NB,loigiai={Vận tốc $v(t)=s'(t)=-t^{2}+10t$, gia tốc $a(t)=v'(t)=-2t+10$. $v'(t)=0\Leftrightarrow t=5$; $v$ đồng biến trên $(0;5)$, nghịch biến trên $(5;11)$ nên $v_{\max}=v(5)=25$. Vận tốc $v(t)=s'(t)=-t^{2}+10t$, gia tốc $a(t)=v'(t)=-2t+10$. Do đó $a(8)=-6$. Vận tốc $v(t)=s'(t)=-t^{2}+10t$, gia tốc $a(t)=v'(t)=-2t+10$. Do đó $v(8)=16$. Vận tốc $v(t)=s'(t)=-t^{2}+10t$, gia tốc $a(t)=v'(t)=-2t+10$. $v'(t)=0\Leftrightarrow t=5$ và $v'$ đổi dấu từ dương sang âm nên vận tốc lớn nhất tại $t=5$. a) Đúng. b) Đúng. c) Đúng. d) Sai: giá trị đúng là $5\,\mathrm{s}$.}]{Một vật chuyển động theo quy luật $s(t)=-\dfrac{1}{3}t^3+5t^2$, trong đó $t$ (giây) là khoảng thời gian tính từ lúc vật bắt đầu chuyển động và $s$ (mét) là quãng đường vật đi được trong khoảng thời gian đó, $0\le t\le 11$. Xét tính đúng sai của các mệnh đề sau:}{%
  \y{Vận tốc lớn nhất của vật là $25\,\mathrm{m/s}$.}{d}
  \y{Gia tốc của vật tại thời điểm $t=8$ giây là $-6\,\mathrm{m/s^2}$.}{d}
  \y{Vận tốc của vật tại thời điểm $t=8$ giây là $16\,\mathrm{m/s}$.}{d}
  \y{Vật đạt vận tốc lớn nhất tại thời điểm $t=10\,\mathrm{s}$.}{s}}

\cauDS[id=C1-B5-D01-NB-042,muc=NB,loigiai={Vận tốc $v(t)=s'(t)=-2t^{2}+20t$, gia tốc $a(t)=v'(t)=-4t+20$. Do đó $a(5)=0$. Vận tốc $v(t)=s'(t)=-2t^{2}+20t$, gia tốc $a(t)=v'(t)=-4t+20$. $v'(t)=0\Leftrightarrow t=5$ và $v'$ đổi dấu từ dương sang âm nên vận tốc lớn nhất tại $t=5$. Quãng đường $s(5)=\dfrac{500}{3}$. Vận tốc $v(t)=s'(t)=-2t^{2}+20t$, gia tốc $a(t)=v'(t)=-4t+20$. $v'(t)=0\Leftrightarrow t=5$; $v$ đồng biến trên $(0;5)$, nghịch biến trên $(5;10)$ nên $v_{\max}=v(5)=50$. a) Sai: giá trị đúng là $5\,\mathrm{s}$. b) Đúng. c) Đúng. d) Đúng.}]{Một vật chuyển động theo quy luật $s(t)=-\dfrac{2}{3}t^3+10t^2$, trong đó $t$ (giây) là khoảng thời gian tính từ lúc vật bắt đầu chuyển động và $s$ (mét) là quãng đường vật đi được trong khoảng thời gian đó, $0\le t\le 10$. Xét tính đúng sai của các mệnh đề sau:}{%
  \y{Vật đạt vận tốc lớn nhất tại thời điểm $t=9\,\mathrm{s}$.}{s}
  \y{Gia tốc của vật tại thời điểm $t=5$ giây là $0\,\mathrm{m/s^2}$.}{d}
  \y{Vận tốc lớn nhất của vật là $50\,\mathrm{m/s}$.}{d}
  \y{Quãng đường vật đi được trong $5$ giây đầu là $\dfrac{500}{3}\,\mathrm{m}$.}{d}}

\cauDS[id=C1-B5-D01-NB-043,muc=NB,loigiai={Vận tốc $v(t)=s'(t)=30-4t$. Gia tốc $a(t)=v'(t)=-4$. Vận tốc $v(t)=s'(t)=30-4t$. Ô tô dừng tại $t=\dfrac{15}{2}$, quãng đường $s(\dfrac{15}{2})=\dfrac{225}{2}$. Vận tốc $v(t)=s'(t)=30-4t$. $v(2)=22$. Vận tốc $v(t)=s'(t)=30-4t$. Ô tô dừng khi $v(t)=0\Leftrightarrow t=\dfrac{15}{2}$. a) Sai: giá trị đúng là $\dfrac{15}{2}\,\mathrm{s}$. b) Đúng. c) Đúng. d) Đúng.}]{Một ô tô đang chạy với vận tốc $30$ m/s thì người lái đạp phanh. Từ thời điểm đó, ô tô chuyển động chậm dần đều và quãng đường đi được sau $t$ giây là $s(t)=30t-2t^2$ (mét). Xét tính đúng sai của các mệnh đề sau:}{%
  \y{Ô tô dừng hẳn sau $\dfrac{17}{2}\,\mathrm{s}$ kể từ lúc đạp phanh.}{s}
  \y{Gia tốc của ô tô là $-4\,\mathrm{m/s^2}$.}{d}
  \y{Quãng đường ô tô đi được từ lúc đạp phanh đến khi dừng là $\dfrac{225}{2}\,\mathrm{m}$.}{d}
  \y{Sau $2$ giây kể từ lúc đạp phanh, vận tốc ô tô là $22\,\mathrm{m/s}$.}{d}}

\cauDS[id=C1-B5-D01-NB-044,muc=NB,loigiai={Ta có $h(2)=62$. Ta có $v(t)=h'(t)=40-10t$. $h'(t)=0\Leftrightarrow t=4$. Ta có $v(t)=h'(t)=40-10t$. $h'(t)=0\Leftrightarrow t=4$, $h$ đạt giá trị lớn nhất bằng $h(4)=82$. Ta có $v(t)=h'(t)=40-10t$. Do đó $v(2)=20$. a) Đúng. b) Sai: giá trị đúng là $82\,\mathrm{m}$. c) Sai: giá trị đúng là $4\,\mathrm{s}$. d) Đúng.}]{Một quả bóng được ném thẳng đứng lên trên từ độ cao $2$ m với vận tốc ban đầu $40$ m/s. Độ cao của quả bóng (tính bằng mét) sau $t$ giây là $h(t)=2+40t-5t^2$ (bỏ qua sức cản không khí, lấy $g=10$ m/s$^2$). Xét tính đúng sai của các mệnh đề sau:}{%
  \y{Vận tốc của quả bóng tại thời điểm $t=2$ giây là $20\,\mathrm{m/s}$.}{d}
  \y{Độ cao lớn nhất của quả bóng là $92\,\mathrm{m}$.}{s}
  \y{Quả bóng đạt độ cao lớn nhất tại thời điểm $t=8\,\mathrm{s}$.}{s}
  \y{Sau $2$ giây quả bóng ở độ cao $62\,\mathrm{m}$.}{d}}

\cauDS[id=C1-B5-D01-NB-045,muc=NB,loigiai={Vận tốc $v(t)=s'(t)=25-2t$. Ô tô dừng tại $t=\dfrac{25}{2}$, quãng đường $s(\dfrac{25}{2})=\dfrac{625}{4}$. Vận tốc $v(t)=s'(t)=25-2t$. Gia tốc $a(t)=v'(t)=-2$. Vận tốc $v(t)=s'(t)=25-2t$. $v(2)=21$. Vận tốc $v(t)=s'(t)=25-2t$. Ô tô dừng khi $v(t)=0\Leftrightarrow t=\dfrac{25}{2}$. a) Đúng. b) Đúng. c) Sai: giá trị đúng là $21\,\mathrm{m/s}$. d) Đúng.}]{Một ô tô đang chạy với vận tốc $25$ m/s thì người lái đạp phanh. Từ thời điểm đó, ô tô chuyển động chậm dần đều và quãng đường đi được sau $t$ giây là $s(t)=25t-t^2$ (mét). Xét tính đúng sai của các mệnh đề sau:}{%
  \y{Gia tốc của ô tô là $-2\,\mathrm{m/s^2}$.}{d}
  \y{Quãng đường ô tô đi được từ lúc đạp phanh đến khi dừng là $\dfrac{625}{4}\,\mathrm{m}$.}{d}
  \y{Sau $2$ giây kể từ lúc đạp phanh, vận tốc ô tô là $42\,\mathrm{m/s}$.}{s}
  \y{Ô tô dừng hẳn sau $\dfrac{25}{2}\,\mathrm{s}$ kể từ lúc đạp phanh.}{d}}

\cauDS[id=C1-B5-D01-NB-046,muc=NB,loigiai={Ta có $h(\dfrac{1}{2})=\dfrac{125}{4}$. Ta có $v(t)=h'(t)=25-10t$. Do đó $v(\dfrac{1}{2})=20$. Ta có $v(t)=h'(t)=25-10t$. $h'(t)=0\Leftrightarrow t=\dfrac{5}{2}$, $h$ đạt giá trị lớn nhất bằng $h(\dfrac{5}{2})=\dfrac{205}{4}$. Ta có $v(t)=h'(t)=25-10t$. $h'(t)=0\Leftrightarrow t=\dfrac{5}{2}$. a) Sai: giá trị đúng là $\dfrac{5}{2}\,\mathrm{s}$. b) Đúng. c) Sai: giá trị đúng là $\dfrac{125}{4}\,\mathrm{m}$. d) Đúng.}]{Một quả bóng được ném thẳng đứng lên trên từ độ cao $20$ m với vận tốc ban đầu $25$ m/s. Độ cao của quả bóng (tính bằng mét) sau $t$ giây là $h(t)=20+25t-5t^2$ (bỏ qua sức cản không khí, lấy $g=10$ m/s$^2$). Xét tính đúng sai của các mệnh đề sau:}{%
  \y{Quả bóng đạt độ cao lớn nhất tại thời điểm $t=\dfrac{5}{4}\,\mathrm{s}$.}{s}
  \y{Vận tốc của quả bóng tại thời điểm $t=\dfrac{1}{2}$ giây là $20\,\mathrm{m/s}$.}{d}
  \y{Sau $\dfrac{1}{2}$ giây quả bóng ở độ cao $\dfrac{125}{8}\,\mathrm{m}$.}{s}
  \y{Độ cao lớn nhất của quả bóng là $\dfrac{205}{4}\,\mathrm{m}$.}{d}}

\cauDS[id=C1-B5-D01-NB-047,muc=NB,loigiai={Vận tốc $v(t)=s'(t)=-2t^{2}+8t+3$, gia tốc $a(t)=v'(t)=-4t+8$. Do đó $v(1)=9$. Vận tốc $v(t)=s'(t)=-2t^{2}+8t+3$, gia tốc $a(t)=v'(t)=-4t+8$. Do đó $a(1)=4$. Vận tốc $v(t)=s'(t)=-2t^{2}+8t+3$, gia tốc $a(t)=v'(t)=-4t+8$. $v'(t)=0\Leftrightarrow t=2$; $v$ đồng biến trên $(0;2)$, nghịch biến trên $(2;7)$ nên $v_{\max}=v(2)=11$. Vận tốc $v(t)=s'(t)=-2t^{2}+8t+3$, gia tốc $a(t)=v'(t)=-4t+8$. $v'(t)=0\Leftrightarrow t=2$ và $v'$ đổi dấu từ dương sang âm nên vận tốc lớn nhất tại $t=2$. a) Sai: giá trị đúng là $2\,\mathrm{s}$. b) Sai: giá trị đúng là $9\,\mathrm{m/s}$. c) Sai: giá trị đúng là $11\,\mathrm{m/s}$. d) Đúng.}]{Một vật chuyển động theo quy luật $s(t)=-\dfrac{2}{3}t^3+4t^2+3t$, trong đó $t$ (giây) là khoảng thời gian tính từ lúc vật bắt đầu chuyển động và $s$ (mét) là quãng đường vật đi được trong khoảng thời gian đó, $0\le t\le 7$. Xét tính đúng sai của các mệnh đề sau:}{%
  \y{Vật đạt vận tốc lớn nhất tại thời điểm $t=3\,\mathrm{s}$.}{s}
  \y{Vận tốc của vật tại thời điểm $t=1$ giây là $10\,\mathrm{m/s}$.}{s}
  \y{Vận tốc lớn nhất của vật là $22\,\mathrm{m/s}$.}{s}
  \y{Gia tốc của vật tại thời điểm $t=1$ giây là $4\,\mathrm{m/s^2}$.}{d}}

\cauDS[id=C1-B5-D01-NB-048,muc=NB,loigiai={Ta có $h(2)=40$. Ta có $v(t)=h'(t)=25-10t$. $h'(t)=0\Leftrightarrow t=\dfrac{5}{2}$. Ta có $v(t)=h'(t)=25-10t$. Do đó $v(2)=5$. Ta có $v(t)=h'(t)=25-10t$. $h'(t)=0\Leftrightarrow t=\dfrac{5}{2}$, $h$ đạt giá trị lớn nhất bằng $h(\dfrac{5}{2})=\dfrac{165}{4}$. a) Đúng. b) Sai: giá trị đúng là $\dfrac{165}{4}\,\mathrm{m}$. c) Sai: giá trị đúng là $5\,\mathrm{m/s}$. d) Sai: giá trị đúng là $\dfrac{5}{2}\,\mathrm{s}$.}]{Một quả bóng được ném thẳng đứng lên trên từ độ cao $10$ m với vận tốc ban đầu $25$ m/s. Độ cao của quả bóng (tính bằng mét) sau $t$ giây là $h(t)=10+25t-5t^2$ (bỏ qua sức cản không khí, lấy $g=10$ m/s$^2$). Xét tính đúng sai của các mệnh đề sau:}{%
  \y{Sau $2$ giây quả bóng ở độ cao $40\,\mathrm{m}$.}{d}
  \y{Độ cao lớn nhất của quả bóng là $\dfrac{165}{8}\,\mathrm{m}$.}{s}
  \y{Vận tốc của quả bóng tại thời điểm $t=2$ giây là $9\,\mathrm{m/s}$.}{s}
  \y{Quả bóng đạt độ cao lớn nhất tại thời điểm $t=\dfrac{5}{4}\,\mathrm{s}$.}{s}}

\cauTLN[id=C1-B5-D01-NB-049,muc=NB,loigiai={Vận tốc $v(t)=s'(t)=25-2t$. Ô tô dừng tại $t=\dfrac{25}{2}$, quãng đường $s(\dfrac{25}{2})=\dfrac{625}{4}$. Kết quả: $156$.}]{Một ô tô đang chạy với vận tốc $25$ m/s thì người lái đạp phanh. Từ thời điểm đó, ô tô chuyển động chậm dần đều và quãng đường đi được sau $t$ giây là $s(t)=25t-t^2$ (mét). Tìm quãng đường ô tô đi được từ lúc đạp phanh đến khi dừng hẳn (đơn vị: $\mathrm{m}$) (làm tròn kết quả đến hàng đơn vị).}{156}

\cauTLN[id=C1-B5-D01-NB-050,muc=NB,loigiai={Vận tốc $v(t)=s'(t)=-3t^{2}+36t$, gia tốc $a(t)=v'(t)=-6t+36$. Do đó $a(2)=24$. Kết quả: $24$.}]{Một vật chuyển động theo quy luật $s(t)=-\dfrac{3}{3}t^3+18t^2$, trong đó $t$ (giây) là khoảng thời gian tính từ lúc vật bắt đầu chuyển động và $s$ (mét) là quãng đường vật đi được trong khoảng thời gian đó, $0\le t\le 12$. Tìm gia tốc của vật tại thời điểm $t=2$ giây (đơn vị: $\mathrm{m/s^2}$).}{24}

\cauTLN[id=C1-B5-D01-NB-051,muc=NB,loigiai={Vận tốc $v(t)=s'(t)=30-4t$. $v(2)=22$. Kết quả: $22$.}]{Một ô tô đang chạy với vận tốc $30$ m/s thì người lái đạp phanh. Từ thời điểm đó, ô tô chuyển động chậm dần đều và quãng đường đi được sau $t$ giây là $s(t)=30t-2t^2$ (mét). Tìm vận tốc của ô tô sau $2$ giây kể từ lúc đạp phanh (đơn vị: $\mathrm{m/s}$).}{22}

\cauTLN[id=C1-B5-D01-NB-052,muc=NB,loigiai={Ta có $v(t)=h'(t)=30-10t$. $h'(t)=0\Leftrightarrow t=3$. Kết quả: $3$.}]{Một quả bóng được ném thẳng đứng lên trên từ độ cao $1$ m với vận tốc ban đầu $30$ m/s. Độ cao của quả bóng (tính bằng mét) sau $t$ giây là $h(t)=1+30t-5t^2$ (bỏ qua sức cản không khí, lấy $g=10$ m/s$^2$). Tìm thời điểm quả bóng đạt độ cao lớn nhất (đơn vị: $\mathrm{s}$).}{3}

\cauTLN[id=C1-B5-D01-NB-053,muc=NB,loigiai={Ta có $v(t)=h'(t)=18-10t$. $h'(t)=0\Leftrightarrow t=\dfrac{9}{5}$, $h$ đạt giá trị lớn nhất bằng $h(\dfrac{9}{5})=\dfrac{91}{5}$. Kết quả: $18{,}2$.}]{Một quả bóng được ném thẳng đứng lên trên từ độ cao $2$ m với vận tốc ban đầu $18$ m/s. Độ cao của quả bóng (tính bằng mét) sau $t$ giây là $h(t)=2+18t-5t^2$ (bỏ qua sức cản không khí, lấy $g=10$ m/s$^2$). Tìm độ cao lớn nhất mà quả bóng đạt được (đơn vị: $\mathrm{m}$).}{18,2}

\cauTLN[id=C1-B5-D01-NB-054,muc=NB,loigiai={Ta có $v(t)=h'(t)=12-10t$. Do đó $v(\dfrac{1}{2})=7$. Kết quả: $7$.}]{Một quả bóng được ném thẳng đứng lên trên từ độ cao $10$ m với vận tốc ban đầu $12$ m/s. Độ cao của quả bóng (tính bằng mét) sau $t$ giây là $h(t)=10+12t-5t^2$ (bỏ qua sức cản không khí, lấy $g=10$ m/s$^2$). Tìm vận tốc của quả bóng tại thời điểm $t=\dfrac{1}{2}$ giây (đơn vị: $\mathrm{m/s}$).}{7}

\cauTLN[id=C1-B5-D01-NB-055,muc=NB,loigiai={Vận tốc $v(t)=s'(t)=20-6t$. Gia tốc $a(t)=v'(t)=-6$. Kết quả: $-6$.}]{Một ô tô đang chạy với vận tốc $20$ m/s thì người lái đạp phanh. Từ thời điểm đó, ô tô chuyển động chậm dần đều và quãng đường đi được sau $t$ giây là $s(t)=20t-3t^2$ (mét). Tìm gia tốc của ô tô (đơn vị: $\mathrm{m/s^2}$).}{-6}

\cauTLN[id=C1-B5-D01-NB-056,muc=NB,loigiai={Vận tốc $v(t)=s'(t)=-3t^{2}+48t$, gia tốc $a(t)=v'(t)=-6t+48$. Do đó $a(11)=-18$. Kết quả: $-18$.}]{Một vật chuyển động theo quy luật $s(t)=-\dfrac{3}{3}t^3+24t^2$, trong đó $t$ (giây) là khoảng thời gian tính từ lúc vật bắt đầu chuyển động và $s$ (mét) là quãng đường vật đi được trong khoảng thời gian đó, $0\le t\le 16$. Tìm gia tốc của vật tại thời điểm $t=11$ giây (đơn vị: $\mathrm{m/s^2}$).}{-18}

\cauTLN[id=C1-B5-D01-NB-057,muc=NB,loigiai={Vận tốc $v(t)=s'(t)=10-3t$. Gia tốc $a(t)=v'(t)=-3$. Kết quả: $-3$.}]{Một ô tô đang chạy với vận tốc $10$ m/s thì người lái đạp phanh. Từ thời điểm đó, ô tô chuyển động chậm dần đều và quãng đường đi được sau $t$ giây là $s(t)=10t-\dfrac{3}{2}t^2$ (mét). Tìm gia tốc của ô tô (đơn vị: $\mathrm{m/s^2}$).}{-3}

\cauTLN[id=C1-B5-D01-NB-058,muc=NB,loigiai={Vận tốc $v(t)=s'(t)=-2t^{2}+12t$, gia tốc $a(t)=v'(t)=-4t+12$. $v'(t)=0\Leftrightarrow t=3$; $v$ đồng biến trên $(0;3)$, nghịch biến trên $(3;6)$ nên $v_{\max}=v(3)=18$. Kết quả: $18$.}]{Một vật chuyển động theo quy luật $s(t)=-\dfrac{2}{3}t^3+6t^2$, trong đó $t$ (giây) là khoảng thời gian tính từ lúc vật bắt đầu chuyển động và $s$ (mét) là quãng đường vật đi được trong khoảng thời gian đó, $0\le t\le 6$. Tìm vận tốc lớn nhất của vật trong khoảng thời gian trên (đơn vị: $\mathrm{m/s}$).}{18}

\cauTLN[id=C1-B5-D01-NB-059,muc=NB,loigiai={Ta có $v(t)=h'(t)=40-10t$. $h'(t)=0\Leftrightarrow t=4$, $h$ đạt giá trị lớn nhất bằng $h(4)=85$. Kết quả: $85$.}]{Một quả bóng được ném thẳng đứng lên trên từ độ cao $5$ m với vận tốc ban đầu $40$ m/s. Độ cao của quả bóng (tính bằng mét) sau $t$ giây là $h(t)=5+40t-5t^2$ (bỏ qua sức cản không khí, lấy $g=10$ m/s$^2$). Tìm độ cao lớn nhất mà quả bóng đạt được (đơn vị: $\mathrm{m}$).}{85}

\cauTLN[id=C1-B5-D01-NB-060,muc=NB,loigiai={Ta có $v(t)=h'(t)=18-10t$. $h'(t)=0\Leftrightarrow t=\dfrac{9}{5}$. Kết quả: $1{,}8$.}]{Một quả bóng được ném thẳng đứng lên trên từ độ cao $5$ m với vận tốc ban đầu $18$ m/s. Độ cao của quả bóng (tính bằng mét) sau $t$ giây là $h(t)=5+18t-5t^2$ (bỏ qua sức cản không khí, lấy $g=10$ m/s$^2$). Tìm thời điểm quả bóng đạt độ cao lớn nhất (đơn vị: $\mathrm{s}$).}{1,8}

\cauTN[id=C1-B5-D01-TH-001,muc=TH,loigiai={Ta có $v(t)=h'(t)=30-10t$. $h'(t)=0\Leftrightarrow t=3$, $h$ đạt giá trị lớn nhất bằng $h(3)=45$.}]{Một quả bóng được ném thẳng đứng lên trên từ độ cao $0$ m với vận tốc ban đầu $30$ m/s. Độ cao của quả bóng (tính bằng mét) sau $t$ giây là $h(t)=30t-5t^2$ (bỏ qua sức cản không khí, lấy $g=10$ m/s$^2$). Độ cao lớn nhất mà quả bóng đạt được bằng}
  {$45\,\mathrm{m}$.}
  {$40\,\mathrm{m}$.}{$55\,\mathrm{m}$.}{$90\,\mathrm{m}$.}

\cauTN[id=C1-B5-D01-TH-002,muc=TH,loigiai={Vận tốc $v(t)=s'(t)=-2t^{2}+8t$, gia tốc $a(t)=v'(t)=-4t+8$. $v'(t)=0\Leftrightarrow t=2$ và $v'$ đổi dấu từ dương sang âm nên vận tốc lớn nhất tại $t=2$.}]{Một vật chuyển động theo quy luật $s(t)=-\dfrac{2}{3}t^3+4t^2$, trong đó $t$ (giây) là khoảng thời gian tính từ lúc vật bắt đầu chuyển động và $s$ (mét) là quãng đường vật đi được trong khoảng thời gian đó, $0\le t\le 4$. Thời điểm vật đạt vận tốc lớn nhất là}
  {$2\,\mathrm{s}$.}
  {$8\,\mathrm{s}$.}{$4\,\mathrm{s}$.}{$1\,\mathrm{s}$.}

\cauTN[id=C1-B5-D01-TH-003,muc=TH,loigiai={Vận tốc $v(t)=s'(t)=-2t^{2}+12t+2$, gia tốc $a(t)=v'(t)=-4t+12$. $v'(t)=0\Leftrightarrow t=3$; $v$ đồng biến trên $(0;3)$, nghịch biến trên $(3;5)$ nên $v_{\max}=v(3)=20$.}]{Một vật chuyển động theo quy luật $s(t)=-\dfrac{2}{3}t^3+6t^2+2t$, trong đó $t$ (giây) là khoảng thời gian tính từ lúc vật bắt đầu chuyển động và $s$ (mét) là quãng đường vật đi được trong khoảng thời gian đó, $0\le t\le 5$. Vận tốc lớn nhất của vật trong khoảng thời gian trên bằng}
  {$20\,\mathrm{m/s}$.}
  {$3\,\mathrm{m/s}$.}{$12\,\mathrm{m/s}$.}{$18\,\mathrm{m/s}$.}

\cauTN[id=C1-B5-D01-TH-004,muc=TH,loigiai={Ta có $v(t)=h'(t)=18-10t$. Do đó $v(2)=-2$.}]{Một quả bóng được ném thẳng đứng lên trên từ độ cao $15$ m với vận tốc ban đầu $18$ m/s. Độ cao của quả bóng (tính bằng mét) sau $t$ giây là $h(t)=15+18t-5t^2$ (bỏ qua sức cản không khí, lấy $g=10$ m/s$^2$). Vận tốc của quả bóng tại thời điểm $t=2$ giây bằng}
  {$-2\,\mathrm{m/s}$.}
  {$2\,\mathrm{m/s}$.}{$-4\,\mathrm{m/s}$.}{$-1\,\mathrm{m/s}$.}

\cauTN[id=C1-B5-D01-TH-005,muc=TH,loigiai={Vận tốc $v(t)=x'(t)=-\dfrac{3 \pi}{2}\sin(\frac{\pi}{2} t)$ (cm/s), gia tốc $a(t)=v'(t)=-\dfrac{3 \pi^{2}}{4}\cos(\frac{\pi}{2} t)$ (cm/s$^2$). $a(\dfrac{1}{3})=- \dfrac{3 \sqrt{3} \pi^{2}}{8}$.}]{Một vật dao động điều hòa có phương trình li độ $x(t)=3\cos(\frac{\pi}{2} t)$ (cm), với $t$ tính bằng giây. Gia tốc của vật tại thời điểm $t=\dfrac{1}{3}$ giây bằng}
  {$- \dfrac{3 \sqrt{3} \pi^{2}}{8}\,\mathrm{cm/s^2}$.}
  {$- \dfrac{\sqrt{3} \pi^{2}}{4}\,\mathrm{cm/s^2}$.}{$\dfrac{3 \sqrt{3} \pi^{2}}{8}\,\mathrm{cm/s^2}$.}{$- \dfrac{3 \sqrt{3} \pi^{2}}{16}\,\mathrm{cm/s^2}$.}

\cauTN[id=C1-B5-D01-TH-006,muc=TH,loigiai={Vận tốc $v(t)=x'(t)=-8 \pi\sin(\pi t)$ (cm/s), gia tốc $a(t)=v'(t)=-8 \pi^{2}\cos(\pi t)$ (cm/s$^2$). Vì $|\sin(\pi t)|\le1$ nên $|v(t)|\le 8 \pi$, dấu bằng xảy ra khi $\sin(\pi t)=\pm1$.}]{Một vật dao động điều hòa có phương trình li độ $x(t)=8\cos(\pi t)$ (cm), với $t$ tính bằng giây. Tốc độ cực đại của vật (giá trị lớn nhất của $|v(t)|$) bằng}
  {$8 \pi\,\mathrm{cm/s}$.}
  {$4 \pi\,\mathrm{cm/s}$.}{$-1 + 8 \pi\,\mathrm{cm/s}$.}{$12 \pi\,\mathrm{cm/s}$.}

\cauTN[id=C1-B5-D01-TH-007,muc=TH,loigiai={Vận tốc $v(t)=s'(t)=-t^{2}+8t+5$, gia tốc $a(t)=v'(t)=-2t+8$. $v'(t)=0\Leftrightarrow t=4$; $v$ đồng biến trên $(0;4)$, nghịch biến trên $(4;8)$ nên $v_{\max}=v(4)=21$.}]{Một vật chuyển động theo quy luật $s(t)=-\dfrac{1}{3}t^3+4t^2+5t$, trong đó $t$ (giây) là khoảng thời gian tính từ lúc vật bắt đầu chuyển động và $s$ (mét) là quãng đường vật đi được trong khoảng thời gian đó, $0\le t\le 8$. Vận tốc lớn nhất của vật trong khoảng thời gian trên bằng}
  {$21\,\mathrm{m/s}$.}
  {$4\,\mathrm{m/s}$.}{$5\,\mathrm{m/s}$.}{$16\,\mathrm{m/s}$.}

\cauTN[id=C1-B5-D01-TH-008,muc=TH,loigiai={Vận tốc $v(t)=s'(t)=-2t^{2}+12t+2$, gia tốc $a(t)=v'(t)=-4t+12$. $v'(t)=0\Leftrightarrow t=3$ và $v'$ đổi dấu từ dương sang âm nên vận tốc lớn nhất tại $t=3$.}]{Một vật chuyển động theo quy luật $s(t)=-\dfrac{2}{3}t^3+6t^2+2t$, trong đó $t$ (giây) là khoảng thời gian tính từ lúc vật bắt đầu chuyển động và $s$ (mét) là quãng đường vật đi được trong khoảng thời gian đó, $0\le t\le 6$. Thời điểm vật đạt vận tốc lớn nhất là}
  {$3\,\mathrm{s}$.}
  {$20\,\mathrm{s}$.}{$6\,\mathrm{s}$.}{$2\,\mathrm{s}$.}

\cauTN[id=C1-B5-D01-TH-009,muc=TH,loigiai={Vận tốc $v(t)=s'(t)=12-3t$. Ô tô dừng khi $v(t)=0\Leftrightarrow t=4$.}]{Một ô tô đang chạy với vận tốc $12$ m/s thì người lái đạp phanh. Từ thời điểm đó, ô tô chuyển động chậm dần đều và quãng đường đi được sau $t$ giây là $s(t)=12t-\dfrac{3}{2}t^2$ (mét). Thời gian từ lúc đạp phanh đến khi ô tô dừng hẳn bằng}
  {$4\,\mathrm{s}$.}
  {$6\,\mathrm{s}$.}{$5\,\mathrm{s}$.}{$2\,\mathrm{s}$.}

\cauTN[id=C1-B5-D01-TH-010,muc=TH,loigiai={Vận tốc $v(t)=x'(t)=-4 \pi\sin(\pi t)$ (cm/s), gia tốc $a(t)=v'(t)=-4 \pi^{2}\cos(\pi t)$ (cm/s$^2$). $v(\dfrac{1}{6})=- 2 \pi$.}]{Một vật dao động điều hòa có phương trình li độ $x(t)=4\cos(\pi t)$ (cm), với $t$ tính bằng giây. Vận tốc của vật tại thời điểm $t=\dfrac{1}{6}$ giây bằng}
  {$- 2 \pi\,\mathrm{cm/s}$.}
  {$- 3 \pi\,\mathrm{cm/s}$.}{$- \pi\,\mathrm{cm/s}$.}{$1 - 2 \pi\,\mathrm{cm/s}$.}

\cauTN[id=C1-B5-D01-TH-011,muc=TH,loigiai={Vận tốc $v(t)=x'(t)=-2 \pi\sin(\pi t)$ (cm/s), gia tốc $a(t)=v'(t)=-2 \pi^{2}\cos(\pi t)$ (cm/s$^2$). $a(\dfrac{1}{3})=- \pi^{2}$.}]{Một vật dao động điều hòa có phương trình li độ $x(t)=2\cos(\pi t)$ (cm), với $t$ tính bằng giây. Gia tốc của vật tại thời điểm $t=\dfrac{1}{3}$ giây bằng}
  {$- \pi^{2}\,\mathrm{cm/s^2}$.}
  {$- \dfrac{3 \pi^{2}}{2}\,\mathrm{cm/s^2}$.}{$1 - \pi^{2}\,\mathrm{cm/s^2}$.}{$- 2 \pi^{2}\,\mathrm{cm/s^2}$.}

\cauTN[id=C1-B5-D01-TH-012,muc=TH,loigiai={Ta có $v(t)=h'(t)=10-10t$. $h'(t)=0\Leftrightarrow t=1$.}]{Một quả bóng được ném thẳng đứng lên trên từ độ cao $15$ m với vận tốc ban đầu $10$ m/s. Độ cao của quả bóng (tính bằng mét) sau $t$ giây là $h(t)=15+10t-5t^2$ (bỏ qua sức cản không khí, lấy $g=10$ m/s$^2$). Thời điểm quả bóng đạt độ cao lớn nhất là}
  {$1\,\mathrm{s}$.}
  {$5\,\mathrm{s}$.}{$2\,\mathrm{s}$.}{$3\,\mathrm{s}$.}

\cauTN[id=C1-B5-D01-TH-013,muc=TH,loigiai={Vận tốc $v(t)=s'(t)=16-6t$. Gia tốc $a(t)=v'(t)=-6$.}]{Một ô tô đang chạy với vận tốc $16$ m/s thì người lái đạp phanh. Từ thời điểm đó, ô tô chuyển động chậm dần đều và quãng đường đi được sau $t$ giây là $s(t)=16t-3t^2$ (mét). Gia tốc của ô tô bằng}
  {$-6\,\mathrm{m/s^2}$.}
  {$-8\,\mathrm{m/s^2}$.}{$6\,\mathrm{m/s^2}$.}{$-2\,\mathrm{m/s^2}$.}

\cauTN[id=C1-B5-D01-TH-014,muc=TH,loigiai={Ta có $v(t)=h'(t)=12-10t$. Do đó $v(2)=-8$.}]{Một quả bóng được ném thẳng đứng lên trên từ độ cao $20$ m với vận tốc ban đầu $12$ m/s. Độ cao của quả bóng (tính bằng mét) sau $t$ giây là $h(t)=20+12t-5t^2$ (bỏ qua sức cản không khí, lấy $g=10$ m/s$^2$). Vận tốc của quả bóng tại thời điểm $t=2$ giây bằng}
  {$-8\,\mathrm{m/s}$.}
  {$8\,\mathrm{m/s}$.}{$-6\,\mathrm{m/s}$.}{$-9\,\mathrm{m/s}$.}

\cauTN[id=C1-B5-D01-TH-015,muc=TH,loigiai={Ta có $h(\dfrac{1}{2})=\dfrac{85}{4}$.}]{Một quả bóng được ném thẳng đứng lên trên từ độ cao $15$ m với vận tốc ban đầu $15$ m/s. Độ cao của quả bóng (tính bằng mét) sau $t$ giây là $h(t)=15+15t-5t^2$ (bỏ qua sức cản không khí, lấy $g=10$ m/s$^2$). Độ cao của quả bóng tại thời điểm $t=\dfrac{1}{2}$ giây bằng}
  {$\dfrac{85}{4}\,\mathrm{m}$.}
  {$\dfrac{85}{2}\,\mathrm{m}$.}{$\dfrac{85}{8}\,\mathrm{m}$.}{$\dfrac{81}{4}\,\mathrm{m}$.}

\cauTN[id=C1-B5-D01-TH-016,muc=TH,loigiai={Ta có $v(t)=h'(t)=12-10t$. $h'(t)=0\Leftrightarrow t=\dfrac{6}{5}$.}]{Một quả bóng được ném thẳng đứng lên trên từ độ cao $20$ m với vận tốc ban đầu $12$ m/s. Độ cao của quả bóng (tính bằng mét) sau $t$ giây là $h(t)=20+12t-5t^2$ (bỏ qua sức cản không khí, lấy $g=10$ m/s$^2$). Thời điểm quả bóng đạt độ cao lớn nhất là}
  {$\dfrac{6}{5}\,\mathrm{s}$.}
  {$\dfrac{3}{5}\,\mathrm{s}$.}{$\dfrac{5}{6}\,\mathrm{s}$.}{$\dfrac{11}{5}\,\mathrm{s}$.}

\cauTN[id=C1-B5-D01-TH-017,muc=TH,loigiai={Vận tốc $v(t)=s'(t)=-t^{2}+10t$, gia tốc $a(t)=v'(t)=-2t+10$. $v'(t)=0\Leftrightarrow t=5$ và $v'$ đổi dấu từ dương sang âm nên vận tốc lớn nhất tại $t=5$.}]{Một vật chuyển động theo quy luật $s(t)=-\dfrac{1}{3}t^3+5t^2$, trong đó $t$ (giây) là khoảng thời gian tính từ lúc vật bắt đầu chuyển động và $s$ (mét) là quãng đường vật đi được trong khoảng thời gian đó, $0\le t\le 8$. Thời điểm vật đạt vận tốc lớn nhất là}
  {$5\,\mathrm{s}$.}
  {$25\,\mathrm{s}$.}{$10\,\mathrm{s}$.}{$1\,\mathrm{s}$.}

\cauTN[id=C1-B5-D01-TH-018,muc=TH,loigiai={Ta có $x(\dfrac{1}{2})=4 \sqrt{2}$.}]{Một vật dao động điều hòa có phương trình li độ $x(t)=8\cos(\frac{\pi}{2} t)$ (cm), với $t$ tính bằng giây. Li độ của vật tại thời điểm $t=\dfrac{1}{2}$ giây bằng}
  {$4 \sqrt{2}\,\mathrm{cm}$.}
  {$2 \sqrt{2}\,\mathrm{cm}$.}{$8 \sqrt{2}\,\mathrm{cm}$.}{$-1 + 4 \sqrt{2}\,\mathrm{cm}$.}

\cauTN[id=C1-B5-D01-TH-019,muc=TH,loigiai={Vận tốc $v(t)=x'(t)=-4 \pi\sin(\frac{\pi}{2} t)$ (cm/s), gia tốc $a(t)=v'(t)=-2 \pi^{2}\cos(\frac{\pi}{2} t)$ (cm/s$^2$). $a(\dfrac{1}{3})=- \sqrt{3} \pi^{2}$.}]{Một vật dao động điều hòa có phương trình li độ $x(t)=8\cos(\frac{\pi}{2} t)$ (cm), với $t$ tính bằng giây. Gia tốc của vật tại thời điểm $t=\dfrac{1}{3}$ giây bằng}
  {$- \sqrt{3} \pi^{2}\,\mathrm{cm/s^2}$.}
  {$- \sqrt{3} \pi^{2} + 1\,\mathrm{cm/s^2}$.}{$- \dfrac{\sqrt{3} \pi^{2}}{2}\,\mathrm{cm/s^2}$.}{$- 2 \sqrt{3} \pi^{2}\,\mathrm{cm/s^2}$.}

\cauTN[id=C1-B5-D01-TH-020,muc=TH,loigiai={Vận tốc $v(t)=x'(t)=-8 \pi\sin(4 \pi t)$ (cm/s), gia tốc $a(t)=v'(t)=-32 \pi^{2}\cos(4 \pi t)$ (cm/s$^2$). $a(\dfrac{1}{4})=32 \pi^{2}$.}]{Một vật dao động điều hòa có phương trình li độ $x(t)=2\cos(4 \pi t)$ (cm), với $t$ tính bằng giây. Gia tốc của vật tại thời điểm $t=\dfrac{1}{4}$ giây bằng}
  {$32 \pi^{2}\,\mathrm{cm/s^2}$.}
  {$48 \pi^{2}\,\mathrm{cm/s^2}$.}{$16 \pi^{2}\,\mathrm{cm/s^2}$.}{$-1 + 32 \pi^{2}\,\mathrm{cm/s^2}$.}

\cauTN[id=C1-B5-D01-TH-021,muc=TH,loigiai={Ta có $v(t)=h'(t)=18-10t$. Do đó $v(\dfrac{1}{2})=13$.}]{Một quả bóng được ném thẳng đứng lên trên từ độ cao $0$ m với vận tốc ban đầu $18$ m/s. Độ cao của quả bóng (tính bằng mét) sau $t$ giây là $h(t)=18t-5t^2$ (bỏ qua sức cản không khí, lấy $g=10$ m/s$^2$). Vận tốc của quả bóng tại thời điểm $t=\dfrac{1}{2}$ giây bằng}
  {$13\,\mathrm{m/s}$.}
  {$23\,\mathrm{m/s}$.}{$26\,\mathrm{m/s}$.}{$3\,\mathrm{m/s}$.}

\cauTN[id=C1-B5-D01-TH-022,muc=TH,loigiai={Vận tốc $v(t)=s'(t)=-2t^{2}+36t+5$, gia tốc $a(t)=v'(t)=-4t+36$. $v'(t)=0\Leftrightarrow t=9$; $v$ đồng biến trên $(0;9)$, nghịch biến trên $(9;13)$ nên $v_{\max}=v(9)=167$.}]{Một vật chuyển động theo quy luật $s(t)=-\dfrac{2}{3}t^3+18t^2+5t$, trong đó $t$ (giây) là khoảng thời gian tính từ lúc vật bắt đầu chuyển động và $s$ (mét) là quãng đường vật đi được trong khoảng thời gian đó, $0\le t\le 13$. Vận tốc lớn nhất của vật trong khoảng thời gian trên bằng}
  {$167\,\mathrm{m/s}$.}
  {$9\,\mathrm{m/s}$.}{$135\,\mathrm{m/s}$.}{$67\,\mathrm{m/s}$.}

\cauTN[id=C1-B5-D01-TH-023,muc=TH,loigiai={Vận tốc $v(t)=s'(t)=20-4t$. Ô tô dừng tại $t=5$, quãng đường $s(5)=50$.}]{Một ô tô đang chạy với vận tốc $20$ m/s thì người lái đạp phanh. Từ thời điểm đó, ô tô chuyển động chậm dần đều và quãng đường đi được sau $t$ giây là $s(t)=20t-2t^2$ (mét). Quãng đường ô tô đi được từ lúc đạp phanh đến khi dừng hẳn bằng}
  {$50\,\mathrm{m}$.}
  {$100\,\mathrm{m}$.}{$45\,\mathrm{m}$.}{$55\,\mathrm{m}$.}

\cauTN[id=C1-B5-D01-TH-024,muc=TH,loigiai={Vận tốc $v(t)=s'(t)=24-4t$. Gia tốc $a(t)=v'(t)=-4$.}]{Một ô tô đang chạy với vận tốc $24$ m/s thì người lái đạp phanh. Từ thời điểm đó, ô tô chuyển động chậm dần đều và quãng đường đi được sau $t$ giây là $s(t)=24t-2t^2$ (mét). Gia tốc của ô tô bằng}
  {$-4\,\mathrm{m/s^2}$.}
  {$-8\,\mathrm{m/s^2}$.}{$0\,\mathrm{m/s^2}$.}{$-2\,\mathrm{m/s^2}$.}

\cauTN[id=C1-B5-D01-TH-025,muc=TH,loigiai={Vận tốc $v(t)=s'(t)=10-2t$. Ô tô dừng tại $t=5$, quãng đường $s(5)=25$.}]{Một ô tô đang chạy với vận tốc $10$ m/s thì người lái đạp phanh. Từ thời điểm đó, ô tô chuyển động chậm dần đều và quãng đường đi được sau $t$ giây là $s(t)=10t-t^2$ (mét). Quãng đường ô tô đi được từ lúc đạp phanh đến khi dừng hẳn bằng}
  {$25\,\mathrm{m}$.}
  {$50\,\mathrm{m}$.}{$21\,\mathrm{m}$.}{$29\,\mathrm{m}$.}

\cauTN[id=C1-B5-D01-TH-026,muc=TH,loigiai={Vận tốc $v(t)=x'(t)=-20 \pi\sin(4 \pi t)$ (cm/s), gia tốc $a(t)=v'(t)=-80 \pi^{2}\cos(4 \pi t)$ (cm/s$^2$). $a(\dfrac{1}{2})=- 80 \pi^{2}$.}]{Một vật dao động điều hòa có phương trình li độ $x(t)=5\cos(4 \pi t)$ (cm), với $t$ tính bằng giây. Gia tốc của vật tại thời điểm $t=\dfrac{1}{2}$ giây bằng}
  {$- 80 \pi^{2}\,\mathrm{cm/s^2}$.}
  {$- 40 \pi^{2}\,\mathrm{cm/s^2}$.}{$1 - 80 \pi^{2}\,\mathrm{cm/s^2}$.}{$- 160 \pi^{2}\,\mathrm{cm/s^2}$.}

\cauTN[id=C1-B5-D01-TH-027,muc=TH,loigiai={Ta có $x(\dfrac{1}{3})=-\dfrac{3}{2}$.}]{Một vật dao động điều hòa có phương trình li độ $x(t)=3\cos(4 \pi t)$ (cm), với $t$ tính bằng giây. Li độ của vật tại thời điểm $t=\dfrac{1}{3}$ giây bằng}
  {$-\dfrac{3}{2}\,\mathrm{cm}$.}
  {$-\dfrac{2}{3}\,\mathrm{cm}$.}{$-3\,\mathrm{cm}$.}{$\dfrac{3}{2}\,\mathrm{cm}$.}

\cauTN[id=C1-B5-D01-TH-028,muc=TH,loigiai={Vận tốc $v(t)=s'(t)=24-2t$. Gia tốc $a(t)=v'(t)=-2$.}]{Một ô tô đang chạy với vận tốc $24$ m/s thì người lái đạp phanh. Từ thời điểm đó, ô tô chuyển động chậm dần đều và quãng đường đi được sau $t$ giây là $s(t)=24t-t^2$ (mét). Gia tốc của ô tô bằng}
  {$-2\,\mathrm{m/s^2}$.}
  {$-4\,\mathrm{m/s^2}$.}{$2\,\mathrm{m/s^2}$.}{$-6\,\mathrm{m/s^2}$.}

\cauTN[id=C1-B5-D01-TH-029,muc=TH,loigiai={Ta có $h(2)=41$.}]{Một quả bóng được ném thẳng đứng lên trên từ độ cao $1$ m với vận tốc ban đầu $30$ m/s. Độ cao của quả bóng (tính bằng mét) sau $t$ giây là $h(t)=1+30t-5t^2$ (bỏ qua sức cản không khí, lấy $g=10$ m/s$^2$). Độ cao của quả bóng tại thời điểm $t=2$ giây bằng}
  {$41\,\mathrm{m}$.}
  {$43\,\mathrm{m}$.}{$40\,\mathrm{m}$.}{$82\,\mathrm{m}$.}

\cauTN[id=C1-B5-D01-TH-030,muc=TH,loigiai={Vận tốc $v(t)=s'(t)=25-6t$. $v(1)=19$.}]{Một ô tô đang chạy với vận tốc $25$ m/s thì người lái đạp phanh. Từ thời điểm đó, ô tô chuyển động chậm dần đều và quãng đường đi được sau $t$ giây là $s(t)=25t-3t^2$ (mét). Vận tốc của ô tô sau $1$ giây kể từ lúc đạp phanh bằng}
  {$19\,\mathrm{m/s}$.}
  {$15\,\mathrm{m/s}$.}{$21\,\mathrm{m/s}$.}{$17\,\mathrm{m/s}$.}

\cauTN[id=C1-B5-D01-TH-031,muc=TH,loigiai={Quãng đường $s(1)=17$.}]{Một vật chuyển động theo quy luật $s(t)=-\dfrac{3}{3}t^3+15t^2+3t$, trong đó $t$ (giây) là khoảng thời gian tính từ lúc vật bắt đầu chuyển động và $s$ (mét) là quãng đường vật đi được trong khoảng thời gian đó, $0\le t\le 8$. Quãng đường vật đi được trong $1$ giây đầu tiên bằng}
  {$17\,\mathrm{m}$.}
  {$27\,\mathrm{m}$.}{$22\,\mathrm{m}$.}{$12\,\mathrm{m}$.}

\cauTN[id=C1-B5-D01-TH-032,muc=TH,loigiai={Vận tốc $v(t)=s'(t)=20-2t$. Ô tô dừng tại $t=10$, quãng đường $s(10)=100$.}]{Một ô tô đang chạy với vận tốc $20$ m/s thì người lái đạp phanh. Từ thời điểm đó, ô tô chuyển động chậm dần đều và quãng đường đi được sau $t$ giây là $s(t)=20t-t^2$ (mét). Quãng đường ô tô đi được từ lúc đạp phanh đến khi dừng hẳn bằng}
  {$100\,\mathrm{m}$.}
  {$200\,\mathrm{m}$.}{$120\,\mathrm{m}$.}{$110\,\mathrm{m}$.}

\cauTN[id=C1-B5-D01-TH-033,muc=TH,loigiai={Ta có $v(t)=h'(t)=12-10t$. $h'(t)=0\Leftrightarrow t=\dfrac{6}{5}$, $h$ đạt giá trị lớn nhất bằng $h(\dfrac{6}{5})=\dfrac{46}{5}$.}]{Một quả bóng được ném thẳng đứng lên trên từ độ cao $2$ m với vận tốc ban đầu $12$ m/s. Độ cao của quả bóng (tính bằng mét) sau $t$ giây là $h(t)=2+12t-5t^2$ (bỏ qua sức cản không khí, lấy $g=10$ m/s$^2$). Độ cao lớn nhất mà quả bóng đạt được bằng}
  {$\dfrac{46}{5}\,\mathrm{m}$.}
  {$\dfrac{5}{46}\,\mathrm{m}$.}{$\dfrac{92}{5}\,\mathrm{m}$.}{$\dfrac{23}{5}\,\mathrm{m}$.}

\cauTN[id=C1-B5-D01-TH-034,muc=TH,loigiai={Vận tốc $v(t)=x'(t)=-\dfrac{5 \pi}{2}\sin(\frac{\pi}{2} t)$ (cm/s), gia tốc $a(t)=v'(t)=-\dfrac{5 \pi^{2}}{4}\cos(\frac{\pi}{2} t)$ (cm/s$^2$). $a(\dfrac{1}{4})=- \dfrac{5 \pi^{2} \sqrt{\sqrt{2} + 2}}{8}$.}]{Một vật dao động điều hòa có phương trình li độ $x(t)=5\cos(\frac{\pi}{2} t)$ (cm), với $t$ tính bằng giây. Gia tốc của vật tại thời điểm $t=\dfrac{1}{4}$ giây bằng}
  {$- \dfrac{5 \pi^{2} \sqrt{\sqrt{2} + 2}}{8}\,\mathrm{cm/s^2}$.}
  {$- \dfrac{5 \pi^{2} \sqrt{\sqrt{2} + 2}}{16}\,\mathrm{cm/s^2}$.}{$- \dfrac{5 \pi^{2} \sqrt{\sqrt{2} + 2}}{8} - 1\,\mathrm{cm/s^2}$.}{$\dfrac{5 \pi^{2} \sqrt{\sqrt{2} + 2}}{8}\,\mathrm{cm/s^2}$.}

\cauTN[id=C1-B5-D01-TH-035,muc=TH,loigiai={Quãng đường $s(1)=\dfrac{58}{3}$.}]{Một vật chuyển động theo quy luật $s(t)=-\dfrac{2}{3}t^3+20t^2$, trong đó $t$ (giây) là khoảng thời gian tính từ lúc vật bắt đầu chuyển động và $s$ (mét) là quãng đường vật đi được trong khoảng thời gian đó, $0\le t\le 12$. Quãng đường vật đi được trong $1$ giây đầu tiên bằng}
  {$\dfrac{58}{3}\,\mathrm{m}$.}
  {$\dfrac{29}{3}\,\mathrm{m}$.}{$\dfrac{61}{3}\,\mathrm{m}$.}{$\dfrac{55}{3}\,\mathrm{m}$.}

\cauTN[id=C1-B5-D01-TH-036,muc=TH,loigiai={Vận tốc $v(t)=s'(t)=10-2t$. Gia tốc $a(t)=v'(t)=-2$.}]{Một ô tô đang chạy với vận tốc $10$ m/s thì người lái đạp phanh. Từ thời điểm đó, ô tô chuyển động chậm dần đều và quãng đường đi được sau $t$ giây là $s(t)=10t-t^2$ (mét). Gia tốc của ô tô bằng}
  {$-2\,\mathrm{m/s^2}$.}
  {$-6\,\mathrm{m/s^2}$.}{$-1\,\mathrm{m/s^2}$.}{$2\,\mathrm{m/s^2}$.}

\cauTN[id=C1-B5-D01-TH-037,muc=TH,loigiai={Vận tốc $v(t)=s'(t)=25-6t$. Ô tô dừng khi $v(t)=0\Leftrightarrow t=\dfrac{25}{6}$.}]{Một ô tô đang chạy với vận tốc $25$ m/s thì người lái đạp phanh. Từ thời điểm đó, ô tô chuyển động chậm dần đều và quãng đường đi được sau $t$ giây là $s(t)=25t-3t^2$ (mét). Thời gian từ lúc đạp phanh đến khi ô tô dừng hẳn bằng}
  {$\dfrac{25}{6}\,\mathrm{s}$.}
  {$\dfrac{25}{12}\,\mathrm{s}$.}{$25\,\mathrm{s}$.}{$\dfrac{19}{6}\,\mathrm{s}$.}

\cauTN[id=C1-B5-D01-TH-038,muc=TH,loigiai={Vận tốc $v(t)=s'(t)=16-6t$. Ô tô dừng khi $v(t)=0\Leftrightarrow t=\dfrac{8}{3}$.}]{Một ô tô đang chạy với vận tốc $16$ m/s thì người lái đạp phanh. Từ thời điểm đó, ô tô chuyển động chậm dần đều và quãng đường đi được sau $t$ giây là $s(t)=16t-3t^2$ (mét). Thời gian từ lúc đạp phanh đến khi ô tô dừng hẳn bằng}
  {$\dfrac{8}{3}\,\mathrm{s}$.}
  {$\dfrac{5}{3}\,\mathrm{s}$.}{$\dfrac{3}{8}\,\mathrm{s}$.}{$\dfrac{16}{3}\,\mathrm{s}$.}

\cauTN[id=C1-B5-D01-TH-039,muc=TH,loigiai={Vận tốc $v(t)=s'(t)=-t^{2}+6t+5$, gia tốc $a(t)=v'(t)=-2t+6$. Do đó $a(4)=-2$.}]{Một vật chuyển động theo quy luật $s(t)=-\dfrac{1}{3}t^3+3t^2+5t$, trong đó $t$ (giây) là khoảng thời gian tính từ lúc vật bắt đầu chuyển động và $s$ (mét) là quãng đường vật đi được trong khoảng thời gian đó, $0\le t\le 6$. Gia tốc của vật tại thời điểm $t=4$ giây bằng}
  {$-2\,\mathrm{m/s^2}$.}
  {$-4\,\mathrm{m/s^2}$.}{$-6\,\mathrm{m/s^2}$.}{$2\,\mathrm{m/s^2}$.}

\cauTN[id=C1-B5-D01-TH-040,muc=TH,loigiai={Vận tốc $v(t)=x'(t)=-3 \pi\sin(\pi t)$ (cm/s), gia tốc $a(t)=v'(t)=-3 \pi^{2}\cos(\pi t)$ (cm/s$^2$). $v(\dfrac{1}{3})=- \dfrac{3 \sqrt{3} \pi}{2}$.}]{Một vật dao động điều hòa có phương trình li độ $x(t)=3\cos(\pi t)$ (cm), với $t$ tính bằng giây. Vận tốc của vật tại thời điểm $t=\dfrac{1}{3}$ giây bằng}
  {$- \dfrac{3 \sqrt{3} \pi}{2}\,\mathrm{cm/s}$.}
  {$- \dfrac{3 \sqrt{3} \pi}{4}\,\mathrm{cm/s}$.}{$- \dfrac{9 \sqrt{3} \pi}{4}\,\mathrm{cm/s}$.}{$- \dfrac{3 \sqrt{3} \pi}{2} + 1\,\mathrm{cm/s}$.}

\cauTN[id=C1-B5-D01-TH-041,muc=TH,loigiai={Ta có $v(t)=h'(t)=40-10t$. Do đó $v(\dfrac{1}{2})=35$.}]{Một quả bóng được ném thẳng đứng lên trên từ độ cao $5$ m với vận tốc ban đầu $40$ m/s. Độ cao của quả bóng (tính bằng mét) sau $t$ giây là $h(t)=5+40t-5t^2$ (bỏ qua sức cản không khí, lấy $g=10$ m/s$^2$). Vận tốc của quả bóng tại thời điểm $t=\dfrac{1}{2}$ giây bằng}
  {$35\,\mathrm{m/s}$.}
  {$39\,\mathrm{m/s}$.}{$70\,\mathrm{m/s}$.}{$37\,\mathrm{m/s}$.}

\cauTN[id=C1-B5-D01-TH-042,muc=TH,loigiai={Ta có $x(\dfrac{1}{3})=4$.}]{Một vật dao động điều hòa có phương trình li độ $x(t)=8\cos(\pi t)$ (cm), với $t$ tính bằng giây. Li độ của vật tại thời điểm $t=\dfrac{1}{3}$ giây bằng}
  {$4\,\mathrm{cm}$.}
  {$8\,\mathrm{cm}$.}{$5\,\mathrm{cm}$.}{$2\,\mathrm{cm}$.}

\cauTN[id=C1-B5-D01-TH-043,muc=TH,loigiai={Vận tốc $v(t)=s'(t)=16-5t$. Ô tô dừng tại $t=\dfrac{16}{5}$, quãng đường $s(\dfrac{16}{5})=\dfrac{128}{5}$.}]{Một ô tô đang chạy với vận tốc $16$ m/s thì người lái đạp phanh. Từ thời điểm đó, ô tô chuyển động chậm dần đều và quãng đường đi được sau $t$ giây là $s(t)=16t-\dfrac{5}{2}t^2$ (mét). Quãng đường ô tô đi được từ lúc đạp phanh đến khi dừng hẳn bằng}
  {$\dfrac{128}{5}\,\mathrm{m}$.}
  {$\dfrac{256}{5}\,\mathrm{m}$.}{$128\,\mathrm{m}$.}{$\dfrac{64}{5}\,\mathrm{m}$.}

\cauTN[id=C1-B5-D01-TH-044,muc=TH,loigiai={Ta có $h(2)=21$.}]{Một quả bóng được ném thẳng đứng lên trên từ độ cao $1$ m với vận tốc ban đầu $20$ m/s. Độ cao của quả bóng (tính bằng mét) sau $t$ giây là $h(t)=1+20t-5t^2$ (bỏ qua sức cản không khí, lấy $g=10$ m/s$^2$). Độ cao của quả bóng tại thời điểm $t=2$ giây bằng}
  {$21\,\mathrm{m}$.}
  {$42\,\mathrm{m}$.}{$19\,\mathrm{m}$.}{$23\,\mathrm{m}$.}

\cauTN[id=C1-B5-D01-TH-045,muc=TH,loigiai={Vận tốc $v(t)=x'(t)=-8 \pi\sin(4 \pi t)$ (cm/s), gia tốc $a(t)=v'(t)=-32 \pi^{2}\cos(4 \pi t)$ (cm/s$^2$). $v(\dfrac{1}{6})=- 4 \sqrt{3} \pi$.}]{Một vật dao động điều hòa có phương trình li độ $x(t)=2\cos(4 \pi t)$ (cm), với $t$ tính bằng giây. Vận tốc của vật tại thời điểm $t=\dfrac{1}{6}$ giây bằng}
  {$- 4 \sqrt{3} \pi\,\mathrm{cm/s}$.}
  {$- 6 \sqrt{3} \pi\,\mathrm{cm/s}$.}{$- 2 \sqrt{3} \pi\,\mathrm{cm/s}$.}{$- 8 \sqrt{3} \pi\,\mathrm{cm/s}$.}

\cauTN[id=C1-B5-D01-TH-046,muc=TH,loigiai={Vận tốc $v(t)=s'(t)=10-3t$. $v(1)=7$.}]{Một ô tô đang chạy với vận tốc $10$ m/s thì người lái đạp phanh. Từ thời điểm đó, ô tô chuyển động chậm dần đều và quãng đường đi được sau $t$ giây là $s(t)=10t-\dfrac{3}{2}t^2$ (mét). Vận tốc của ô tô sau $1$ giây kể từ lúc đạp phanh bằng}
  {$7\,\mathrm{m/s}$.}
  {$3\,\mathrm{m/s}$.}{$11\,\mathrm{m/s}$.}{$14\,\mathrm{m/s}$.}

\cauTN[id=C1-B5-D01-TH-047,muc=TH,loigiai={Ta có $v(t)=h'(t)=40-10t$. $h'(t)=0\Leftrightarrow t=4$, $h$ đạt giá trị lớn nhất bằng $h(4)=82$.}]{Một quả bóng được ném thẳng đứng lên trên từ độ cao $2$ m với vận tốc ban đầu $40$ m/s. Độ cao của quả bóng (tính bằng mét) sau $t$ giây là $h(t)=2+40t-5t^2$ (bỏ qua sức cản không khí, lấy $g=10$ m/s$^2$). Độ cao lớn nhất mà quả bóng đạt được bằng}
  {$82\,\mathrm{m}$.}
  {$164\,\mathrm{m}$.}{$41\,\mathrm{m}$.}{$84\,\mathrm{m}$.}

\cauTN[id=C1-B5-D01-TH-048,muc=TH,loigiai={Vận tốc $v(t)=x'(t)=-8 \pi\sin(4 \pi t)$ (cm/s), gia tốc $a(t)=v'(t)=-32 \pi^{2}\cos(4 \pi t)$ (cm/s$^2$). $a(\dfrac{1}{6})=16 \pi^{2}$.}]{Một vật dao động điều hòa có phương trình li độ $x(t)=2\cos(4 \pi t)$ (cm), với $t$ tính bằng giây. Gia tốc của vật tại thời điểm $t=\dfrac{1}{6}$ giây bằng}
  {$16 \pi^{2}\,\mathrm{cm/s^2}$.}
  {$8 \pi^{2}\,\mathrm{cm/s^2}$.}{$24 \pi^{2}\,\mathrm{cm/s^2}$.}{$1 + 16 \pi^{2}\,\mathrm{cm/s^2}$.}

\cauDS[id=C1-B5-D01-TH-049,muc=TH,loigiai={Ta có $h(1)=20$. Ta có $v(t)=h'(t)=25-10t$. $h'(t)=0\Leftrightarrow t=\dfrac{5}{2}$. Ta có $v(t)=h'(t)=25-10t$. Do đó $v(1)=15$. Ta có $v(t)=h'(t)=25-10t$. $h'(t)=0\Leftrightarrow t=\dfrac{5}{2}$, $h$ đạt giá trị lớn nhất bằng $h(\dfrac{5}{2})=\dfrac{125}{4}$. a) Đúng. b) Sai: giá trị đúng là $\dfrac{5}{2}\,\mathrm{s}$. c) Sai: giá trị đúng là $\dfrac{125}{4}\,\mathrm{m}$. d) Sai: giá trị đúng là $20\,\mathrm{m}$.}]{Một quả bóng được ném thẳng đứng lên trên từ độ cao $0$ m với vận tốc ban đầu $25$ m/s. Độ cao của quả bóng (tính bằng mét) sau $t$ giây là $h(t)=25t-5t^2$ (bỏ qua sức cản không khí, lấy $g=10$ m/s$^2$). Xét tính đúng sai của các mệnh đề sau:}{%
  \y{Vận tốc của quả bóng tại thời điểm $t=1$ giây là $15\,\mathrm{m/s}$.}{d}
  \y{Quả bóng đạt độ cao lớn nhất tại thời điểm $t=\dfrac{5}{4}\,\mathrm{s}$.}{s}
  \y{Độ cao lớn nhất của quả bóng là $\dfrac{129}{4}\,\mathrm{m}$.}{s}
  \y{Sau $1$ giây quả bóng ở độ cao $16\,\mathrm{m}$.}{s}}

\cauDS[id=C1-B5-D01-TH-050,muc=TH,loigiai={Vận tốc $v(t)=x'(t)=-6 \pi\sin(\pi t)$ (cm/s), gia tốc $a(t)=v'(t)=-6 \pi^{2}\cos(\pi t)$ (cm/s$^2$). $a(\dfrac{1}{3})=- 3 \pi^{2}$. Ta có $x(\dfrac{1}{3})=3$. Vận tốc $v(t)=x'(t)=-6 \pi\sin(\pi t)$ (cm/s), gia tốc $a(t)=v'(t)=-6 \pi^{2}\cos(\pi t)$ (cm/s$^2$). $v(\dfrac{1}{3})=- 3 \sqrt{3} \pi$. Vận tốc $v(t)=x'(t)=-6 \pi\sin(\pi t)$ (cm/s), gia tốc $a(t)=v'(t)=-6 \pi^{2}\cos(\pi t)$ (cm/s$^2$). Vì $|\sin(\pi t)|\le1$ nên $|v(t)|\le 6 \pi$, dấu bằng xảy ra khi $\sin(\pi t)=\pm1$. a) Sai: giá trị đúng là $6 \pi\,\mathrm{cm/s}$. b) Sai: giá trị đúng là $- 3 \sqrt{3} \pi\,\mathrm{cm/s}$. c) Đúng. d) Đúng.}]{Một vật dao động điều hòa có phương trình li độ $x(t)=6\cos(\pi t)$ (cm), với $t$ tính bằng giây. Xét tính đúng sai của các mệnh đề sau:}{%
  \y{Tốc độ cực đại của vật là $12 \pi\,\mathrm{cm/s}$.}{s}
  \y{Vận tốc của vật tại thời điểm $t=\dfrac{1}{3}$ giây là $- 3 \sqrt{3} \pi - 1\,\mathrm{cm/s}$.}{s}
  \y{Li độ của vật tại thời điểm $t=\dfrac{1}{3}$ giây là $3\,\mathrm{cm}$.}{d}
  \y{Gia tốc của vật tại thời điểm $t=\dfrac{1}{3}$ giây là $- 3 \pi^{2}\,\mathrm{cm/s^2}$.}{d}}

\cauDS[id=C1-B5-D01-TH-051,muc=TH,loigiai={Vận tốc $v(t)=x'(t)=-40 \pi\sin(4 \pi t)$ (cm/s), gia tốc $a(t)=v'(t)=-160 \pi^{2}\cos(4 \pi t)$ (cm/s$^2$). $a(\dfrac{1}{6})=80 \pi^{2}$. Vận tốc $v(t)=x'(t)=-40 \pi\sin(4 \pi t)$ (cm/s), gia tốc $a(t)=v'(t)=-160 \pi^{2}\cos(4 \pi t)$ (cm/s$^2$). $v(\dfrac{1}{6})=- 20 \sqrt{3} \pi$. Ta có $x(\dfrac{1}{6})=-5$. Vận tốc $v(t)=x'(t)=-40 \pi\sin(4 \pi t)$ (cm/s), gia tốc $a(t)=v'(t)=-160 \pi^{2}\cos(4 \pi t)$ (cm/s$^2$). Vì $|\sin(4 \pi t)|\le1$ nên $|v(t)|\le 40 \pi$, dấu bằng xảy ra khi $\sin(4 \pi t)=\pm1$. a) Đúng. b) Đúng. c) Sai: giá trị đúng là $-5\,\mathrm{cm}$. d) Sai: giá trị đúng là $- 20 \sqrt{3} \pi\,\mathrm{cm/s}$.}]{Một vật dao động điều hòa có phương trình li độ $x(t)=10\cos(4 \pi t)$ (cm), với $t$ tính bằng giây. Xét tính đúng sai của các mệnh đề sau:}{%
  \y{Tốc độ cực đại của vật là $40 \pi\,\mathrm{cm/s}$.}{d}
  \y{Gia tốc của vật tại thời điểm $t=\dfrac{1}{6}$ giây là $80 \pi^{2}\,\mathrm{cm/s^2}$.}{d}
  \y{Li độ của vật tại thời điểm $t=\dfrac{1}{6}$ giây là $5\,\mathrm{cm}$.}{s}
  \y{Vận tốc của vật tại thời điểm $t=\dfrac{1}{6}$ giây là $- 20 \sqrt{3} \pi + 1\,\mathrm{cm/s}$.}{s}}

\cauDS[id=C1-B5-D01-TH-052,muc=TH,loigiai={Ta có $x(\dfrac{1}{3})=\sqrt{3}$. Vận tốc $v(t)=x'(t)=-\pi\sin(\frac{\pi}{2} t)$ (cm/s), gia tốc $a(t)=v'(t)=-\dfrac{\pi^{2}}{2}\cos(\frac{\pi}{2} t)$ (cm/s$^2$). Vì $|\sin(\frac{\pi}{2} t)|\le1$ nên $|v(t)|\le \pi$, dấu bằng xảy ra khi $\sin(\frac{\pi}{2} t)=\pm1$. Vận tốc $v(t)=x'(t)=-\pi\sin(\frac{\pi}{2} t)$ (cm/s), gia tốc $a(t)=v'(t)=-\dfrac{\pi^{2}}{2}\cos(\frac{\pi}{2} t)$ (cm/s$^2$). $a(\dfrac{1}{3})=- \dfrac{\sqrt{3} \pi^{2}}{4}$. Vận tốc $v(t)=x'(t)=-\pi\sin(\frac{\pi}{2} t)$ (cm/s), gia tốc $a(t)=v'(t)=-\dfrac{\pi^{2}}{2}\cos(\frac{\pi}{2} t)$ (cm/s$^2$). $v(\dfrac{1}{3})=- \dfrac{\pi}{2}$. a) Đúng. b) Sai: giá trị đúng là $\pi\,\mathrm{cm/s}$. c) Đúng. d) Đúng.}]{Một vật dao động điều hòa có phương trình li độ $x(t)=2\cos(\frac{\pi}{2} t)$ (cm), với $t$ tính bằng giây. Xét tính đúng sai của các mệnh đề sau:}{%
  \y{Gia tốc của vật tại thời điểm $t=\dfrac{1}{3}$ giây là $- \dfrac{\sqrt{3} \pi^{2}}{4}\,\mathrm{cm/s^2}$.}{d}
  \y{Tốc độ cực đại của vật là $1 + \pi\,\mathrm{cm/s}$.}{s}
  \y{Vận tốc của vật tại thời điểm $t=\dfrac{1}{3}$ giây là $- \dfrac{\pi}{2}\,\mathrm{cm/s}$.}{d}
  \y{Li độ của vật tại thời điểm $t=\dfrac{1}{3}$ giây là $\sqrt{3}\,\mathrm{cm}$.}{d}}

\cauDS[id=C1-B5-D01-TH-053,muc=TH,loigiai={Ta có $v(t)=h'(t)=20-10t$. $h'(t)=0\Leftrightarrow t=2$, $h$ đạt giá trị lớn nhất bằng $h(2)=21$. Ta có $v(t)=h'(t)=20-10t$. Do đó $v(2)=0$. Ta có $h(2)=21$. Ta có $v(t)=h'(t)=20-10t$. $h'(t)=0\Leftrightarrow t=2$. a) Đúng. b) Đúng. c) Sai: giá trị đúng là $2\,\mathrm{s}$. d) Sai: giá trị đúng là $21\,\mathrm{m}$.}]{Một quả bóng được ném thẳng đứng lên trên từ độ cao $1$ m với vận tốc ban đầu $20$ m/s. Độ cao của quả bóng (tính bằng mét) sau $t$ giây là $h(t)=1+20t-5t^2$ (bỏ qua sức cản không khí, lấy $g=10$ m/s$^2$). Xét tính đúng sai của các mệnh đề sau:}{%
  \y{Vận tốc của quả bóng tại thời điểm $t=2$ giây là $0\,\mathrm{m/s}$.}{d}
  \y{Sau $2$ giây quả bóng ở độ cao $21\,\mathrm{m}$.}{d}
  \y{Quả bóng đạt độ cao lớn nhất tại thời điểm $t=6\,\mathrm{s}$.}{s}
  \y{Độ cao lớn nhất của quả bóng là $11\,\mathrm{m}$.}{s}}

\cauDS[id=C1-B5-D01-TH-054,muc=TH,loigiai={Vận tốc $v(t)=x'(t)=-20 \pi\sin(4 \pi t)$ (cm/s), gia tốc $a(t)=v'(t)=-80 \pi^{2}\cos(4 \pi t)$ (cm/s$^2$). $v(\dfrac{1}{6})=- 10 \sqrt{3} \pi$. Vận tốc $v(t)=x'(t)=-20 \pi\sin(4 \pi t)$ (cm/s), gia tốc $a(t)=v'(t)=-80 \pi^{2}\cos(4 \pi t)$ (cm/s$^2$). $a(\dfrac{1}{6})=40 \pi^{2}$. Vận tốc $v(t)=x'(t)=-20 \pi\sin(4 \pi t)$ (cm/s), gia tốc $a(t)=v'(t)=-80 \pi^{2}\cos(4 \pi t)$ (cm/s$^2$). Vì $|\sin(4 \pi t)|\le1$ nên $|v(t)|\le 20 \pi$, dấu bằng xảy ra khi $\sin(4 \pi t)=\pm1$. Ta có $x(\dfrac{1}{6})=-\dfrac{5}{2}$. a) Sai: giá trị đúng là $40 \pi^{2}\,\mathrm{cm/s^2}$. b) Đúng. c) Sai: giá trị đúng là $- 10 \sqrt{3} \pi\,\mathrm{cm/s}$. d) Đúng.}]{Một vật dao động điều hòa có phương trình li độ $x(t)=5\cos(4 \pi t)$ (cm), với $t$ tính bằng giây. Xét tính đúng sai của các mệnh đề sau:}{%
  \y{Gia tốc của vật tại thời điểm $t=\dfrac{1}{6}$ giây là $1 + 40 \pi^{2}\,\mathrm{cm/s^2}$.}{s}
  \y{Li độ của vật tại thời điểm $t=\dfrac{1}{6}$ giây là $-\dfrac{5}{2}\,\mathrm{cm}$.}{d}
  \y{Vận tốc của vật tại thời điểm $t=\dfrac{1}{6}$ giây là $- 10 \sqrt{3} \pi + 1\,\mathrm{cm/s}$.}{s}
  \y{Tốc độ cực đại của vật là $20 \pi\,\mathrm{cm/s}$.}{d}}

\cauDS[id=C1-B5-D01-TH-055,muc=TH,loigiai={Ta có $h(\dfrac{1}{2})=\dfrac{51}{4}$. Ta có $v(t)=h'(t)=18-10t$. $h'(t)=0\Leftrightarrow t=\dfrac{9}{5}$, $h$ đạt giá trị lớn nhất bằng $h(\dfrac{9}{5})=\dfrac{106}{5}$. Ta có $v(t)=h'(t)=18-10t$. Do đó $v(\dfrac{1}{2})=13$. Ta có $v(t)=h'(t)=18-10t$. $h'(t)=0\Leftrightarrow t=\dfrac{9}{5}$. a) Sai: giá trị đúng là $13\,\mathrm{m/s}$. b) Sai: giá trị đúng là $\dfrac{51}{4}\,\mathrm{m}$. c) Đúng. d) Sai: giá trị đúng là $\dfrac{9}{5}\,\mathrm{s}$.}]{Một quả bóng được ném thẳng đứng lên trên từ độ cao $5$ m với vận tốc ban đầu $18$ m/s. Độ cao của quả bóng (tính bằng mét) sau $t$ giây là $h(t)=5+18t-5t^2$ (bỏ qua sức cản không khí, lấy $g=10$ m/s$^2$). Xét tính đúng sai của các mệnh đề sau:}{%
  \y{Vận tốc của quả bóng tại thời điểm $t=\dfrac{1}{2}$ giây là $17\,\mathrm{m/s}$.}{s}
  \y{Sau $\dfrac{1}{2}$ giây quả bóng ở độ cao $\dfrac{4}{51}\,\mathrm{m}$.}{s}
  \y{Độ cao lớn nhất của quả bóng là $\dfrac{106}{5}\,\mathrm{m}$.}{d}
  \y{Quả bóng đạt độ cao lớn nhất tại thời điểm $t=\dfrac{18}{5}\,\mathrm{s}$.}{s}}

\cauDS[id=C1-B5-D01-TH-056,muc=TH,loigiai={Vận tốc $v(t)=x'(t)=-16 \pi\sin(2 \pi t)$ (cm/s), gia tốc $a(t)=v'(t)=-32 \pi^{2}\cos(2 \pi t)$ (cm/s$^2$). $v(\dfrac{1}{12})=- 8 \pi$. Vận tốc $v(t)=x'(t)=-16 \pi\sin(2 \pi t)$ (cm/s), gia tốc $a(t)=v'(t)=-32 \pi^{2}\cos(2 \pi t)$ (cm/s$^2$). $a(\dfrac{1}{12})=- 16 \sqrt{3} \pi^{2}$. Vận tốc $v(t)=x'(t)=-16 \pi\sin(2 \pi t)$ (cm/s), gia tốc $a(t)=v'(t)=-32 \pi^{2}\cos(2 \pi t)$ (cm/s$^2$). Vì $|\sin(2 \pi t)|\le1$ nên $|v(t)|\le 16 \pi$, dấu bằng xảy ra khi $\sin(2 \pi t)=\pm1$. Ta có $x(\dfrac{1}{12})=4 \sqrt{3}$. a) Đúng. b) Sai: giá trị đúng là $- 16 \sqrt{3} \pi^{2}\,\mathrm{cm/s^2}$. c) Đúng. d) Đúng.}]{Một vật dao động điều hòa có phương trình li độ $x(t)=8\cos(2 \pi t)$ (cm), với $t$ tính bằng giây. Xét tính đúng sai của các mệnh đề sau:}{%
  \y{Vận tốc của vật tại thời điểm $t=\dfrac{1}{12}$ giây là $- 8 \pi\,\mathrm{cm/s}$.}{d}
  \y{Gia tốc của vật tại thời điểm $t=\dfrac{1}{12}$ giây là $- 8 \sqrt{3} \pi^{2}\,\mathrm{cm/s^2}$.}{s}
  \y{Li độ của vật tại thời điểm $t=\dfrac{1}{12}$ giây là $4 \sqrt{3}\,\mathrm{cm}$.}{d}
  \y{Tốc độ cực đại của vật là $16 \pi\,\mathrm{cm/s}$.}{d}}

\cauDS[id=C1-B5-D01-TH-057,muc=TH,loigiai={Ta có $v(t)=h'(t)=12-10t$. Do đó $v(\dfrac{1}{2})=7$. Ta có $v(t)=h'(t)=12-10t$. $h'(t)=0\Leftrightarrow t=\dfrac{6}{5}$. Ta có $h(\dfrac{1}{2})=\dfrac{59}{4}$. Ta có $v(t)=h'(t)=12-10t$. $h'(t)=0\Leftrightarrow t=\dfrac{6}{5}$, $h$ đạt giá trị lớn nhất bằng $h(\dfrac{6}{5})=\dfrac{86}{5}$. a) Đúng. b) Sai: giá trị đúng là $\dfrac{6}{5}\,\mathrm{s}$. c) Sai: giá trị đúng là $\dfrac{86}{5}\,\mathrm{m}$. d) Đúng.}]{Một quả bóng được ném thẳng đứng lên trên từ độ cao $10$ m với vận tốc ban đầu $12$ m/s. Độ cao của quả bóng (tính bằng mét) sau $t$ giây là $h(t)=10+12t-5t^2$ (bỏ qua sức cản không khí, lấy $g=10$ m/s$^2$). Xét tính đúng sai của các mệnh đề sau:}{%
  \y{Sau $\dfrac{1}{2}$ giây quả bóng ở độ cao $\dfrac{59}{4}\,\mathrm{m}$.}{d}
  \y{Quả bóng đạt độ cao lớn nhất tại thời điểm $t=\dfrac{11}{5}\,\mathrm{s}$.}{s}
  \y{Độ cao lớn nhất của quả bóng là $\dfrac{172}{5}\,\mathrm{m}$.}{s}
  \y{Vận tốc của quả bóng tại thời điểm $t=\dfrac{1}{2}$ giây là $7\,\mathrm{m/s}$.}{d}}

\cauDS[id=C1-B5-D01-TH-058,muc=TH,loigiai={Ta có $v(t)=h'(t)=12-10t$. $h'(t)=0\Leftrightarrow t=\dfrac{6}{5}$, $h$ đạt giá trị lớn nhất bằng $h(\dfrac{6}{5})=\dfrac{41}{5}$. Ta có $v(t)=h'(t)=12-10t$. $h'(t)=0\Leftrightarrow t=\dfrac{6}{5}$. Ta có $v(t)=h'(t)=12-10t$. Do đó $v(2)=-8$. Ta có $h(2)=5$. a) Đúng. b) Sai: giá trị đúng là $\dfrac{41}{5}\,\mathrm{m}$. c) Đúng. d) Đúng.}]{Một quả bóng được ném thẳng đứng lên trên từ độ cao $1$ m với vận tốc ban đầu $12$ m/s. Độ cao của quả bóng (tính bằng mét) sau $t$ giây là $h(t)=1+12t-5t^2$ (bỏ qua sức cản không khí, lấy $g=10$ m/s$^2$). Xét tính đúng sai của các mệnh đề sau:}{%
  \y{Vận tốc của quả bóng tại thời điểm $t=2$ giây là $-8\,\mathrm{m/s}$.}{d}
  \y{Độ cao lớn nhất của quả bóng là $\dfrac{5}{41}\,\mathrm{m}$.}{s}
  \y{Quả bóng đạt độ cao lớn nhất tại thời điểm $t=\dfrac{6}{5}\,\mathrm{s}$.}{d}
  \y{Sau $2$ giây quả bóng ở độ cao $5\,\mathrm{m}$.}{d}}

\cauDS[id=C1-B5-D01-TH-059,muc=TH,loigiai={Vận tốc $v(t)=s'(t)=15-6t$. Gia tốc $a(t)=v'(t)=-6$. Vận tốc $v(t)=s'(t)=15-6t$. $v(1)=9$. Vận tốc $v(t)=s'(t)=15-6t$. Ô tô dừng tại $t=\dfrac{5}{2}$, quãng đường $s(\dfrac{5}{2})=\dfrac{75}{4}$. Vận tốc $v(t)=s'(t)=15-6t$. Ô tô dừng khi $v(t)=0\Leftrightarrow t=\dfrac{5}{2}$. a) Sai: giá trị đúng là $\dfrac{75}{4}\,\mathrm{m}$. b) Sai: giá trị đúng là $-6\,\mathrm{m/s^2}$. c) Sai: giá trị đúng là $9\,\mathrm{m/s}$. d) Đúng.}]{Một ô tô đang chạy với vận tốc $15$ m/s thì người lái đạp phanh. Từ thời điểm đó, ô tô chuyển động chậm dần đều và quãng đường đi được sau $t$ giây là $s(t)=15t-3t^2$ (mét). Xét tính đúng sai của các mệnh đề sau:}{%
  \y{Quãng đường ô tô đi được từ lúc đạp phanh đến khi dừng là $\dfrac{75}{8}\,\mathrm{m}$.}{s}
  \y{Gia tốc của ô tô là $-4\,\mathrm{m/s^2}$.}{s}
  \y{Sau $1$ giây kể từ lúc đạp phanh, vận tốc ô tô là $5\,\mathrm{m/s}$.}{s}
  \y{Ô tô dừng hẳn sau $\dfrac{5}{2}\,\mathrm{s}$ kể từ lúc đạp phanh.}{d}}

\cauDS[id=C1-B5-D01-TH-060,muc=TH,loigiai={Vận tốc $v(t)=s'(t)=-t^{2}+20t$, gia tốc $a(t)=v'(t)=-2t+20$. Do đó $a(5)=10$. Vận tốc $v(t)=s'(t)=-t^{2}+20t$, gia tốc $a(t)=v'(t)=-2t+20$. $v'(t)=0\Leftrightarrow t=10$; $v$ đồng biến trên $(0;10)$, nghịch biến trên $(10;20)$ nên $v_{\max}=v(10)=100$. Vận tốc $v(t)=s'(t)=-t^{2}+20t$, gia tốc $a(t)=v'(t)=-2t+20$. $v'(t)=0\Leftrightarrow t=10$ và $v'$ đổi dấu từ dương sang âm nên vận tốc lớn nhất tại $t=10$. Vận tốc $v(t)=s'(t)=-t^{2}+20t$, gia tốc $a(t)=v'(t)=-2t+20$. Do đó $v(5)=75$. a) Sai: giá trị đúng là $10\,\mathrm{m/s^2}$. b) Đúng. c) Đúng. d) Đúng.}]{Một vật chuyển động theo quy luật $s(t)=-\dfrac{1}{3}t^3+10t^2$, trong đó $t$ (giây) là khoảng thời gian tính từ lúc vật bắt đầu chuyển động và $s$ (mét) là quãng đường vật đi được trong khoảng thời gian đó, $0\le t\le 20$. Xét tính đúng sai của các mệnh đề sau:}{%
  \y{Gia tốc của vật tại thời điểm $t=5$ giây là $6\,\mathrm{m/s^2}$.}{s}
  \y{Vận tốc lớn nhất của vật là $100\,\mathrm{m/s}$.}{d}
  \y{Vận tốc của vật tại thời điểm $t=5$ giây là $75\,\mathrm{m/s}$.}{d}
  \y{Vật đạt vận tốc lớn nhất tại thời điểm $t=10\,\mathrm{s}$.}{d}}

\cauDS[id=C1-B5-D01-TH-061,muc=TH,loigiai={Vận tốc $v(t)=s'(t)=-2t^{2}+8t+2$, gia tốc $a(t)=v'(t)=-4t+8$. $v'(t)=0\Leftrightarrow t=2$ và $v'$ đổi dấu từ dương sang âm nên vận tốc lớn nhất tại $t=2$. Vận tốc $v(t)=s'(t)=-2t^{2}+8t+2$, gia tốc $a(t)=v'(t)=-4t+8$. $v'(t)=0\Leftrightarrow t=2$; $v$ đồng biến trên $(0;2)$, nghịch biến trên $(2;6)$ nên $v_{\max}=v(2)=10$. Vận tốc $v(t)=s'(t)=-2t^{2}+8t+2$, gia tốc $a(t)=v'(t)=-4t+8$. Do đó $a(2)=0$. Quãng đường $s(2)=\dfrac{44}{3}$. a) Sai: giá trị đúng là $0\,\mathrm{m/s^2}$. b) Đúng. c) Sai: giá trị đúng là $10\,\mathrm{m/s}$. d) Đúng.}]{Một vật chuyển động theo quy luật $s(t)=-\dfrac{2}{3}t^3+4t^2+2t$, trong đó $t$ (giây) là khoảng thời gian tính từ lúc vật bắt đầu chuyển động và $s$ (mét) là quãng đường vật đi được trong khoảng thời gian đó, $0\le t\le 6$. Xét tính đúng sai của các mệnh đề sau:}{%
  \y{Gia tốc của vật tại thời điểm $t=2$ giây là $2\,\mathrm{m/s^2}$.}{s}
  \y{Vật đạt vận tốc lớn nhất tại thời điểm $t=2\,\mathrm{s}$.}{d}
  \y{Vận tốc lớn nhất của vật là $12\,\mathrm{m/s}$.}{s}
  \y{Quãng đường vật đi được trong $2$ giây đầu là $\dfrac{44}{3}\,\mathrm{m}$.}{d}}

\cauDS[id=C1-B5-D01-TH-062,muc=TH,loigiai={Ta có $v(t)=h'(t)=10-10t$. $h'(t)=0\Leftrightarrow t=1$, $h$ đạt giá trị lớn nhất bằng $h(1)=5$. Ta có $v(t)=h'(t)=10-10t$. Do đó $v(2)=-10$. Ta có $v(t)=h'(t)=10-10t$. $h'(t)=0\Leftrightarrow t=1$. Ta có $h(2)=0$. a) Sai: giá trị đúng là $1\,\mathrm{s}$. b) Sai: giá trị đúng là $5\,\mathrm{m}$. c) Đúng. d) Đúng.}]{Một quả bóng được ném thẳng đứng lên trên từ độ cao $0$ m với vận tốc ban đầu $10$ m/s. Độ cao của quả bóng (tính bằng mét) sau $t$ giây là $h(t)=10t-5t^2$ (bỏ qua sức cản không khí, lấy $g=10$ m/s$^2$). Xét tính đúng sai của các mệnh đề sau:}{%
  \y{Quả bóng đạt độ cao lớn nhất tại thời điểm $t=3\,\mathrm{s}$.}{s}
  \y{Độ cao lớn nhất của quả bóng là $4\,\mathrm{m}$.}{s}
  \y{Vận tốc của quả bóng tại thời điểm $t=2$ giây là $-10\,\mathrm{m/s}$.}{d}
  \y{Sau $2$ giây quả bóng ở độ cao $0\,\mathrm{m}$.}{d}}

\cauDS[id=C1-B5-D01-TH-063,muc=TH,loigiai={Vận tốc $v(t)=x'(t)=-10 \pi\sin(\pi t)$ (cm/s), gia tốc $a(t)=v'(t)=-10 \pi^{2}\cos(\pi t)$ (cm/s$^2$). Vì $|\sin(\pi t)|\le1$ nên $|v(t)|\le 10 \pi$, dấu bằng xảy ra khi $\sin(\pi t)=\pm1$. Ta có $x(\dfrac{1}{12})=\dfrac{5 \sqrt{2}}{2} + \dfrac{5 \sqrt{6}}{2}$. Vận tốc $v(t)=x'(t)=-10 \pi\sin(\pi t)$ (cm/s), gia tốc $a(t)=v'(t)=-10 \pi^{2}\cos(\pi t)$ (cm/s$^2$). $v(\dfrac{1}{12})=\dfrac{5 \pi (- \sqrt{6} + \sqrt{2})}{2}$. Vận tốc $v(t)=x'(t)=-10 \pi\sin(\pi t)$ (cm/s), gia tốc $a(t)=v'(t)=-10 \pi^{2}\cos(\pi t)$ (cm/s$^2$). $a(\dfrac{1}{12})=\dfrac{5 \pi^{2} (- \sqrt{6} - \sqrt{2})}{2}$. a) Đúng. b) Đúng. c) Sai: giá trị đúng là $\dfrac{5 \sqrt{2}}{2} + \dfrac{5 \sqrt{6}}{2}\,\mathrm{cm}$. d) Sai: giá trị đúng là $\dfrac{5 \pi (- \sqrt{6} + \sqrt{2})}{2}\,\mathrm{cm/s}$.}]{Một vật dao động điều hòa có phương trình li độ $x(t)=10\cos(\pi t)$ (cm), với $t$ tính bằng giây. Xét tính đúng sai của các mệnh đề sau:}{%
  \y{Gia tốc của vật tại thời điểm $t=\dfrac{1}{12}$ giây là $\dfrac{5 \pi^{2} (- \sqrt{6} - \sqrt{2})}{2}\,\mathrm{cm/s^2}$.}{d}
  \y{Tốc độ cực đại của vật là $10 \pi\,\mathrm{cm/s}$.}{d}
  \y{Li độ của vật tại thời điểm $t=\dfrac{1}{12}$ giây là $5 \sqrt{2} + 5 \sqrt{6}\,\mathrm{cm}$.}{s}
  \y{Vận tốc của vật tại thời điểm $t=\dfrac{1}{12}$ giây là $\dfrac{15 \pi (- \sqrt{6} + \sqrt{2})}{4}\,\mathrm{cm/s}$.}{s}}

\cauDS[id=C1-B5-D01-TH-064,muc=TH,loigiai={Vận tốc $v(t)=x'(t)=-3 \pi\sin(\frac{\pi}{2} t)$ (cm/s), gia tốc $a(t)=v'(t)=-\dfrac{3 \pi^{2}}{2}\cos(\frac{\pi}{2} t)$ (cm/s$^2$). $a(\dfrac{1}{12})=\dfrac{3 \pi^{2} (- \sqrt{3} \sqrt{\sqrt{2} + 2} - \sqrt{2 - \sqrt{2}})}{8}$. Ta có $x(\dfrac{1}{12})=\dfrac{3 \sqrt{2 - \sqrt{2}}}{2} + \dfrac{3 \sqrt{3 \sqrt{2} + 6}}{2}$. Vận tốc $v(t)=x'(t)=-3 \pi\sin(\frac{\pi}{2} t)$ (cm/s), gia tốc $a(t)=v'(t)=-\dfrac{3 \pi^{2}}{2}\cos(\frac{\pi}{2} t)$ (cm/s$^2$). Vì $|\sin(\frac{\pi}{2} t)|\le1$ nên $|v(t)|\le 3 \pi$, dấu bằng xảy ra khi $\sin(\frac{\pi}{2} t)=\pm1$. Vận tốc $v(t)=x'(t)=-3 \pi\sin(\frac{\pi}{2} t)$ (cm/s), gia tốc $a(t)=v'(t)=-\dfrac{3 \pi^{2}}{2}\cos(\frac{\pi}{2} t)$ (cm/s$^2$). $v(\dfrac{1}{12})=\dfrac{3 \pi (- \sqrt{\sqrt{2} + 2} + \sqrt{6 - 3 \sqrt{2}})}{4}$. a) Đúng. b) Sai: giá trị đúng là $\dfrac{3 \pi^{2} (- \sqrt{3} \sqrt{\sqrt{2} + 2} - \sqrt{2 - \sqrt{2}})}{8}\,\mathrm{cm/s^2}$. c) Sai: giá trị đúng là $3 \pi\,\mathrm{cm/s}$. d) Sai: giá trị đúng là $\dfrac{3 \pi (- \sqrt{\sqrt{2} + 2} + \sqrt{6 - 3 \sqrt{2}})}{4}\,\mathrm{cm/s}$.}]{Một vật dao động điều hòa có phương trình li độ $x(t)=6\cos(\frac{\pi}{2} t)$ (cm), với $t$ tính bằng giây. Xét tính đúng sai của các mệnh đề sau:}{%
  \y{Li độ của vật tại thời điểm $t=\dfrac{1}{12}$ giây là $\dfrac{3 \sqrt{2 - \sqrt{2}}}{2} + \dfrac{3 \sqrt{3 \sqrt{2} + 6}}{2}\,\mathrm{cm}$.}{d}
  \y{Gia tốc của vật tại thời điểm $t=\dfrac{1}{12}$ giây là $- \dfrac{3 \pi^{2} (\sqrt{2 - \sqrt{2}} + \sqrt{3} \sqrt{\sqrt{2} + 2})}{8} + 1\,\mathrm{cm/s^2}$.}{s}
  \y{Tốc độ cực đại của vật là $\dfrac{3 \pi}{2}\,\mathrm{cm/s}$.}{s}
  \y{Vận tốc của vật tại thời điểm $t=\dfrac{1}{12}$ giây là $\dfrac{3 \pi (- \sqrt{\sqrt{2} + 2} + \sqrt{6 - 3 \sqrt{2}})}{4} + 1\,\mathrm{cm/s}$.}{s}}

\cauTLN[id=C1-B5-D01-TH-065,muc=TH,loigiai={Vận tốc $v(t)=s'(t)=24-3t$. $v(2)=18$. Kết quả: $18$.}]{Một ô tô đang chạy với vận tốc $24$ m/s thì người lái đạp phanh. Từ thời điểm đó, ô tô chuyển động chậm dần đều và quãng đường đi được sau $t$ giây là $s(t)=24t-\dfrac{3}{2}t^2$ (mét). Tìm vận tốc của ô tô sau $2$ giây kể từ lúc đạp phanh (đơn vị: $\mathrm{m/s}$).}{18}

\cauTLN[id=C1-B5-D01-TH-066,muc=TH,loigiai={Vận tốc $v(t)=x'(t)=-32 \pi\sin(4 \pi t)$ (cm/s), gia tốc $a(t)=v'(t)=-128 \pi^{2}\cos(4 \pi t)$ (cm/s$^2$). Vì $|\sin(4 \pi t)|\le1$ nên $|v(t)|\le 32 \pi$, dấu bằng xảy ra khi $\sin(4 \pi t)=\pm1$. Kết quả: $101$.}]{Một vật dao động điều hòa có phương trình li độ $x(t)=8\cos(4 \pi t)$ (cm), với $t$ tính bằng giây. Tìm tốc độ cực đại của vật (giá trị lớn nhất của $|v(t)|$) (đơn vị: $\mathrm{cm/s}$) (làm tròn kết quả đến hàng đơn vị).}{101}

\cauTLN[id=C1-B5-D01-TH-067,muc=TH,loigiai={Ta có $v(t)=h'(t)=10-10t$. $h'(t)=0\Leftrightarrow t=1$. Kết quả: $1$.}]{Một quả bóng được ném thẳng đứng lên trên từ độ cao $0$ m với vận tốc ban đầu $10$ m/s. Độ cao của quả bóng (tính bằng mét) sau $t$ giây là $h(t)=10t-5t^2$ (bỏ qua sức cản không khí, lấy $g=10$ m/s$^2$). Tìm thời điểm quả bóng đạt độ cao lớn nhất (đơn vị: $\mathrm{s}$).}{1}

\cauTLN[id=C1-B5-D01-TH-068,muc=TH,loigiai={Vận tốc $v(t)=s'(t)=-3t^{2}+24t+2$, gia tốc $a(t)=v'(t)=-6t+24$. Do đó $v(1)=23$. Kết quả: $23$.}]{Một vật chuyển động theo quy luật $s(t)=-\dfrac{3}{3}t^3+12t^2+2t$, trong đó $t$ (giây) là khoảng thời gian tính từ lúc vật bắt đầu chuyển động và $s$ (mét) là quãng đường vật đi được trong khoảng thời gian đó, $0\le t\le 6$. Tìm vận tốc của vật tại thời điểm $t=1$ giây (đơn vị: $\mathrm{m/s}$).}{23}

\cauTLN[id=C1-B5-D01-TH-069,muc=TH,loigiai={Vận tốc $v(t)=x'(t)=-5 \pi\sin(\pi t)$ (cm/s), gia tốc $a(t)=v'(t)=-5 \pi^{2}\cos(\pi t)$ (cm/s$^2$). $a(\dfrac{1}{3})=- \dfrac{5 \pi^{2}}{2}$. Kết quả: $-25$.}]{Một vật dao động điều hòa có phương trình li độ $x(t)=5\cos(\pi t)$ (cm), với $t$ tính bằng giây. Tìm gia tốc của vật tại thời điểm $t=\dfrac{1}{3}$ giây (đơn vị: $\mathrm{cm/s^2}$) (làm tròn kết quả đến hàng đơn vị).}{-25}

\cauTLN[id=C1-B5-D01-TH-070,muc=TH,loigiai={Vận tốc $v(t)=s'(t)=25-3t$. Ô tô dừng khi $v(t)=0\Leftrightarrow t=\dfrac{25}{3}$. Kết quả: $8{,}33$.}]{Một ô tô đang chạy với vận tốc $25$ m/s thì người lái đạp phanh. Từ thời điểm đó, ô tô chuyển động chậm dần đều và quãng đường đi được sau $t$ giây là $s(t)=25t-\dfrac{3}{2}t^2$ (mét). Tìm thời gian từ lúc đạp phanh đến khi ô tô dừng hẳn (đơn vị: $\mathrm{s}$) (làm tròn kết quả đến hàng phần trăm).}{8,33}

\cauTLN[id=C1-B5-D01-TH-071,muc=TH,loigiai={Ta có $x(\dfrac{1}{2})=0$. Kết quả: $0$.}]{Một vật dao động điều hòa có phương trình li độ $x(t)=10\cos(\pi t)$ (cm), với $t$ tính bằng giây. Tìm li độ của vật tại thời điểm $t=\dfrac{1}{2}$ giây (đơn vị: $\mathrm{cm}$).}{0}

\cauTLN[id=C1-B5-D01-TH-072,muc=TH,loigiai={Vận tốc $v(t)=s'(t)=-t^{2}+24t+2$, gia tốc $a(t)=v'(t)=-2t+24$. Do đó $a(9)=6$. Kết quả: $6$.}]{Một vật chuyển động theo quy luật $s(t)=-\dfrac{1}{3}t^3+12t^2+2t$, trong đó $t$ (giây) là khoảng thời gian tính từ lúc vật bắt đầu chuyển động và $s$ (mét) là quãng đường vật đi được trong khoảng thời gian đó, $0\le t\le 14$. Tìm gia tốc của vật tại thời điểm $t=9$ giây (đơn vị: $\mathrm{m/s^2}$).}{6}

\cauTLN[id=C1-B5-D01-TH-073,muc=TH,loigiai={Ta có $x(\dfrac{1}{4})=0$. Kết quả: $0$.}]{Một vật dao động điều hòa có phương trình li độ $x(t)=5\cos(2 \pi t)$ (cm), với $t$ tính bằng giây. Tìm li độ của vật tại thời điểm $t=\dfrac{1}{4}$ giây (đơn vị: $\mathrm{cm}$).}{0}

\cauTLN[id=C1-B5-D01-TH-074,muc=TH,loigiai={Vận tốc $v(t)=s'(t)=30-2t$. $v(2)=26$. Kết quả: $26$.}]{Một ô tô đang chạy với vận tốc $30$ m/s thì người lái đạp phanh. Từ thời điểm đó, ô tô chuyển động chậm dần đều và quãng đường đi được sau $t$ giây là $s(t)=30t-t^2$ (mét). Tìm vận tốc của ô tô sau $2$ giây kể từ lúc đạp phanh (đơn vị: $\mathrm{m/s}$).}{26}

\cauTLN[id=C1-B5-D01-TH-075,muc=TH,loigiai={Vận tốc $v(t)=s'(t)=20-5t$. Gia tốc $a(t)=v'(t)=-5$. Kết quả: $-5$.}]{Một ô tô đang chạy với vận tốc $20$ m/s thì người lái đạp phanh. Từ thời điểm đó, ô tô chuyển động chậm dần đều và quãng đường đi được sau $t$ giây là $s(t)=20t-\dfrac{5}{2}t^2$ (mét). Tìm gia tốc của ô tô (đơn vị: $\mathrm{m/s^2}$).}{-5}

\cauTLN[id=C1-B5-D01-TH-076,muc=TH,loigiai={Ta có $h(1)=22$. Kết quả: $22$.}]{Một quả bóng được ném thẳng đứng lên trên từ độ cao $2$ m với vận tốc ban đầu $25$ m/s. Độ cao của quả bóng (tính bằng mét) sau $t$ giây là $h(t)=2+25t-5t^2$ (bỏ qua sức cản không khí, lấy $g=10$ m/s$^2$). Tìm độ cao của quả bóng tại thời điểm $t=1$ giây (đơn vị: $\mathrm{m}$).}{22}

\cauTLN[id=C1-B5-D01-TH-077,muc=TH,loigiai={Vận tốc $v(t)=x'(t)=-3 \pi\sin(\pi t)$ (cm/s), gia tốc $a(t)=v'(t)=-3 \pi^{2}\cos(\pi t)$ (cm/s$^2$). $a(\dfrac{1}{6})=- \dfrac{3 \sqrt{3} \pi^{2}}{2}$. Kết quả: $-26$.}]{Một vật dao động điều hòa có phương trình li độ $x(t)=3\cos(\pi t)$ (cm), với $t$ tính bằng giây. Tìm gia tốc của vật tại thời điểm $t=\dfrac{1}{6}$ giây (đơn vị: $\mathrm{cm/s^2}$) (làm tròn kết quả đến hàng đơn vị).}{-26}

\cauTLN[id=C1-B5-D01-TH-078,muc=TH,loigiai={Ta có $v(t)=h'(t)=25-10t$. $h'(t)=0\Leftrightarrow t=\dfrac{5}{2}$, $h$ đạt giá trị lớn nhất bằng $h(\dfrac{5}{2})=\dfrac{133}{4}$. Kết quả: $33{,}3$.}]{Một quả bóng được ném thẳng đứng lên trên từ độ cao $2$ m với vận tốc ban đầu $25$ m/s. Độ cao của quả bóng (tính bằng mét) sau $t$ giây là $h(t)=2+25t-5t^2$ (bỏ qua sức cản không khí, lấy $g=10$ m/s$^2$). Tìm độ cao lớn nhất mà quả bóng đạt được (đơn vị: $\mathrm{m}$) (làm tròn kết quả đến hàng phần mười).}{33,3}

\cauTLN[id=C1-B5-D01-TH-079,muc=TH,loigiai={Vận tốc $v(t)=x'(t)=-2 \pi\sin(\frac{\pi}{2} t)$ (cm/s), gia tốc $a(t)=v'(t)=-\pi^{2}\cos(\frac{\pi}{2} t)$ (cm/s$^2$). $a(\dfrac{1}{4})=- \dfrac{\pi^{2} \sqrt{\sqrt{2} + 2}}{2}$. Kết quả: $-9{,}1$.}]{Một vật dao động điều hòa có phương trình li độ $x(t)=4\cos(\frac{\pi}{2} t)$ (cm), với $t$ tính bằng giây. Tìm gia tốc của vật tại thời điểm $t=\dfrac{1}{4}$ giây (đơn vị: $\mathrm{cm/s^2}$) (làm tròn kết quả đến hàng phần mười).}{-9,1}

\cauTLN[id=C1-B5-D01-TH-080,muc=TH,loigiai={Quãng đường $s(5)=650$. Kết quả: $650$.}]{Một vật chuyển động theo quy luật $s(t)=-\dfrac{3}{3}t^3+30t^2+5t$, trong đó $t$ (giây) là khoảng thời gian tính từ lúc vật bắt đầu chuyển động và $s$ (mét) là quãng đường vật đi được trong khoảng thời gian đó, $0\le t\le 20$. Tìm quãng đường vật đi được trong $5$ giây đầu tiên (đơn vị: $\mathrm{m}$).}{650}

\cauTN[id=C1-B5-D01-VD-001,muc=VD,loigiai={Vận tốc $v(t)=s'(t)=-2t^{2}+36t+5$, gia tốc $a(t)=v'(t)=-4t+36$. Do đó $a(9)=0$.}]{Một vật chuyển động theo quy luật $s(t)=-\dfrac{2}{3}t^3+18t^2+5t$, trong đó $t$ (giây) là khoảng thời gian tính từ lúc vật bắt đầu chuyển động và $s$ (mét) là quãng đường vật đi được trong khoảng thời gian đó, $0\le t\le 12$. Gia tốc của vật tại thời điểm $t=9$ giây bằng}
  {$0\,\mathrm{m/s^2}$.}
  {$2\,\mathrm{m/s^2}$.}{$-2\,\mathrm{m/s^2}$.}{$-4\,\mathrm{m/s^2}$.}

\cauTN[id=C1-B5-D01-VD-002,muc=VD,loigiai={Vận tốc $v(t)=x'(t)=-8 \pi\sin(4 \pi t)$ (cm/s), gia tốc $a(t)=v'(t)=-32 \pi^{2}\cos(4 \pi t)$ (cm/s$^2$). Vì $|\sin(4 \pi t)|\le1$ nên $|v(t)|\le 8 \pi$, dấu bằng xảy ra khi $\sin(4 \pi t)=\pm1$.}]{Một vật dao động điều hòa có phương trình li độ $x(t)=2\cos(4 \pi t)$ (cm), với $t$ tính bằng giây. Tốc độ cực đại của vật (giá trị lớn nhất của $|v(t)|$) bằng}
  {$8 \pi\,\mathrm{cm/s}$.}
  {$-1 + 8 \pi\,\mathrm{cm/s}$.}{$12 \pi\,\mathrm{cm/s}$.}{$16 \pi\,\mathrm{cm/s}$.}

\cauTN[id=C1-B5-D01-VD-003,muc=VD,loigiai={Vận tốc $v(t)=x'(t)=3t^{2}-18t+15=3(t-1)(t-5)$. $v'(t)=6t-18=0\Leftrightarrow t=3$, $v_{\min}=v(3)=-12$.}]{Một chất điểm chuyển động thẳng có phương trình $x(t)=t^{3}-9t^{2}+15t$ (mét), với $t$ (giây), $0\le t\le 7$. Vận tốc nhỏ nhất của chất điểm (vận tốc có thể âm) bằng}
  {$-12\,\mathrm{m/s}$.}
  {$-6\,\mathrm{m/s}$.}{$-24\,\mathrm{m/s}$.}{$-10\,\mathrm{m/s}$.}

\cauTN[id=C1-B5-D01-VD-004,muc=VD,loigiai={Vận tốc $v(t)=s'(t)=-t^{2}+16t+5$, gia tốc $a(t)=v'(t)=-2t+16$. $v'(t)=0\Leftrightarrow t=8$ và $v'$ đổi dấu từ dương sang âm nên vận tốc lớn nhất tại $t=8$.}]{Một vật chuyển động theo quy luật $s(t)=-\dfrac{1}{3}t^3+8t^2+5t$, trong đó $t$ (giây) là khoảng thời gian tính từ lúc vật bắt đầu chuyển động và $s$ (mét) là quãng đường vật đi được trong khoảng thời gian đó, $0\le t\le 14$. Thời điểm vật đạt vận tốc lớn nhất là}
  {$8\,\mathrm{s}$.}
  {$69\,\mathrm{s}$.}{$16\,\mathrm{s}$.}{$7\,\mathrm{s}$.}

\cauTN[id=C1-B5-D01-VD-005,muc=VD,loigiai={Vận tốc $v(t)=s'(t)=-3t^{2}+18t+2$, gia tốc $a(t)=v'(t)=-6t+18$. $v'(t)=0\Leftrightarrow t=3$ và $v'$ đổi dấu từ dương sang âm nên vận tốc lớn nhất tại $t=3$.}]{Một vật chuyển động theo quy luật $s(t)=-\dfrac{3}{3}t^3+9t^2+2t$, trong đó $t$ (giây) là khoảng thời gian tính từ lúc vật bắt đầu chuyển động và $s$ (mét) là quãng đường vật đi được trong khoảng thời gian đó, $0\le t\le 6$. Thời điểm vật đạt vận tốc lớn nhất là}
  {$3\,\mathrm{s}$.}
  {$29\,\mathrm{s}$.}{$6\,\mathrm{s}$.}{$7\,\mathrm{s}$.}

\cauTN[id=C1-B5-D01-VD-006,muc=VD,loigiai={Vận tốc $v(t)=x'(t)=-40 \pi\sin(4 \pi t)$ (cm/s), gia tốc $a(t)=v'(t)=-160 \pi^{2}\cos(4 \pi t)$ (cm/s$^2$). $v(\dfrac{1}{2})=0$.}]{Một vật dao động điều hòa có phương trình li độ $x(t)=10\cos(4 \pi t)$ (cm), với $t$ tính bằng giây. Vận tốc của vật tại thời điểm $t=\dfrac{1}{2}$ giây bằng}
  {$0\,\mathrm{cm/s}$.}
  {$2\,\mathrm{cm/s}$.}{$4\,\mathrm{cm/s}$.}{$-4\,\mathrm{cm/s}$.}

\cauTN[id=C1-B5-D01-VD-007,muc=VD,loigiai={Vận tốc $v(t)=x'(t)=3t^{2}-12t+9=3(t-1)(t-3)$. Gia tốc $a(t)=6t-12$, $a(3)=6$.}]{Một chất điểm chuyển động thẳng có phương trình $x(t)=t^{3}-6t^{2}+9t$ (mét), với $t$ (giây), $0\le t\le 4$. Gia tốc của chất điểm tại thời điểm vận tốc bằng $0$ lần thứ hai bằng}
  {$6\,\mathrm{m/s^2}$.}
  {$3\,\mathrm{m/s^2}$.}{$8\,\mathrm{m/s^2}$.}{$2\,\mathrm{m/s^2}$.}

\cauTN[id=C1-B5-D01-VD-008,muc=VD,loigiai={Ta có $x(\dfrac{1}{3})=\dfrac{5}{2}$.}]{Một vật dao động điều hòa có phương trình li độ $x(t)=5\cos(\pi t)$ (cm), với $t$ tính bằng giây. Li độ của vật tại thời điểm $t=\dfrac{1}{3}$ giây bằng}
  {$\dfrac{5}{2}\,\mathrm{cm}$.}
  {$\dfrac{2}{5}\,\mathrm{cm}$.}{$\dfrac{5}{4}\,\mathrm{cm}$.}{$5\,\mathrm{cm}$.}

\cauTN[id=C1-B5-D01-VD-009,muc=VD,loigiai={Vận tốc $v(t)=s'(t)=-2t^{2}+32t$, gia tốc $a(t)=v'(t)=-4t+32$. Do đó $v(10)=120$.}]{Một vật chuyển động theo quy luật $s(t)=-\dfrac{2}{3}t^3+16t^2$, trong đó $t$ (giây) là khoảng thời gian tính từ lúc vật bắt đầu chuyển động và $s$ (mét) là quãng đường vật đi được trong khoảng thời gian đó, $0\le t\le 13$. Vận tốc của vật tại thời điểm $t=10$ giây bằng}
  {$120\,\mathrm{m/s}$.}
  {$160\,\mathrm{m/s}$.}{$240\,\mathrm{m/s}$.}{$60\,\mathrm{m/s}$.}

\cauTN[id=C1-B5-D01-VD-010,muc=VD,loigiai={Vận tốc $v(t)=x'(t)=-20 \pi\sin(4 \pi t)$ (cm/s), gia tốc $a(t)=v'(t)=-80 \pi^{2}\cos(4 \pi t)$ (cm/s$^2$). $v(\dfrac{1}{6})=- 10 \sqrt{3} \pi$.}]{Một vật dao động điều hòa có phương trình li độ $x(t)=5\cos(4 \pi t)$ (cm), với $t$ tính bằng giây. Vận tốc của vật tại thời điểm $t=\dfrac{1}{6}$ giây bằng}
  {$- 10 \sqrt{3} \pi\,\mathrm{cm/s}$.}
  {$- 5 \sqrt{3} \pi\,\mathrm{cm/s}$.}{$- 15 \sqrt{3} \pi\,\mathrm{cm/s}$.}{$- 20 \sqrt{3} \pi\,\mathrm{cm/s}$.}

\cauTN[id=C1-B5-D01-VD-011,muc=VD,loigiai={Vận tốc $v(t)=x'(t)=-16 \pi\sin(4 \pi t)$ (cm/s), gia tốc $a(t)=v'(t)=-64 \pi^{2}\cos(4 \pi t)$ (cm/s$^2$). $a(\dfrac{1}{6})=32 \pi^{2}$.}]{Một vật dao động điều hòa có phương trình li độ $x(t)=4\cos(4 \pi t)$ (cm), với $t$ tính bằng giây. Gia tốc của vật tại thời điểm $t=\dfrac{1}{6}$ giây bằng}
  {$32 \pi^{2}\,\mathrm{cm/s^2}$.}
  {$\dfrac{64 \pi^{2}}{3}\,\mathrm{cm/s^2}$.}{$-1 + 32 \pi^{2}\,\mathrm{cm/s^2}$.}{$1 + 32 \pi^{2}\,\mathrm{cm/s^2}$.}

\cauTN[id=C1-B5-D01-VD-012,muc=VD,loigiai={Vận tốc $v(t)=s'(t)=-t^{2}+18t+5$, gia tốc $a(t)=v'(t)=-2t+18$. $v'(t)=0\Leftrightarrow t=9$; $v$ đồng biến trên $(0;9)$, nghịch biến trên $(9;18)$ nên $v_{\max}=v(9)=86$.}]{Một vật chuyển động theo quy luật $s(t)=-\dfrac{1}{3}t^3+9t^2+5t$, trong đó $t$ (giây) là khoảng thời gian tính từ lúc vật bắt đầu chuyển động và $s$ (mét) là quãng đường vật đi được trong khoảng thời gian đó, $0\le t\le 18$. Vận tốc lớn nhất của vật trong khoảng thời gian trên bằng}
  {$86\,\mathrm{m/s}$.}
  {$9\,\mathrm{m/s}$.}{$5\,\mathrm{m/s}$.}{$87\,\mathrm{m/s}$.}

\cauTN[id=C1-B5-D01-VD-013,muc=VD,loigiai={Vận tốc $v(t)=x'(t)=-32 \pi\sin(4 \pi t)$ (cm/s), gia tốc $a(t)=v'(t)=-128 \pi^{2}\cos(4 \pi t)$ (cm/s$^2$). Vì $|\sin(4 \pi t)|\le1$ nên $|v(t)|\le 32 \pi$, dấu bằng xảy ra khi $\sin(4 \pi t)=\pm1$.}]{Một vật dao động điều hòa có phương trình li độ $x(t)=8\cos(4 \pi t)$ (cm), với $t$ tính bằng giây. Tốc độ cực đại của vật (giá trị lớn nhất của $|v(t)|$) bằng}
  {$32 \pi\,\mathrm{cm/s}$.}
  {$-1 + 32 \pi\,\mathrm{cm/s}$.}{$16 \pi\,\mathrm{cm/s}$.}{$48 \pi\,\mathrm{cm/s}$.}

\cauTN[id=C1-B5-D01-VD-014,muc=VD,loigiai={Vận tốc $v(t)=x'(t)=3t^{2}-24t+45=3(t-3)(t-5)$. $v'(t)=6t-24=0\Leftrightarrow t=4$, $v_{\min}=v(4)=-3$.}]{Một chất điểm chuyển động thẳng có phương trình $x(t)=t^{3}-12t^{2}+45t$ (mét), với $t$ (giây), $0\le t\le 7$. Vận tốc nhỏ nhất của chất điểm (vận tốc có thể âm) bằng}
  {$-3\,\mathrm{m/s}$.}
  {$-5\,\mathrm{m/s}$.}{$3\,\mathrm{m/s}$.}{$-4\,\mathrm{m/s}$.}

\cauTN[id=C1-B5-D01-VD-015,muc=VD,loigiai={Vận tốc $v(t)=s'(t)=-t^{2}+20t$, gia tốc $a(t)=v'(t)=-2t+20$. Do đó $a(10)=0$.}]{Một vật chuyển động theo quy luật $s(t)=-\dfrac{1}{3}t^3+10t^2$, trong đó $t$ (giây) là khoảng thời gian tính từ lúc vật bắt đầu chuyển động và $s$ (mét) là quãng đường vật đi được trong khoảng thời gian đó, $0\le t\le 12$. Gia tốc của vật tại thời điểm $t=10$ giây bằng}
  {$0\,\mathrm{m/s^2}$.}
  {$-1\,\mathrm{m/s^2}$.}{$2\,\mathrm{m/s^2}$.}{$-2\,\mathrm{m/s^2}$.}

\cauTN[id=C1-B5-D01-VD-016,muc=VD,loigiai={Vận tốc $v(t)=s'(t)=-t^{2}+18t$, gia tốc $a(t)=v'(t)=-2t+18$. $v'(t)=0\Leftrightarrow t=9$ và $v'$ đổi dấu từ dương sang âm nên vận tốc lớn nhất tại $t=9$.}]{Một vật chuyển động theo quy luật $s(t)=-\dfrac{1}{3}t^3+9t^2$, trong đó $t$ (giây) là khoảng thời gian tính từ lúc vật bắt đầu chuyển động và $s$ (mét) là quãng đường vật đi được trong khoảng thời gian đó, $0\le t\le 13$. Thời điểm vật đạt vận tốc lớn nhất là}
  {$9\,\mathrm{s}$.}
  {$81\,\mathrm{s}$.}{$18\,\mathrm{s}$.}{$11\,\mathrm{s}$.}

\cauTN[id=C1-B5-D01-VD-017,muc=VD,loigiai={Vận tốc $v(t)=x'(t)=3t^{2}-18t+24=3(t-2)(t-4)$. $v'(t)=6t-18=0\Leftrightarrow t=3$, $v_{\min}=v(3)=-3$.}]{Một chất điểm chuyển động thẳng có phương trình $x(t)=t^{3}-9t^{2}+24t$ (mét), với $t$ (giây), $0\le t\le 5$. Vận tốc nhỏ nhất của chất điểm (vận tốc có thể âm) bằng}
  {$-3\,\mathrm{m/s}$.}
  {$-2\,\mathrm{m/s}$.}{$-1\,\mathrm{m/s}$.}{$3\,\mathrm{m/s}$.}

\cauTN[id=C1-B5-D01-VD-018,muc=VD,loigiai={Vận tốc $v(t)=s'(t)=-t^{2}+16t$, gia tốc $a(t)=v'(t)=-2t+16$. $v'(t)=0\Leftrightarrow t=8$; $v$ đồng biến trên $(0;8)$, nghịch biến trên $(8;16)$ nên $v_{\max}=v(8)=64$.}]{Một vật chuyển động theo quy luật $s(t)=-\dfrac{1}{3}t^3+8t^2$, trong đó $t$ (giây) là khoảng thời gian tính từ lúc vật bắt đầu chuyển động và $s$ (mét) là quãng đường vật đi được trong khoảng thời gian đó, $0\le t\le 16$. Vận tốc lớn nhất của vật trong khoảng thời gian trên bằng}
  {$64\,\mathrm{m/s}$.}
  {$8\,\mathrm{m/s}$.}{$68\,\mathrm{m/s}$.}{$62\,\mathrm{m/s}$.}

\cauTN[id=C1-B5-D01-VD-019,muc=VD,loigiai={Vận tốc $v(t)=s'(t)=-t^{2}+12t+3$, gia tốc $a(t)=v'(t)=-2t+12$. Do đó $a(1)=10$.}]{Một vật chuyển động theo quy luật $s(t)=-\dfrac{1}{3}t^3+6t^2+3t$, trong đó $t$ (giây) là khoảng thời gian tính từ lúc vật bắt đầu chuyển động và $s$ (mét) là quãng đường vật đi được trong khoảng thời gian đó, $0\le t\le 12$. Gia tốc của vật tại thời điểm $t=1$ giây bằng}
  {$10\,\mathrm{m/s^2}$.}
  {$6\,\mathrm{m/s^2}$.}{$8\,\mathrm{m/s^2}$.}{$14\,\mathrm{m/s^2}$.}

\cauTN[id=C1-B5-D01-VD-020,muc=VD,loigiai={Vận tốc $v(t)=x'(t)=-\dfrac{3 \pi}{2}\sin(\frac{\pi}{2} t)$ (cm/s), gia tốc $a(t)=v'(t)=-\dfrac{3 \pi^{2}}{4}\cos(\frac{\pi}{2} t)$ (cm/s$^2$). $a(\dfrac{1}{4})=- \dfrac{3 \pi^{2} \sqrt{\sqrt{2} + 2}}{8}$.}]{Một vật dao động điều hòa có phương trình li độ $x(t)=3\cos(\frac{\pi}{2} t)$ (cm), với $t$ tính bằng giây. Gia tốc của vật tại thời điểm $t=\dfrac{1}{4}$ giây bằng}
  {$- \dfrac{3 \pi^{2} \sqrt{\sqrt{2} + 2}}{8}\,\mathrm{cm/s^2}$.}
  {$- \dfrac{3 \pi^{2} \sqrt{\sqrt{2} + 2}}{8} + 1\,\mathrm{cm/s^2}$.}{$- \dfrac{3 \pi^{2} \sqrt{\sqrt{2} + 2}}{16}\,\mathrm{cm/s^2}$.}{$- \dfrac{\pi^{2} \sqrt{\sqrt{2} + 2}}{4}\,\mathrm{cm/s^2}$.}

\cauTN[id=C1-B5-D01-VD-021,muc=VD,loigiai={Vận tốc $v(t)=x'(t)=3t^{2}-18t+24=3(t-2)(t-4)$. Chất điểm đổi chiều tại $t=2$ và $t=4$. Ta có $x(0)=0$, $x(2)=20$, $x(4)=16$, $x(5)=20$, nên quãng đường là $|20|+|-4|+|4|=28$.}]{Một chất điểm chuyển động thẳng có phương trình $x(t)=t^{3}-9t^{2}+24t$ (mét), với $t$ (giây), $0\le t\le 5$. Tổng quãng đường chất điểm đi được trong $5$ giây đầu bằng}
  {$28\,\mathrm{m}$.}
  {$33\,\mathrm{m}$.}{$56\,\mathrm{m}$.}{$14\,\mathrm{m}$.}

\cauTN[id=C1-B5-D01-VD-022,muc=VD,loigiai={Vận tốc $v(t)=s'(t)=-2t^{2}+16t$, gia tốc $a(t)=v'(t)=-4t+16$. Do đó $v(5)=30$.}]{Một vật chuyển động theo quy luật $s(t)=-\dfrac{2}{3}t^3+8t^2$, trong đó $t$ (giây) là khoảng thời gian tính từ lúc vật bắt đầu chuyển động và $s$ (mét) là quãng đường vật đi được trong khoảng thời gian đó, $0\le t\le 8$. Vận tốc của vật tại thời điểm $t=5$ giây bằng}
  {$30\,\mathrm{m/s}$.}
  {$31\,\mathrm{m/s}$.}{$28\,\mathrm{m/s}$.}{$60\,\mathrm{m/s}$.}

\cauTN[id=C1-B5-D01-VD-023,muc=VD,loigiai={Vận tốc $v(t)=s'(t)=-t^{2}+8t$, gia tốc $a(t)=v'(t)=-2t+8$. Do đó $v(1)=7$.}]{Một vật chuyển động theo quy luật $s(t)=-\dfrac{1}{3}t^3+4t^2$, trong đó $t$ (giây) là khoảng thời gian tính từ lúc vật bắt đầu chuyển động và $s$ (mét) là quãng đường vật đi được trong khoảng thời gian đó, $0\le t\le 8$. Vận tốc của vật tại thời điểm $t=1$ giây bằng}
  {$7\,\mathrm{m/s}$.}
  {$14\,\mathrm{m/s}$.}{$9\,\mathrm{m/s}$.}{$5\,\mathrm{m/s}$.}

\cauTN[id=C1-B5-D01-VD-024,muc=VD,loigiai={Ta có $x(\dfrac{1}{4})=0$.}]{Một vật dao động điều hòa có phương trình li độ $x(t)=8\cos(2 \pi t)$ (cm), với $t$ tính bằng giây. Li độ của vật tại thời điểm $t=\dfrac{1}{4}$ giây bằng}
  {$0\,\mathrm{cm}$.}
  {$4\,\mathrm{cm}$.}{$-4\,\mathrm{cm}$.}{$-2\,\mathrm{cm}$.}

\cauTN[id=C1-B5-D01-VD-025,muc=VD,loigiai={Vận tốc $v(t)=s'(t)=-t^{2}+18t+3$, gia tốc $a(t)=v'(t)=-2t+18$. Do đó $v(13)=68$.}]{Một vật chuyển động theo quy luật $s(t)=-\dfrac{1}{3}t^3+9t^2+3t$, trong đó $t$ (giây) là khoảng thời gian tính từ lúc vật bắt đầu chuyển động và $s$ (mét) là quãng đường vật đi được trong khoảng thời gian đó, $0\le t\le 15$. Vận tốc của vật tại thời điểm $t=13$ giây bằng}
  {$68\,\mathrm{m/s}$.}
  {$63\,\mathrm{m/s}$.}{$78\,\mathrm{m/s}$.}{$58\,\mathrm{m/s}$.}

\cauTN[id=C1-B5-D01-VD-026,muc=VD,loigiai={Vận tốc $v(t)=x'(t)=3t^{2}-18t+15=3(t-1)(t-5)$. $v(t)=0\Leftrightarrow t=1$ hoặc $t=5$; chất điểm dừng lần đầu tại $t=1$.}]{Một chất điểm chuyển động thẳng có phương trình $x(t)=t^{3}-9t^{2}+15t$ (mét), với $t$ (giây), $0\le t\le 6$. Thời điểm đầu tiên chất điểm dừng lại (đổi chiều chuyển động) là}
  {$1\,\mathrm{s}$.}
  {$5\,\mathrm{s}$.}{$2\,\mathrm{s}$.}{$3\,\mathrm{s}$.}

\cauTN[id=C1-B5-D01-VD-027,muc=VD,loigiai={Vận tốc $v(t)=s'(t)=-3t^{2}+12t+5$, gia tốc $a(t)=v'(t)=-6t+12$. Do đó $a(5)=-18$.}]{Một vật chuyển động theo quy luật $s(t)=-\dfrac{3}{3}t^3+6t^2+5t$, trong đó $t$ (giây) là khoảng thời gian tính từ lúc vật bắt đầu chuyển động và $s$ (mét) là quãng đường vật đi được trong khoảng thời gian đó, $0\le t\le 6$. Gia tốc của vật tại thời điểm $t=5$ giây bằng}
  {$-18\,\mathrm{m/s^2}$.}
  {$-20\,\mathrm{m/s^2}$.}{$-19\,\mathrm{m/s^2}$.}{$-16\,\mathrm{m/s^2}$.}

\cauTN[id=C1-B5-D01-VD-028,muc=VD,loigiai={Vận tốc $v(t)=x'(t)=-12 \pi\sin(2 \pi t)$ (cm/s), gia tốc $a(t)=v'(t)=-24 \pi^{2}\cos(2 \pi t)$ (cm/s$^2$). $v(\dfrac{1}{6})=- 6 \sqrt{3} \pi$.}]{Một vật dao động điều hòa có phương trình li độ $x(t)=6\cos(2 \pi t)$ (cm), với $t$ tính bằng giây. Vận tốc của vật tại thời điểm $t=\dfrac{1}{6}$ giây bằng}
  {$- 6 \sqrt{3} \pi\,\mathrm{cm/s}$.}
  {$- 9 \sqrt{3} \pi\,\mathrm{cm/s}$.}{$- 6 \sqrt{3} \pi - 1\,\mathrm{cm/s}$.}{$6 \sqrt{3} \pi\,\mathrm{cm/s}$.}

\cauTN[id=C1-B5-D01-VD-029,muc=VD,loigiai={Vận tốc $v(t)=s'(t)=-t^{2}+18t$, gia tốc $a(t)=v'(t)=-2t+18$. $v'(t)=0\Leftrightarrow t=9$ và $v'$ đổi dấu từ dương sang âm nên vận tốc lớn nhất tại $t=9$.}]{Một vật chuyển động theo quy luật $s(t)=-\dfrac{1}{3}t^3+9t^2$, trong đó $t$ (giây) là khoảng thời gian tính từ lúc vật bắt đầu chuyển động và $s$ (mét) là quãng đường vật đi được trong khoảng thời gian đó, $0\le t\le 12$. Thời điểm vật đạt vận tốc lớn nhất là}
  {$9\,\mathrm{s}$.}
  {$81\,\mathrm{s}$.}{$18\,\mathrm{s}$.}{$8\,\mathrm{s}$.}

\cauTN[id=C1-B5-D01-VD-030,muc=VD,loigiai={Vận tốc $v(t)=s'(t)=-t^{2}+8t+2$, gia tốc $a(t)=v'(t)=-2t+8$. $v'(t)=0\Leftrightarrow t=4$ và $v'$ đổi dấu từ dương sang âm nên vận tốc lớn nhất tại $t=4$.}]{Một vật chuyển động theo quy luật $s(t)=-\dfrac{1}{3}t^3+4t^2+2t$, trong đó $t$ (giây) là khoảng thời gian tính từ lúc vật bắt đầu chuyển động và $s$ (mét) là quãng đường vật đi được trong khoảng thời gian đó, $0\le t\le 8$. Thời điểm vật đạt vận tốc lớn nhất là}
  {$4\,\mathrm{s}$.}
  {$18\,\mathrm{s}$.}{$8\,\mathrm{s}$.}{$5\,\mathrm{s}$.}

\cauTN[id=C1-B5-D01-VD-031,muc=VD,loigiai={Vận tốc $v(t)=x'(t)=-4 \pi\sin(\frac{\pi}{2} t)$ (cm/s), gia tốc $a(t)=v'(t)=-2 \pi^{2}\cos(\frac{\pi}{2} t)$ (cm/s$^2$). $a(\dfrac{1}{2})=- \sqrt{2} \pi^{2}$.}]{Một vật dao động điều hòa có phương trình li độ $x(t)=8\cos(\frac{\pi}{2} t)$ (cm), với $t$ tính bằng giây. Gia tốc của vật tại thời điểm $t=\dfrac{1}{2}$ giây bằng}
  {$- \sqrt{2} \pi^{2}\,\mathrm{cm/s^2}$.}
  {$- \dfrac{\sqrt{2} \pi^{2}}{2}\,\mathrm{cm/s^2}$.}{$- \sqrt{2} \pi^{2} + 1\,\mathrm{cm/s^2}$.}{$- \sqrt{2} \pi^{2} - 1\,\mathrm{cm/s^2}$.}

\cauTN[id=C1-B5-D01-VD-032,muc=VD,loigiai={Vận tốc $v(t)=x'(t)=-5 \pi\sin(\pi t)$ (cm/s), gia tốc $a(t)=v'(t)=-5 \pi^{2}\cos(\pi t)$ (cm/s$^2$). $v(\dfrac{1}{6})=- \dfrac{5 \pi}{2}$.}]{Một vật dao động điều hòa có phương trình li độ $x(t)=5\cos(\pi t)$ (cm), với $t$ tính bằng giây. Vận tốc của vật tại thời điểm $t=\dfrac{1}{6}$ giây bằng}
  {$- \dfrac{5 \pi}{2}\,\mathrm{cm/s}$.}
  {$- \dfrac{5 \pi}{2} - 1\,\mathrm{cm/s}$.}{$- \dfrac{5 \pi}{3}\,\mathrm{cm/s}$.}{$- \dfrac{5 \pi}{4}\,\mathrm{cm/s}$.}

\cauTN[id=C1-B5-D01-VD-033,muc=VD,loigiai={Vận tốc $v(t)=s'(t)=-t^{2}+16t+5$, gia tốc $a(t)=v'(t)=-2t+16$. $v'(t)=0\Leftrightarrow t=8$ và $v'$ đổi dấu từ dương sang âm nên vận tốc lớn nhất tại $t=8$.}]{Một vật chuyển động theo quy luật $s(t)=-\dfrac{1}{3}t^3+8t^2+5t$, trong đó $t$ (giây) là khoảng thời gian tính từ lúc vật bắt đầu chuyển động và $s$ (mét) là quãng đường vật đi được trong khoảng thời gian đó, $0\le t\le 13$. Thời điểm vật đạt vận tốc lớn nhất là}
  {$8\,\mathrm{s}$.}
  {$69\,\mathrm{s}$.}{$16\,\mathrm{s}$.}{$6\,\mathrm{s}$.}

\cauTN[id=C1-B5-D01-VD-034,muc=VD,loigiai={Quãng đường $s(10)=800$.}]{Một vật chuyển động theo quy luật $s(t)=-\dfrac{3}{3}t^3+18t^2$, trong đó $t$ (giây) là khoảng thời gian tính từ lúc vật bắt đầu chuyển động và $s$ (mét) là quãng đường vật đi được trong khoảng thời gian đó, $0\le t\le 12$. Quãng đường vật đi được trong $10$ giây đầu tiên bằng}
  {$800\,\mathrm{m}$.}
  {$1600\,\mathrm{m}$.}{$850\,\mathrm{m}$.}{$700\,\mathrm{m}$.}

\cauTN[id=C1-B5-D01-VD-035,muc=VD,loigiai={Vận tốc $v(t)=s'(t)=-2t^{2}+12t+5$, gia tốc $a(t)=v'(t)=-4t+12$. $v'(t)=0\Leftrightarrow t=3$; $v$ đồng biến trên $(0;3)$, nghịch biến trên $(3;6)$ nên $v_{\max}=v(3)=23$.}]{Một vật chuyển động theo quy luật $s(t)=-\dfrac{2}{3}t^3+6t^2+5t$, trong đó $t$ (giây) là khoảng thời gian tính từ lúc vật bắt đầu chuyển động và $s$ (mét) là quãng đường vật đi được trong khoảng thời gian đó, $0\le t\le 6$. Vận tốc lớn nhất của vật trong khoảng thời gian trên bằng}
  {$23\,\mathrm{m/s}$.}
  {$3\,\mathrm{m/s}$.}{$5\,\mathrm{m/s}$.}{$19\,\mathrm{m/s}$.}

\cauTN[id=C1-B5-D01-VD-036,muc=VD,loigiai={Vận tốc $v(t)=s'(t)=-3t^{2}+36t+2$, gia tốc $a(t)=v'(t)=-6t+36$. $v'(t)=0\Leftrightarrow t=6$ và $v'$ đổi dấu từ dương sang âm nên vận tốc lớn nhất tại $t=6$.}]{Một vật chuyển động theo quy luật $s(t)=-\dfrac{3}{3}t^3+18t^2+2t$, trong đó $t$ (giây) là khoảng thời gian tính từ lúc vật bắt đầu chuyển động và $s$ (mét) là quãng đường vật đi được trong khoảng thời gian đó, $0\le t\le 12$. Thời điểm vật đạt vận tốc lớn nhất là}
  {$6\,\mathrm{s}$.}
  {$110\,\mathrm{s}$.}{$12\,\mathrm{s}$.}{$4\,\mathrm{s}$.}

\cauDS[id=C1-B5-D01-VD-037,muc=VD,loigiai={Vận tốc $v(t)=x'(t)=-32 \pi\sin(4 \pi t)$ (cm/s), gia tốc $a(t)=v'(t)=-128 \pi^{2}\cos(4 \pi t)$ (cm/s$^2$). $a(\dfrac{1}{2})=- 128 \pi^{2}$. Vận tốc $v(t)=x'(t)=-32 \pi\sin(4 \pi t)$ (cm/s), gia tốc $a(t)=v'(t)=-128 \pi^{2}\cos(4 \pi t)$ (cm/s$^2$). $v(\dfrac{1}{2})=0$. Ta có $x(\dfrac{1}{2})=8$. Vận tốc $v(t)=x'(t)=-32 \pi\sin(4 \pi t)$ (cm/s), gia tốc $a(t)=v'(t)=-128 \pi^{2}\cos(4 \pi t)$ (cm/s$^2$). Vì $|\sin(4 \pi t)|\le1$ nên $|v(t)|\le 32 \pi$, dấu bằng xảy ra khi $\sin(4 \pi t)=\pm1$. a) Sai: giá trị đúng là $8\,\mathrm{cm}$. b) Sai: giá trị đúng là $0\,\mathrm{cm/s}$. c) Đúng. d) Đúng.}]{Một vật dao động điều hòa có phương trình li độ $x(t)=8\cos(4 \pi t)$ (cm), với $t$ tính bằng giây. Xét tính đúng sai của các mệnh đề sau:}{%
  \y{Li độ của vật tại thời điểm $t=\dfrac{1}{2}$ giây là $12\,\mathrm{cm}$.}{s}
  \y{Vận tốc của vật tại thời điểm $t=\dfrac{1}{2}$ giây là $4\,\mathrm{cm/s}$.}{s}
  \y{Gia tốc của vật tại thời điểm $t=\dfrac{1}{2}$ giây là $- 128 \pi^{2}\,\mathrm{cm/s^2}$.}{d}
  \y{Tốc độ cực đại của vật là $32 \pi\,\mathrm{cm/s}$.}{d}}

\cauDS[id=C1-B5-D01-VD-038,muc=VD,loigiai={Vận tốc $v(t)=s'(t)=-2t^{2}+40t+5$, gia tốc $a(t)=v'(t)=-4t+40$. $v'(t)=0\Leftrightarrow t=10$ và $v'$ đổi dấu từ dương sang âm nên vận tốc lớn nhất tại $t=10$. Vận tốc $v(t)=s'(t)=-2t^{2}+40t+5$, gia tốc $a(t)=v'(t)=-4t+40$. Do đó $a(3)=28$. Vận tốc $v(t)=s'(t)=-2t^{2}+40t+5$, gia tốc $a(t)=v'(t)=-4t+40$. $v'(t)=0\Leftrightarrow t=10$; $v$ đồng biến trên $(0;10)$, nghịch biến trên $(10;15)$ nên $v_{\max}=v(10)=205$. Vận tốc $v(t)=s'(t)=-2t^{2}+40t+5$, gia tốc $a(t)=v'(t)=-4t+40$. Do đó $v(3)=107$. a) Đúng. b) Đúng. c) Sai: giá trị đúng là $107\,\mathrm{m/s}$. d) Sai: giá trị đúng là $205\,\mathrm{m/s}$.}]{Một vật chuyển động theo quy luật $s(t)=-\dfrac{2}{3}t^3+20t^2+5t$, trong đó $t$ (giây) là khoảng thời gian tính từ lúc vật bắt đầu chuyển động và $s$ (mét) là quãng đường vật đi được trong khoảng thời gian đó, $0\le t\le 15$. Xét tính đúng sai của các mệnh đề sau:}{%
  \y{Gia tốc của vật tại thời điểm $t=3$ giây là $28\,\mathrm{m/s^2}$.}{d}
  \y{Vật đạt vận tốc lớn nhất tại thời điểm $t=10\,\mathrm{s}$.}{d}
  \y{Vận tốc của vật tại thời điểm $t=3$ giây là $214\,\mathrm{m/s}$.}{s}
  \y{Vận tốc lớn nhất của vật là $410\,\mathrm{m/s}$.}{s}}

\cauDS[id=C1-B5-D01-VD-039,muc=VD,loigiai={Vận tốc $v(t)=s'(t)=-t^{2}+16t+3$, gia tốc $a(t)=v'(t)=-2t+16$. Do đó $v(4)=51$. Vận tốc $v(t)=s'(t)=-t^{2}+16t+3$, gia tốc $a(t)=v'(t)=-2t+16$. $v'(t)=0\Leftrightarrow t=8$; $v$ đồng biến trên $(0;8)$, nghịch biến trên $(8;13)$ nên $v_{\max}=v(8)=67$. Vận tốc $v(t)=s'(t)=-t^{2}+16t+3$, gia tốc $a(t)=v'(t)=-2t+16$. $v'(t)=0\Leftrightarrow t=8$ và $v'$ đổi dấu từ dương sang âm nên vận tốc lớn nhất tại $t=8$. Quãng đường $s(4)=\dfrac{356}{3}$. a) Đúng. b) Đúng. c) Đúng. d) Sai: giá trị đúng là $67\,\mathrm{m/s}$.}]{Một vật chuyển động theo quy luật $s(t)=-\dfrac{1}{3}t^3+8t^2+3t$, trong đó $t$ (giây) là khoảng thời gian tính từ lúc vật bắt đầu chuyển động và $s$ (mét) là quãng đường vật đi được trong khoảng thời gian đó, $0\le t\le 13$. Xét tính đúng sai của các mệnh đề sau:}{%
  \y{Quãng đường vật đi được trong $4$ giây đầu là $\dfrac{356}{3}\,\mathrm{m}$.}{d}
  \y{Vật đạt vận tốc lớn nhất tại thời điểm $t=8\,\mathrm{s}$.}{d}
  \y{Vận tốc của vật tại thời điểm $t=4$ giây là $51\,\mathrm{m/s}$.}{d}
  \y{Vận tốc lớn nhất của vật là $63\,\mathrm{m/s}$.}{s}}

\cauDS[id=C1-B5-D01-VD-040,muc=VD,loigiai={Vận tốc $v(t)=s'(t)=-3t^{2}+30t$, gia tốc $a(t)=v'(t)=-6t+30$. $v'(t)=0\Leftrightarrow t=5$; $v$ đồng biến trên $(0;5)$, nghịch biến trên $(5;10)$ nên $v_{\max}=v(5)=75$. Quãng đường $s(2)=52$. Vận tốc $v(t)=s'(t)=-3t^{2}+30t$, gia tốc $a(t)=v'(t)=-6t+30$. $v'(t)=0\Leftrightarrow t=5$ và $v'$ đổi dấu từ dương sang âm nên vận tốc lớn nhất tại $t=5$. Vận tốc $v(t)=s'(t)=-3t^{2}+30t$, gia tốc $a(t)=v'(t)=-6t+30$. Do đó $a(2)=18$. a) Sai: giá trị đúng là $75\,\mathrm{m/s}$. b) Đúng. c) Đúng. d) Đúng.}]{Một vật chuyển động theo quy luật $s(t)=-\dfrac{3}{3}t^3+15t^2$, trong đó $t$ (giây) là khoảng thời gian tính từ lúc vật bắt đầu chuyển động và $s$ (mét) là quãng đường vật đi được trong khoảng thời gian đó, $0\le t\le 10$. Xét tính đúng sai của các mệnh đề sau:}{%
  \y{Vận tốc lớn nhất của vật là $150\,\mathrm{m/s}$.}{s}
  \y{Vật đạt vận tốc lớn nhất tại thời điểm $t=5\,\mathrm{s}$.}{d}
  \y{Quãng đường vật đi được trong $2$ giây đầu là $52\,\mathrm{m}$.}{d}
  \y{Gia tốc của vật tại thời điểm $t=2$ giây là $18\,\mathrm{m/s^2}$.}{d}}

\cauDS[id=C1-B5-D01-VD-041,muc=VD,loigiai={Quãng đường $s(2)=\dfrac{224}{3}$. Vận tốc $v(t)=s'(t)=-2t^{2}+40t$, gia tốc $a(t)=v'(t)=-4t+40$. $v'(t)=0\Leftrightarrow t=10$; $v$ đồng biến trên $(0;10)$, nghịch biến trên $(10;20)$ nên $v_{\max}=v(10)=200$. Vận tốc $v(t)=s'(t)=-2t^{2}+40t$, gia tốc $a(t)=v'(t)=-4t+40$. $v'(t)=0\Leftrightarrow t=10$ và $v'$ đổi dấu từ dương sang âm nên vận tốc lớn nhất tại $t=10$. Vận tốc $v(t)=s'(t)=-2t^{2}+40t$, gia tốc $a(t)=v'(t)=-4t+40$. Do đó $a(2)=32$. a) Đúng. b) Sai: giá trị đúng là $200\,\mathrm{m/s}$. c) Đúng. d) Đúng.}]{Một vật chuyển động theo quy luật $s(t)=-\dfrac{2}{3}t^3+20t^2$, trong đó $t$ (giây) là khoảng thời gian tính từ lúc vật bắt đầu chuyển động và $s$ (mét) là quãng đường vật đi được trong khoảng thời gian đó, $0\le t\le 20$. Xét tính đúng sai của các mệnh đề sau:}{%
  \y{Vật đạt vận tốc lớn nhất tại thời điểm $t=10\,\mathrm{s}$.}{d}
  \y{Vận tốc lớn nhất của vật là $240\,\mathrm{m/s}$.}{s}
  \y{Quãng đường vật đi được trong $2$ giây đầu là $\dfrac{224}{3}\,\mathrm{m}$.}{d}
  \y{Gia tốc của vật tại thời điểm $t=2$ giây là $32\,\mathrm{m/s^2}$.}{d}}

\cauDS[id=C1-B5-D01-VD-042,muc=VD,loigiai={Ta có $x(\dfrac{1}{2})=0$. Vận tốc $v(t)=x'(t)=-5 \pi\sin(\pi t)$ (cm/s), gia tốc $a(t)=v'(t)=-5 \pi^{2}\cos(\pi t)$ (cm/s$^2$). $a(\dfrac{1}{2})=0$. Vận tốc $v(t)=x'(t)=-5 \pi\sin(\pi t)$ (cm/s), gia tốc $a(t)=v'(t)=-5 \pi^{2}\cos(\pi t)$ (cm/s$^2$). Vì $|\sin(\pi t)|\le1$ nên $|v(t)|\le 5 \pi$, dấu bằng xảy ra khi $\sin(\pi t)=\pm1$. Vận tốc $v(t)=x'(t)=-5 \pi\sin(\pi t)$ (cm/s), gia tốc $a(t)=v'(t)=-5 \pi^{2}\cos(\pi t)$ (cm/s$^2$). $v(\dfrac{1}{2})=- 5 \pi$. a) Sai: giá trị đúng là $0\,\mathrm{cm}$. b) Đúng. c) Đúng. d) Đúng.}]{Một vật dao động điều hòa có phương trình li độ $x(t)=5\cos(\pi t)$ (cm), với $t$ tính bằng giây. Xét tính đúng sai của các mệnh đề sau:}{%
  \y{Li độ của vật tại thời điểm $t=\dfrac{1}{2}$ giây là $-2\,\mathrm{cm}$.}{s}
  \y{Tốc độ cực đại của vật là $5 \pi\,\mathrm{cm/s}$.}{d}
  \y{Vận tốc của vật tại thời điểm $t=\dfrac{1}{2}$ giây là $- 5 \pi\,\mathrm{cm/s}$.}{d}
  \y{Gia tốc của vật tại thời điểm $t=\dfrac{1}{2}$ giây là $0\,\mathrm{cm/s^2}$.}{d}}

\cauDS[id=C1-B5-D01-VD-043,muc=VD,loigiai={Vận tốc $v(t)=x'(t)=-16 \pi\sin(4 \pi t)$ (cm/s), gia tốc $a(t)=v'(t)=-64 \pi^{2}\cos(4 \pi t)$ (cm/s$^2$). Vì $|\sin(4 \pi t)|\le1$ nên $|v(t)|\le 16 \pi$, dấu bằng xảy ra khi $\sin(4 \pi t)=\pm1$. Vận tốc $v(t)=x'(t)=-16 \pi\sin(4 \pi t)$ (cm/s), gia tốc $a(t)=v'(t)=-64 \pi^{2}\cos(4 \pi t)$ (cm/s$^2$). $a(\dfrac{1}{4})=64 \pi^{2}$. Vận tốc $v(t)=x'(t)=-16 \pi\sin(4 \pi t)$ (cm/s), gia tốc $a(t)=v'(t)=-64 \pi^{2}\cos(4 \pi t)$ (cm/s$^2$). $v(\dfrac{1}{4})=0$. Ta có $x(\dfrac{1}{4})=-4$. a) Sai: giá trị đúng là $-4\,\mathrm{cm}$. b) Đúng. c) Đúng. d) Đúng.}]{Một vật dao động điều hòa có phương trình li độ $x(t)=4\cos(4 \pi t)$ (cm), với $t$ tính bằng giây. Xét tính đúng sai của các mệnh đề sau:}{%
  \y{Li độ của vật tại thời điểm $t=\dfrac{1}{4}$ giây là $-2\,\mathrm{cm}$.}{s}
  \y{Vận tốc của vật tại thời điểm $t=\dfrac{1}{4}$ giây là $0\,\mathrm{cm/s}$.}{d}
  \y{Tốc độ cực đại của vật là $16 \pi\,\mathrm{cm/s}$.}{d}
  \y{Gia tốc của vật tại thời điểm $t=\dfrac{1}{4}$ giây là $64 \pi^{2}\,\mathrm{cm/s^2}$.}{d}}

\cauDS[id=C1-B5-D01-VD-044,muc=VD,loigiai={Vận tốc $v(t)=s'(t)=-2t^{2}+20t+5$, gia tốc $a(t)=v'(t)=-4t+20$. Do đó $a(5)=0$. Vận tốc $v(t)=s'(t)=-2t^{2}+20t+5$, gia tốc $a(t)=v'(t)=-4t+20$. $v'(t)=0\Leftrightarrow t=5$; $v$ đồng biến trên $(0;5)$, nghịch biến trên $(5;10)$ nên $v_{\max}=v(5)=55$. Quãng đường $s(5)=\dfrac{575}{3}$. Vận tốc $v(t)=s'(t)=-2t^{2}+20t+5$, gia tốc $a(t)=v'(t)=-4t+20$. Do đó $v(5)=55$. a) Đúng. b) Sai: giá trị đúng là $55\,\mathrm{m/s}$. c) Đúng. d) Sai: giá trị đúng là $\dfrac{575}{3}\,\mathrm{m}$.}]{Một vật chuyển động theo quy luật $s(t)=-\dfrac{2}{3}t^3+10t^2+5t$, trong đó $t$ (giây) là khoảng thời gian tính từ lúc vật bắt đầu chuyển động và $s$ (mét) là quãng đường vật đi được trong khoảng thời gian đó, $0\le t\le 10$. Xét tính đúng sai của các mệnh đề sau:}{%
  \y{Vận tốc của vật tại thời điểm $t=5$ giây là $55\,\mathrm{m/s}$.}{d}
  \y{Vận tốc lớn nhất của vật là $110\,\mathrm{m/s}$.}{s}
  \y{Gia tốc của vật tại thời điểm $t=5$ giây là $0\,\mathrm{m/s^2}$.}{d}
  \y{Quãng đường vật đi được trong $5$ giây đầu là $\dfrac{572}{3}\,\mathrm{m}$.}{s}}

\cauDS[id=C1-B5-D01-VD-045,muc=VD,loigiai={Vận tốc $v(t)=s'(t)=-2t^{2}+48t+5$, gia tốc $a(t)=v'(t)=-4t+48$. $v'(t)=0\Leftrightarrow t=12$ và $v'$ đổi dấu từ dương sang âm nên vận tốc lớn nhất tại $t=12$. Vận tốc $v(t)=s'(t)=-2t^{2}+48t+5$, gia tốc $a(t)=v'(t)=-4t+48$. $v'(t)=0\Leftrightarrow t=12$; $v$ đồng biến trên $(0;12)$, nghịch biến trên $(12;15)$ nên $v_{\max}=v(12)=293$. Vận tốc $v(t)=s'(t)=-2t^{2}+48t+5$, gia tốc $a(t)=v'(t)=-4t+48$. Do đó $a(1)=44$. Quãng đường $s(1)=\dfrac{85}{3}$. a) Sai: giá trị đúng là $12\,\mathrm{s}$. b) Sai: giá trị đúng là $\dfrac{85}{3}\,\mathrm{m}$. c) Đúng. d) Đúng.}]{Một vật chuyển động theo quy luật $s(t)=-\dfrac{2}{3}t^3+24t^2+5t$, trong đó $t$ (giây) là khoảng thời gian tính từ lúc vật bắt đầu chuyển động và $s$ (mét) là quãng đường vật đi được trong khoảng thời gian đó, $0\le t\le 15$. Xét tính đúng sai của các mệnh đề sau:}{%
  \y{Vật đạt vận tốc lớn nhất tại thời điểm $t=6\,\mathrm{s}$.}{s}
  \y{Quãng đường vật đi được trong $1$ giây đầu là $85\,\mathrm{m}$.}{s}
  \y{Gia tốc của vật tại thời điểm $t=1$ giây là $44\,\mathrm{m/s^2}$.}{d}
  \y{Vận tốc lớn nhất của vật là $293\,\mathrm{m/s}$.}{d}}

\cauDS[id=C1-B5-D01-VD-046,muc=VD,loigiai={Vận tốc $v(t)=s'(t)=-3t^{2}+54t+3$, gia tốc $a(t)=v'(t)=-6t+54$. Do đó $a(3)=36$. Vận tốc $v(t)=s'(t)=-3t^{2}+54t+3$, gia tốc $a(t)=v'(t)=-6t+54$. Do đó $v(3)=138$. Quãng đường $s(3)=225$. Vận tốc $v(t)=s'(t)=-3t^{2}+54t+3$, gia tốc $a(t)=v'(t)=-6t+54$. $v'(t)=0\Leftrightarrow t=9$ và $v'$ đổi dấu từ dương sang âm nên vận tốc lớn nhất tại $t=9$. a) Sai: giá trị đúng là $138\,\mathrm{m/s}$. b) Đúng. c) Sai: giá trị đúng là $225\,\mathrm{m}$. d) Đúng.}]{Một vật chuyển động theo quy luật $s(t)=-\dfrac{3}{3}t^3+27t^2+3t$, trong đó $t$ (giây) là khoảng thời gian tính từ lúc vật bắt đầu chuyển động và $s$ (mét) là quãng đường vật đi được trong khoảng thời gian đó, $0\le t\le 14$. Xét tính đúng sai của các mệnh đề sau:}{%
  \y{Vận tốc của vật tại thời điểm $t=3$ giây là $158\,\mathrm{m/s}$.}{s}
  \y{Vật đạt vận tốc lớn nhất tại thời điểm $t=9\,\mathrm{s}$.}{d}
  \y{Quãng đường vật đi được trong $3$ giây đầu là $275\,\mathrm{m}$.}{s}
  \y{Gia tốc của vật tại thời điểm $t=3$ giây là $36\,\mathrm{m/s^2}$.}{d}}

\cauDS[id=C1-B5-D01-VD-047,muc=VD,loigiai={Vận tốc $v(t)=s'(t)=-3t^{2}+24t+3$, gia tốc $a(t)=v'(t)=-6t+24$. $v'(t)=0\Leftrightarrow t=4$; $v$ đồng biến trên $(0;4)$, nghịch biến trên $(4;7)$ nên $v_{\max}=v(4)=51$. Vận tốc $v(t)=s'(t)=-3t^{2}+24t+3$, gia tốc $a(t)=v'(t)=-6t+24$. Do đó $a(5)=-6$. Quãng đường $s(5)=190$. Vận tốc $v(t)=s'(t)=-3t^{2}+24t+3$, gia tốc $a(t)=v'(t)=-6t+24$. Do đó $v(5)=48$. a) Sai: giá trị đúng là $-6\,\mathrm{m/s^2}$. b) Đúng. c) Đúng. d) Sai: giá trị đúng là $48\,\mathrm{m/s}$.}]{Một vật chuyển động theo quy luật $s(t)=-\dfrac{3}{3}t^3+12t^2+3t$, trong đó $t$ (giây) là khoảng thời gian tính từ lúc vật bắt đầu chuyển động và $s$ (mét) là quãng đường vật đi được trong khoảng thời gian đó, $0\le t\le 7$. Xét tính đúng sai của các mệnh đề sau:}{%
  \y{Gia tốc của vật tại thời điểm $t=5$ giây là $-3\,\mathrm{m/s^2}$.}{s}
  \y{Vận tốc lớn nhất của vật là $51\,\mathrm{m/s}$.}{d}
  \y{Quãng đường vật đi được trong $5$ giây đầu là $190\,\mathrm{m}$.}{d}
  \y{Vận tốc của vật tại thời điểm $t=5$ giây là $38\,\mathrm{m/s}$.}{s}}

\cauDS[id=C1-B5-D01-VD-048,muc=VD,loigiai={Vận tốc $v(t)=s'(t)=-t^{2}+18t+2$, gia tốc $a(t)=v'(t)=-2t+18$. Do đó $v(17)=19$. Vận tốc $v(t)=s'(t)=-t^{2}+18t+2$, gia tốc $a(t)=v'(t)=-2t+18$. $v'(t)=0\Leftrightarrow t=9$ và $v'$ đổi dấu từ dương sang âm nên vận tốc lớn nhất tại $t=9$. Vận tốc $v(t)=s'(t)=-t^{2}+18t+2$, gia tốc $a(t)=v'(t)=-2t+18$. Do đó $a(17)=-16$. Vận tốc $v(t)=s'(t)=-t^{2}+18t+2$, gia tốc $a(t)=v'(t)=-2t+18$. $v'(t)=0\Leftrightarrow t=9$; $v$ đồng biến trên $(0;9)$, nghịch biến trên $(9;18)$ nên $v_{\max}=v(9)=83$. a) Đúng. b) Sai: giá trị đúng là $19\,\mathrm{m/s}$. c) Sai: giá trị đúng là $-16\,\mathrm{m/s^2}$. d) Đúng.}]{Một vật chuyển động theo quy luật $s(t)=-\dfrac{1}{3}t^3+9t^2+2t$, trong đó $t$ (giây) là khoảng thời gian tính từ lúc vật bắt đầu chuyển động và $s$ (mét) là quãng đường vật đi được trong khoảng thời gian đó, $0\le t\le 18$. Xét tính đúng sai của các mệnh đề sau:}{%
  \y{Vận tốc lớn nhất của vật là $83\,\mathrm{m/s}$.}{d}
  \y{Vận tốc của vật tại thời điểm $t=17$ giây là $23\,\mathrm{m/s}$.}{s}
  \y{Gia tốc của vật tại thời điểm $t=17$ giây là $-32\,\mathrm{m/s^2}$.}{s}
  \y{Vật đạt vận tốc lớn nhất tại thời điểm $t=9\,\mathrm{s}$.}{d}}

\cauTLN[id=C1-B5-D01-VD-049,muc=VD,loigiai={Vận tốc $v(t)=s'(t)=-t^{2}+6t+5$, gia tốc $a(t)=v'(t)=-2t+6$. $v'(t)=0\Leftrightarrow t=3$; $v$ đồng biến trên $(0;3)$, nghịch biến trên $(3;6)$ nên $v_{\max}=v(3)=14$. Kết quả: $14$.}]{Một vật chuyển động theo quy luật $s(t)=-\dfrac{1}{3}t^3+3t^2+5t$, trong đó $t$ (giây) là khoảng thời gian tính từ lúc vật bắt đầu chuyển động và $s$ (mét) là quãng đường vật đi được trong khoảng thời gian đó, $0\le t\le 6$. Tìm vận tốc lớn nhất của vật trong khoảng thời gian trên (đơn vị: $\mathrm{m/s}$).}{14}

\cauTLN[id=C1-B5-D01-VD-050,muc=VD,loigiai={Ta có $x(\dfrac{1}{2})=6$. Kết quả: $6$.}]{Một vật dao động điều hòa có phương trình li độ $x(t)=6\cos(4 \pi t)$ (cm), với $t$ tính bằng giây. Tìm li độ của vật tại thời điểm $t=\dfrac{1}{2}$ giây (đơn vị: $\mathrm{cm}$).}{6}

\cauTLN[id=C1-B5-D01-VD-051,muc=VD,loigiai={Vận tốc $v(t)=s'(t)=-t^{2}+8t+5$, gia tốc $a(t)=v'(t)=-2t+8$. Do đó $a(2)=4$. Kết quả: $4$.}]{Một vật chuyển động theo quy luật $s(t)=-\dfrac{1}{3}t^3+4t^2+5t$, trong đó $t$ (giây) là khoảng thời gian tính từ lúc vật bắt đầu chuyển động và $s$ (mét) là quãng đường vật đi được trong khoảng thời gian đó, $0\le t\le 10$. Tìm gia tốc của vật tại thời điểm $t=2$ giây (đơn vị: $\mathrm{m/s^2}$).}{4}

\cauTLN[id=C1-B5-D01-VD-052,muc=VD,loigiai={Vận tốc $v(t)=s'(t)=-t^{2}+16t+2$, gia tốc $a(t)=v'(t)=-2t+16$. Do đó $v(4)=50$. Kết quả: $50$.}]{Một vật chuyển động theo quy luật $s(t)=-\dfrac{1}{3}t^3+8t^2+2t$, trong đó $t$ (giây) là khoảng thời gian tính từ lúc vật bắt đầu chuyển động và $s$ (mét) là quãng đường vật đi được trong khoảng thời gian đó, $0\le t\le 16$. Tìm vận tốc của vật tại thời điểm $t=4$ giây (đơn vị: $\mathrm{m/s}$).}{50}

\cauTLN[id=C1-B5-D01-VD-053,muc=VD,loigiai={Vận tốc $v(t)=x'(t)=-4 \pi\sin(\pi t)$ (cm/s), gia tốc $a(t)=v'(t)=-4 \pi^{2}\cos(\pi t)$ (cm/s$^2$). Vì $|\sin(\pi t)|\le1$ nên $|v(t)|\le 4 \pi$, dấu bằng xảy ra khi $\sin(\pi t)=\pm1$. Kết quả: $12{,}6$.}]{Một vật dao động điều hòa có phương trình li độ $x(t)=4\cos(\pi t)$ (cm), với $t$ tính bằng giây. Tìm tốc độ cực đại của vật (giá trị lớn nhất của $|v(t)|$) (đơn vị: $\mathrm{cm/s}$) (làm tròn kết quả đến hàng phần mười).}{12,6}

\cauTLN[id=C1-B5-D01-VD-054,muc=VD,loigiai={Vận tốc $v(t)=x'(t)=-16 \pi\sin(2 \pi t)$ (cm/s), gia tốc $a(t)=v'(t)=-32 \pi^{2}\cos(2 \pi t)$ (cm/s$^2$). $v(\dfrac{1}{3})=- 8 \sqrt{3} \pi$. Kết quả: $-44$.}]{Một vật dao động điều hòa có phương trình li độ $x(t)=8\cos(2 \pi t)$ (cm), với $t$ tính bằng giây. Tìm vận tốc của vật tại thời điểm $t=\dfrac{1}{3}$ giây (đơn vị: $\mathrm{cm/s}$) (làm tròn kết quả đến hàng đơn vị).}{-44}

\cauTLN[id=C1-B5-D01-VD-055,muc=VD,loigiai={Vận tốc $v(t)=x'(t)=-2 \pi\sin(\pi t)$ (cm/s), gia tốc $a(t)=v'(t)=-2 \pi^{2}\cos(\pi t)$ (cm/s$^2$). $a(\dfrac{1}{6})=- \sqrt{3} \pi^{2}$. Kết quả: $-17$.}]{Một vật dao động điều hòa có phương trình li độ $x(t)=2\cos(\pi t)$ (cm), với $t$ tính bằng giây. Tìm gia tốc của vật tại thời điểm $t=\dfrac{1}{6}$ giây (đơn vị: $\mathrm{cm/s^2}$) (làm tròn kết quả đến hàng đơn vị).}{-17}

\cauTLN[id=C1-B5-D01-VD-056,muc=VD,loigiai={Vận tốc $v(t)=s'(t)=-t^{2}+16t$, gia tốc $a(t)=v'(t)=-2t+16$. Do đó $a(2)=12$. Kết quả: $12$.}]{Một vật chuyển động theo quy luật $s(t)=-\dfrac{1}{3}t^3+8t^2$, trong đó $t$ (giây) là khoảng thời gian tính từ lúc vật bắt đầu chuyển động và $s$ (mét) là quãng đường vật đi được trong khoảng thời gian đó, $0\le t\le 16$. Tìm gia tốc của vật tại thời điểm $t=2$ giây (đơn vị: $\mathrm{m/s^2}$).}{12}

\cauTLN[id=C1-B5-D01-VD-057,muc=VD,loigiai={Vận tốc $v(t)=x'(t)=-3 \pi\sin(\frac{\pi}{2} t)$ (cm/s), gia tốc $a(t)=v'(t)=-\dfrac{3 \pi^{2}}{2}\cos(\frac{\pi}{2} t)$ (cm/s$^2$). $a(\dfrac{1}{2})=- \dfrac{3 \sqrt{2} \pi^{2}}{4}$. Kết quả: $-10$.}]{Một vật dao động điều hòa có phương trình li độ $x(t)=6\cos(\frac{\pi}{2} t)$ (cm), với $t$ tính bằng giây. Tìm gia tốc của vật tại thời điểm $t=\dfrac{1}{2}$ giây (đơn vị: $\mathrm{cm/s^2}$) (làm tròn kết quả đến hàng đơn vị).}{-10}

\cauTLN[id=C1-B5-D01-VD-058,muc=VD,loigiai={Vận tốc $v(t)=s'(t)=-3t^{2}+30t$, gia tốc $a(t)=v'(t)=-6t+30$. Do đó $a(3)=12$. Kết quả: $12$.}]{Một vật chuyển động theo quy luật $s(t)=-\dfrac{3}{3}t^3+15t^2$, trong đó $t$ (giây) là khoảng thời gian tính từ lúc vật bắt đầu chuyển động và $s$ (mét) là quãng đường vật đi được trong khoảng thời gian đó, $0\le t\le 9$. Tìm gia tốc của vật tại thời điểm $t=3$ giây (đơn vị: $\mathrm{m/s^2}$).}{12}

\cauTLN[id=C1-B5-D01-VD-059,muc=VD,loigiai={Vận tốc $v(t)=x'(t)=-2 \pi\sin(\frac{\pi}{2} t)$ (cm/s), gia tốc $a(t)=v'(t)=-\pi^{2}\cos(\frac{\pi}{2} t)$ (cm/s$^2$). $a(\dfrac{1}{2})=- \dfrac{\sqrt{2} \pi^{2}}{2}$. Kết quả: $-7$.}]{Một vật dao động điều hòa có phương trình li độ $x(t)=4\cos(\frac{\pi}{2} t)$ (cm), với $t$ tính bằng giây. Tìm gia tốc của vật tại thời điểm $t=\dfrac{1}{2}$ giây (đơn vị: $\mathrm{cm/s^2}$) (làm tròn kết quả đến hàng phần mười).}{-7}

\cauTLN[id=C1-B5-D01-VD-060,muc=VD,loigiai={Vận tốc $v(t)=s'(t)=-t^{2}+12t$, gia tốc $a(t)=v'(t)=-2t+12$. Do đó $v(5)=35$. Kết quả: $35$.}]{Một vật chuyển động theo quy luật $s(t)=-\dfrac{1}{3}t^3+6t^2$, trong đó $t$ (giây) là khoảng thời gian tính từ lúc vật bắt đầu chuyển động và $s$ (mét) là quãng đường vật đi được trong khoảng thời gian đó, $0\le t\le 12$. Tìm vận tốc của vật tại thời điểm $t=5$ giây (đơn vị: $\mathrm{m/s}$).}{35}
